\input awebmac
%::::::::::
%atangle.aweb
%::::::::::
% This is ATANGLE.AWEB in text format, as of May 04, 1988.
% This program by D. E. Knuth is not copyrighted and can be used freely.
% Adaption to ADA by U. Schweigert based on TANGLE Version 2.8.
% Version 1.0 was released May 31, 1988.


% Here is TeX material that gets inserted after \input webmaca
\def\hang{\hangindent 3em\indent\ignorespaces}
\font\ninerm=cmr9
\let\mc=\ninerm % medium caps for names like ADA
\def\ADA{{\mc ADA}}
\def\ab{$\.|\ldots\.|$} % ADA brackets (|...|)
\def\v{\.{\char'174}} % vertical (|) in typewriter font
\mathchardef\BA="3224 % double arrow
\def\({} % kludge for alphabetizing certain module names

\def\title{ATANGLE}
\def\contentspagenumber{137} % should be odd
\def\topofcontents{\null\vfill
  \titlefalse % include headline on the contents page
  \def\rheader{\mainfont Appendix E\hfil \contentspagenumber}
  \centerline{\titlefont The {\ttitlefont ATANGLE} processor}
  \vskip 15pt
  \centerline{(based on {\tt TANGLE} Version 2.8)}
  \centerline{{\ninerm[ADA-Version 1.0]}}
  \vfill}
\pageno=\contentspagenumber \advance\pageno by 1


\N1\*.  Introduction.
This program converts an \.{AWEB} file to an \ADA\ file. It was written%
---originally for conversion to {\mc PASCAL}---by D. E. Knuth in September,
1981; a somewhat similar {\mc SAIL} program had been developed in March,
1979. Since this program describes itself, a bootstrapping process involving
hand-translation had to be used to get started. The conversion to \ADA\
has been done by U. Schweigert in May 1988.

For large \.{AWEB} files one should have a large memory, since \.{ATANGLE}
keeps all the \ADA\ text in memory (in an abbreviated form). This program
uses only \ADA-features as described in ``The Programming Language \ADA\
Reference Manual'' and is thus conform with ANSI/MIL-STD-1815A-1983.

The system-dependent portions of \.{ATANGLE} can be identified by looking at
the entries for `system dependencies' in the index below.

The ``banner line'' defined here should be changed whenever \.{ATANGLE} is
modified.

\Y\P\4\D\\{banner}\?\S\?\.{"This\ is\ ATANGLE,\ Version\ 1.1\ (NRaD\ Version)"}%
\?\6
\?\par
\fi

\M2. The program consists of several packages that are declared right now;
each of these packages and either the specification and the body of
the packages are send to a seperate file. The main program itself
is declared later.

\fi

\M3\*. The first package we declare is called \\{int\_number\_io} and enables
us
to output integer numbers. Further we declare a package \\{data\_structures},
which will contain all the global constants, types, variables, and some
procedures used throughout the program; it will (nearly) contain code of
the sections ``The character set,'' ``Data structures,'' and ``Tokens.''

\Y\P\O{data\_structures.ads}\?\6
\?\&{with}\8\p{TEXT\_IO};\6
\&{with}\8\p{TEXT\_IO};\6
\&{with}\8\\{int\_number\_io};\6
\&{package}\8\1\\{data\_structures}\2\8\&{is}\8\1\6
\X10:Types and objects visible from \\{data\_structures}\X\6
\X47:Specification of procedures visible from \\{data\_structures}\X\2\6
\&{end}\8\\{data\_structures};\7
\O{data\_structures.adb}\?\6
\?\&{with}\8\\{input\_output};\6
\&{use}\8\\{input\_output};\6
\&{package}\8\1\&{body}\8\\{data\_structures}\2\8\&{is}\8\1\6
\X48:Procedures in \\{data\_structures}\X\6
\4\&{begin}\6
\X15:Set initial values\X\2\6
\&{end}\8\\{data\_structures};\par
\fi

\M4\*. The following package, called \\{input\_output}, is responsible for the
input/output features neaded by \.{ATANGLE}, and for the error handling.
It will contain code of sections ``Input and output'' and ``Reporting
errors to the user.''

\Y\P\O{input\_output.ads}\?\6
\?\&{with}\8\p{TEXT\_IO};\6
\&{with}\8\\{data\_structures};\6
\&{use}\8\\{data\_structures};\6
\&{package}\8\1\\{input\_output}\2\8\&{is}\8\1\6
\X23:Specification of procedures visible from \\{input\_output}\X\2\6
\&{end}\8\\{input\_output};\7
\O{input\_output.adb}\?\6
\?\&{with}\8\p{TEXT\_IO}\?,\39\?\\{int\_number\_io};\6
\&{package}\8\1\&{body}\8\\{input\_output}\2\8\&{is}\8\1\6
\X24:Procedures in \\{input\_output}\X\2\6
\&{end}\8\\{input\_output};\par
\fi

\M5\*. Additionally we have a package \\{hashing}, which contains the code of
sections ``Searching for identifiers''
and ``Searching for module names.''

\Y\P\O{hashing.ads}\?\6
\?\&{with}\8\\{data\_structures};\6
\&{use}\8\\{data\_structures};\6
\&{package}\8\1\\{hashing}\2\8\&{is}\8\1\6
\X50:Specification of procedures visible from \\{hashing}\X\2\6
\&{end}\8\\{hashing};\7
\O{hashing.adb}\?\6
\?\&{with}\8\p{TEXT\_IO};\6
\&{with}\8\\{input\_output};\6
\&{use}\8\\{input\_output};\6
\&{package}\8\1\&{body}\8\\{hashing}\2\8\&{is}\8\1\6
\X51:Procedures in \\{hashing}\X\2\6
\&{end}\8\\{hashing};\par
\fi

\M6\*. The first of the two major packages will manage the output process of
\.{ATANGLE} and contains code of sections ``Stacks for output,''
``Producing the output,'' and ``The big output switch.''

\Y\P\O{output\_phase.ads}\?\6
\?\&{with}\8\\{data\_structures};\6
\&{use}\8\\{data\_structures};\6
\&{package}\8\1\\{output\_phase}\2\8\&{is}\8\1\6
\&{procedure}\8\1\\{phase\_ii}\2;\2\6
\&{end}\8\\{output\_phase};\7
\O{output\_phase.adb}\?\6
\?\&{with}\8\p{TEXT\_IO}\?,\39\?\\{int\_number\_io};\6
\&{with}\8\\{input\_output};\6
\&{use}\8\\{input\_output};\6
\&{package}\8\1\&{body}\8\\{output\_phase}\2\8\&{is}\8\1\6
\X81:Procedures in \\{output\_phase}\X\2\6
\&{end}\8\\{output\_phase};\par
\fi

\M7\*. The last package will handle \.{ATANGLE}'s input process and contains
code
of sections ``Introduction to the input phase,'' ``Inputting the next
token,'' ``Scanning a macro definition,'' and ``Scanning a module.''

\Y\P\O{input\_phase.ads}\?\6
\?\&{with}\8\\{data\_structures};\6
\&{use}\8\\{data\_structures};\6
\&{with}\8\\{hashing};\6
\&{use}\8\\{hashing};\6
\&{package}\8\1\\{input\_phase}\2\8\&{is}\8\1\6
\&{procedure}\8\1\\{phase\_i}\2;\2\6
\&{end}\8\\{input\_phase};\7
\O{input\_phase.adb}\?\6
\?\&{with}\8\p{TEXT\_IO}\?,\39\?\\{int\_number\_io};\6
\&{with}\8\\{input\_output};\6
\&{use}\8\\{input\_output};\6
\&{package}\8\1\&{body}\8\\{input\_phase}\2\8\&{is}\8\1\6
\X119:Procedures in \\{input\_phase}\X\2\6
\&{end}\8\\{input\_phase};\par
\fi

\M8. Some of this code, delimited by  \&{stat} and \&{tats}, is optionally
included if statistics about \.{ATANGLE}'s memory usage are desired.
Other material is optionally included when a somewhat optimal code is
wanted; such is delimited by  \&{optim} and \&{mitpo}.

\Y\P\4\D\\{stat}\?\S\B\C{change this to `$\\{stat}\equiv\null$'   when
gathering usage statistics}\6
\?\par
\P\4\D\\{tats}\?\S\?\hbox{}\?\T\C{change this to `$\\{tats}\equiv\null$'   when
gathering usage statistics}\6
\?\par
\P\4\F\\{stat}\?\S\?\\{declare}\par
\P\4\F\\{tats}\?\S\?\\{end}\par
\P\4\D\\{optim}\?\S\B\C{change this to `$\\{optim}\equiv\null$' when optimal
 code is desired}\6
\?\par
\P\4\D\\{mitpo}\?\S\?\hbox{}\?\T\C{change this to `$\\{mitpo}\equiv\null$' when
optimal     code is desired}\6
\?\par
\P\4\F\\{optim}\?\S\?\\{declare}\par
\P\4\F\\{mitpo}\?\S\?\\{end}\par
\fi

\M9. Here are some macros for common programming idioms.

\Y\P\4\D\\{incr}\?\1(\?\#\?\2)\S\?\#\?\K\?\#\?+\?1\C{increase a variable by
unity}\par
\P\4\D\\{decr}\?\1(\?\#\?\2)\S\?\#\?\K\?\#\?-\?1\C{decrease a variable by
unity}\par
\P\4\D\\{repeat}\?\S\?\8\&{loop}\par
\P\4\D\\{until}\?\1(\?\#\?\2)\S\&{exit}\8\&{when}\?\8\#;\2\6
\&{end}\8\&{loop}\par
\P\4\F\\{repeat}\?\S\?\\{loop}\par
\P\4\F\\{until}\?\S\?\\{end}\par
\fi

\M10. The following parameters are set big enough to handle \TeX, so they
should be sufficient for most applications of \.{ATANGLE}.

\Y\P\?\4\X10:Types and objects visible from \\{data\_structures}\X\S\?\6
\\{buf\_size}\?:\?\&{constant}\?\8\K\?100;\C{maximum length of input line}\6
\\{max\_bytes}\?:\?\&{constant}\?\8\K\?45\_000;\C{1\?/\?\\{ww} times the number
of bytes in    identifiers, strings, and module names; must be less than 65%
\_536}\6
\\{max\_toks}\?:\?\&{constant}\?\8\K\?50\_000;\C{1\?/\?\\{zz} times the number
of bytes in       compressed \ADA\ code; must be less than 65\_536}\6
\\{max\_names}\?:\?\&{constant}\?\8\K\?4\_000;\C{number of identifiers,
strings, module names;     must be less than 10\_240}\6
\\{max\_texts}\?:\?\&{constant}\?\8\K\?2\_000;\C{number of replacement texts,
must be less     than 10\_240}\6
\\{hash\_size}\?:\?\&{constant}\?\8\K\?353;\C{should be prime}\6
\\{longest\_name}\?:\?\&{constant}\?\8\K\?400;\C{module names shouldn't be
longer than this}\6
\\{line\_length}\?:\?\&{constant}\?\8\K\?72;\C{lines of \ADA\ output have at
most this many                                characters}\6
\\{out\_buf\_size}\?:\?\&{constant}\?\8\K\?144;\C{length of output buffer,
should be twice                                  \\{line\_length}}\6
\\{stack\_size}\?:\?\&{constant}\?\8\K\?50;\C{number of simultaneous levels of
macro                               expansion}\6
\\{max\_id\_length}\?:\?\&{constant}\?\8\K\?\\{line\_length};\C{long
identifiers are chopped to                             this length, which must
not exceed \\{line\_length}}\6
\\{unambig\_length}\?:\?\&{constant}\?\8\K\?\\{max\_id\_length};\C{identifiers
must be unique if                                              chopped to this
length}\par
\A sections~11, 12, 13, 14, 16, 22, 25, 27, 30, 35, 38, 39, 40, 42, 44, 45, 49,
62, 63, 70, 72, 76, 77, 79, 83, 91, 92, 97, 107, 109, 116, 118, 131, 136, 148,
156, and~165.
\U section~3\*.\fi

\M11. A global variable called \\{history} will contain one of four values
at the end of every run: \\{spotless} means that no unusual messages were
printed; \\{harmless\_message} means that a message of possible interest
was printed but no serious errors were detected; \\{error\_message} means that
at least one error was found; \\{fatal\_message} means that the program
terminated abnormally. The value of \\{history} does not influence the
behavior of the program; it is simply computed for the convenience
of systems that might want to use such information.

\Y\P\4\D\\{mark\_harmless}\?\S\?\6
\&{if}\1\1\8\\{history}\?=\?\\{spotless}\2\8\&{then}\6
\\{history}\?\K\?\\{harmless\_message};\2\6
\&{end}\8\2\&{if}\1\par
\P\4\D\\{mark\_error}\?\S\?\\{history}\?\K\?\\{error\_message}\par
\P\4\D\\{mark\_fatal}\?\S\?\\{history}\?\K\?\\{fatal\_message}\par
\Y\P\?\4\X10:Types and objects visible from \\{data\_structures}\X\mathrel{+}\S%
\?\6
\&{type}\8\\{history\_type}\8\&{is}\?\8\1(\?\\{spotless}\?,\39\?\\{harmless%
\_message}\?,\39\?\\{error\_message}\?,\39\?\\{fatal\_message}\?\2)\?;\6
\\{history}\?:\?\\{history\_type}\?\K\?\\{spotless};\C{how bad was this run?}%
\par
\fi

\N12.  The character set.
One of the main goals in the design of \.{AWEB} has been to make it readily
portable between a wide variety of computers. Yet \.{AWEB} by its very
nature must use a greater variety of characters than most computer
programs deal with, and character encoding is one of the areas in which
existing machines differ most widely from each other.

To resolve this problem, all input to \.{AWEAVE} and \.{ATANGLE} is converted
to an internal seven-bit code that is essentially standard ASCII, the
``American Standard Code for Information Interchange.''  The conversion is
done immediately when each character is read in. Conversely, characters are
converted from ASCII to the user's external representation just before
they are output.

Such an internal code is relevant to users of \.{AWEB} only because it is
the code used for preprocessed constants like \.{@'A'}. If you are
writing a program in \.{AWEB} that makes use of such one-character constants,
you should convert your input to ASCII form, like \.{AWEAVE} and \.{ATANGLE}
do. Otherwise \.{AWEB}'s internal coding scheme does not affect you.

Here is a table of the standard visible ASCII codes:
$$\def\:{\char\count255\global\advance\count255 by 1}
\count255='40
\vbox{
\hbox{\hbox to 40pt{\it\hfill0\/\hfill}%
\hbox to 40pt{\it\hfill1\/\hfill}%
\hbox to 40pt{\it\hfill2\/\hfill}%
\hbox to 40pt{\it\hfill3\/\hfill}%
\hbox to 40pt{\it\hfill4\/\hfill}%
\hbox to 40pt{\it\hfill5\/\hfill}%
\hbox to 40pt{\it\hfill6\/\hfill}%
\hbox to 40pt{\it\hfill7\/\hfill}}
\vskip 4pt
\hrule
\def\^{\vrule height 10.5pt depth 4.5pt}
\halign{\hbox to 0pt{\hskip -24pt\bn{8}{#0}\hfill}&\^
\hbox to 40pt{\tt\hfill#\hfill\^}&
&\hbox to 40pt{\tt\hfill#\hfill\^}\cr
04&\:&\:&\:&\:&\:&\:&\:&\:\cr\noalign{\hrule}
05&\:&\:&\:&\:&\:&\:&\:&\:\cr\noalign{\hrule}
06&\:&\:&\:&\:&\:&\:&\:&\:\cr\noalign{\hrule}
07&\:&\:&\:&\:&\:&\:&\:&\:\cr\noalign{\hrule}
10&\:&\:&\:&\:&\:&\:&\:&\:\cr\noalign{\hrule}
11&\:&\:&\:&\:&\:&\:&\:&\:\cr\noalign{\hrule}
12&\:&\:&\:&\:&\:&\:&\:&\:\cr\noalign{\hrule}
13&\:&\:&\:&\:&\:&\:&\:&\:\cr\noalign{\hrule}
14&\:&\:&\:&\:&\:&\:&\:&\:\cr\noalign{\hrule}
15&\:&\:&\:&\:&\:&\:&\:&\:\cr\noalign{\hrule}
16&\:&\:&\:&\:&\:&\:&\:&\:\cr\noalign{\hrule}
17&\:&\:&\:&\:&\:&\:&\:\cr}
\hrule width 280pt}$$
(Actually, of course, code \bn{8}{40} is an invisible blank space.)  Code
\bn{8}{136} was once an upward arrow (\.{\char'13}), and code \bn{8}{137} was
once a left arrow (\.^^X), in olden times when the first draft
of ASCII code was prepared; but \.{AWEB} works with today's standard
ASCII in which those codes represent circumflex and underline as shown.

\Y\P\?\4\X10:Types and objects visible from \\{data\_structures}\X\mathrel{+}\S%
\?\6
\&{subtype}\8\\{ASCII\_code}\8\&{is}\8\p{INTEGER}\8\&{range}\80\?\to\?127;%
\C{seven-bit numbers, a subrange
of the integers}\par
\fi

\M13. We shall use the name \\{text\_char} to stand for the data type of the
characters in the input and output files.  We shall also assume that
\\{text\_char} consists of the elements \\{text\_char}\?'\?\p{FIRST}
through  \\{text\_char}\?'\?\p{LAST}, inclusive. The following
definition should be adjusted if necessary.

\Y\P\?\4\X10:Types and objects visible from \\{data\_structures}\X\mathrel{+}\S%
\?\6
\&{subtype}\8\\{text\_char}\8\&{is}\8\p{CHARACTER};\C{the data type of
characters in text files}\6
\&{subtype}\8\\{text\_file}\8\&{is}\8\p{TEXT\_IO}\?.\?\p{FILE\_TYPE};\C{file
type consisting of \\{text\_char}}\par
\fi

\M14. The \.{AWEAVE} and \.{ATANGLE} processors convert between ASCII code and
the user's external character set by means of arrays \\{xord} and \\{xchr}
that are analogous to \ADA's \p{POS} and \p{VAL} attributes.

\Y\P\?\4\X10:Types and objects visible from \\{data\_structures}\X\mathrel{+}\S%
\?\6
\\{xord}\?:\1\&{array}\8\1(\?\\{text\_char}\?\2)\2\8\&{of}\?\8\\{ASCII\_code};%
\C{specifies conversion of input characters}\6
\\{xchr}\?:\1\&{array}\8\1(\?\\{ASCII\_code}\?\2)\2\8\&{of}\?\8\\{text\_char};%
\C{specifies conversion of output characters}\par
\fi

\M15. If we assume that every system using \.{AWEB} is able to read and write
the
visible characters of standard ASCII (although not necessarily using the
ASCII codes to represent them), the following assignment statements
initialize most of the \\{xchr} array properly, without needing any
system-dependent changes. For example, the statement
\.{xchr(8\#101\#):='A';} that appears in the present \.{AWEB} file might be
encoded in, say, {\mc EBCDIC} code on the external medium on which it
resides, but \.{ATANGLE} will convert from this external code to ASCII and
back again. Therefore the assignment statement \.{XCHR(8\#101\#):='A';}
will appear in the corresponding \ADA\ file, and \ADA\ will compile this
statement so that \\{xchr}\?(\bn{8}{101})\? receives  the character \.A in the
external
(\\{text\_char}) code. Note that it would be quite incorrect to say
\.{xchr(8\#101\#):=@'A';}, because \H{A} is a constant of type
\p{INTEGER}, not \p{CHARACTER}, and because we have \H{A}\?=\?65 regardless of
the
external character set.

\Y\P\?\4\X15:Set initial values\X\S\?\6
\\{xchr}\?\1(\bn{8}{40}\2)\K\?\.{\'\ \'};\5
\\{xchr}\?\1(\bn{8}{41}\2)\K\?\.{\'!\'};\5
\\{xchr}\?\1(\bn{8}{42}\2)\K\?\.{\'"\'};\5
\\{xchr}\?\1(\bn{8}{43}\2)\K\?\.{\'\#\'};\5
\\{xchr}\?\1(\bn{8}{44}\2)\K\?\.{\'\$\'};\5
\\{xchr}\?\1(\bn{8}{45}\2)\K\?\.{\'\%\'};\5
\\{xchr}\?\1(\bn{8}{46}\2)\K\?\.{\'\&\'};\5
\\{xchr}\?\1(\bn{8}{47}\2)\K\?\.{\'\'\'};\6
\\{xchr}\?\1(\bn{8}{50}\2)\K\?\.{\'(\'};\5
\\{xchr}\?\1(\bn{8}{51}\2)\K\?\.{\')\'};\5
\\{xchr}\?\1(\bn{8}{52}\2)\K\?\.{\'*\'};\5
\\{xchr}\?\1(\bn{8}{53}\2)\K\?\.{\'+\'};\5
\\{xchr}\?\1(\bn{8}{54}\2)\K\?\.{\',\'};\5
\\{xchr}\?\1(\bn{8}{55}\2)\K\?\.{\'-\'};\5
\\{xchr}\?\1(\bn{8}{56}\2)\K\?\.{\'.\'};\5
\\{xchr}\?\1(\bn{8}{57}\2)\K\?\.{\'/\'};\6
\\{xchr}\?\1(\bn{8}{60}\2)\K\?\.{\'0\'};\5
\\{xchr}\?\1(\bn{8}{61}\2)\K\?\.{\'1\'};\5
\\{xchr}\?\1(\bn{8}{62}\2)\K\?\.{\'2\'};\5
\\{xchr}\?\1(\bn{8}{63}\2)\K\?\.{\'3\'};\5
\\{xchr}\?\1(\bn{8}{64}\2)\K\?\.{\'4\'};\5
\\{xchr}\?\1(\bn{8}{65}\2)\K\?\.{\'5\'};\5
\\{xchr}\?\1(\bn{8}{66}\2)\K\?\.{\'6\'};\5
\\{xchr}\?\1(\bn{8}{67}\2)\K\?\.{\'7\'};\6
\\{xchr}\?\1(\bn{8}{70}\2)\K\?\.{\'8\'};\5
\\{xchr}\?\1(\bn{8}{71}\2)\K\?\.{\'9\'};\5
\\{xchr}\?\1(\bn{8}{72}\2)\K\?\.{\':\'};\5
\\{xchr}\?\1(\bn{8}{73}\2)\K\?\.{\';\'};\5
\\{xchr}\?\1(\bn{8}{74}\2)\K\?\.{\'<\'};\5
\\{xchr}\?\1(\bn{8}{75}\2)\K\?\.{\'=\'};\5
\\{xchr}\?\1(\bn{8}{76}\2)\K\?\.{\'>\'};\5
\\{xchr}\?\1(\bn{8}{77}\2)\K\?\.{\'?\'};\6
\\{xchr}\?\1(\bn{8}{100}\2)\K\?\.{\'@\'};\5
\\{xchr}\?\1(\bn{8}{101}\2)\K\?\.{\'A\'};\5
\\{xchr}\?\1(\bn{8}{102}\2)\K\?\.{\'B\'};\5
\\{xchr}\?\1(\bn{8}{103}\2)\K\?\.{\'C\'};\5
\\{xchr}\?\1(\bn{8}{104}\2)\K\?\.{\'D\'};\5
\\{xchr}\?\1(\bn{8}{105}\2)\K\?\.{\'E\'};\5
\\{xchr}\?\1(\bn{8}{106}\2)\K\?\.{\'F\'};\5
\\{xchr}\?\1(\bn{8}{107}\2)\K\?\.{\'G\'};\6
\\{xchr}\?\1(\bn{8}{110}\2)\K\?\.{\'H\'};\5
\\{xchr}\?\1(\bn{8}{111}\2)\K\?\.{\'I\'};\5
\\{xchr}\?\1(\bn{8}{112}\2)\K\?\.{\'J\'};\5
\\{xchr}\?\1(\bn{8}{113}\2)\K\?\.{\'K\'};\5
\\{xchr}\?\1(\bn{8}{114}\2)\K\?\.{\'L\'};\5
\\{xchr}\?\1(\bn{8}{115}\2)\K\?\.{\'M\'};\5
\\{xchr}\?\1(\bn{8}{116}\2)\K\?\.{\'N\'};\5
\\{xchr}\?\1(\bn{8}{117}\2)\K\?\.{\'O\'};\6
\\{xchr}\?\1(\bn{8}{120}\2)\K\?\.{\'P\'};\5
\\{xchr}\?\1(\bn{8}{121}\2)\K\?\.{\'Q\'};\5
\\{xchr}\?\1(\bn{8}{122}\2)\K\?\.{\'R\'};\5
\\{xchr}\?\1(\bn{8}{123}\2)\K\?\.{\'S\'};\5
\\{xchr}\?\1(\bn{8}{124}\2)\K\?\.{\'T\'};\5
\\{xchr}\?\1(\bn{8}{125}\2)\K\?\.{\'U\'};\5
\\{xchr}\?\1(\bn{8}{126}\2)\K\?\.{\'V\'};\5
\\{xchr}\?\1(\bn{8}{127}\2)\K\?\.{\'W\'};\6
\\{xchr}\?\1(\bn{8}{130}\2)\K\?\.{\'X\'};\5
\\{xchr}\?\1(\bn{8}{131}\2)\K\?\.{\'Y\'};\5
\\{xchr}\?\1(\bn{8}{132}\2)\K\?\.{\'Z\'};\5
\\{xchr}\?\1(\bn{8}{133}\2)\K\?\.{\'[\'};\5
\\{xchr}\?\1(\bn{8}{134}\2)\K\?\.{\'\\\'};\5
\\{xchr}\?\1(\bn{8}{135}\2)\K\?\.{\']\'};\5
\\{xchr}\?\1(\bn{8}{136}\2)\K\?\.{\'\^\'};\5
\\{xchr}\?\1(\bn{8}{137}\2)\K\?\.{\'\_\'};\6
\\{xchr}\?\1(\bn{8}{140}\2)\K\?\.{\'\`\'};\5
\\{xchr}\?\1(\bn{8}{141}\2)\K\?\.{\'a\'};\5
\\{xchr}\?\1(\bn{8}{142}\2)\K\?\.{\'b\'};\5
\\{xchr}\?\1(\bn{8}{143}\2)\K\?\.{\'c\'};\5
\\{xchr}\?\1(\bn{8}{144}\2)\K\?\.{\'d\'};\5
\\{xchr}\?\1(\bn{8}{145}\2)\K\?\.{\'e\'};\5
\\{xchr}\?\1(\bn{8}{146}\2)\K\?\.{\'f\'};\5
\\{xchr}\?\1(\bn{8}{147}\2)\K\?\.{\'g\'};\6
\\{xchr}\?\1(\bn{8}{150}\2)\K\?\.{\'h\'};\5
\\{xchr}\?\1(\bn{8}{151}\2)\K\?\.{\'i\'};\5
\\{xchr}\?\1(\bn{8}{152}\2)\K\?\.{\'j\'};\5
\\{xchr}\?\1(\bn{8}{153}\2)\K\?\.{\'k\'};\5
\\{xchr}\?\1(\bn{8}{154}\2)\K\?\.{\'l\'};\5
\\{xchr}\?\1(\bn{8}{155}\2)\K\?\.{\'m\'};\5
\\{xchr}\?\1(\bn{8}{156}\2)\K\?\.{\'n\'};\5
\\{xchr}\?\1(\bn{8}{157}\2)\K\?\.{\'o\'};\6
\\{xchr}\?\1(\bn{8}{160}\2)\K\?\.{\'p\'};\5
\\{xchr}\?\1(\bn{8}{161}\2)\K\?\.{\'q\'};\5
\\{xchr}\?\1(\bn{8}{162}\2)\K\?\.{\'r\'};\5
\\{xchr}\?\1(\bn{8}{163}\2)\K\?\.{\'s\'};\5
\\{xchr}\?\1(\bn{8}{164}\2)\K\?\.{\'t\'};\5
\\{xchr}\?\1(\bn{8}{165}\2)\K\?\.{\'u\'};\5
\\{xchr}\?\1(\bn{8}{166}\2)\K\?\.{\'v\'};\5
\\{xchr}\?\1(\bn{8}{167}\2)\K\?\.{\'w\'};\6
\\{xchr}\?\1(\bn{8}{170}\2)\K\?\.{\'x\'};\5
\\{xchr}\?\1(\bn{8}{171}\2)\K\?\.{\'y\'};\5
\\{xchr}\?\1(\bn{8}{172}\2)\K\?\.{\'z\'};\5
\\{xchr}\?\1(\bn{8}{173}\2)\K\?\.{\'\{\'};\5
\\{xchr}\?\1(\bn{8}{174}\2)\K\?\.{\'|\'};\5
\\{xchr}\?\1(\bn{8}{175}\2)\K\?\.{\'\}\'};\5
\\{xchr}\?\1(\bn{8}{176}\2)\K\?\.{\'\~\'};\6
\\{xchr}\?\1(\?0\?\2)\K\?\.{\'\ \'};\5
\\{xchr}\?\1(\bn{8}{177}\2)\K\?\.{\'\ \'};\C{these ASCII codes are not used}\par
\A sections~17, 18, 41, 43, 46, 71, and~142.
\U section~3\*.\fi

\M16. Some of the ASCII codes below \bn{8}{40} have been given symbolic names
in
\.{AWEAVE} and \.{ATANGLE} because they are used with a special meaning.

\Y\P\?\4\X10:Types and objects visible from \\{data\_structures}\X\mathrel{+}\S%
\?\6
\\{and\_sign}\?:\?\&{constant}\8\\{ASCII\_code}\?\K\bn{8}{4}\?;\C{equivalent to
`\.{and}'}\6
\\{not\_sign}\?:\?\&{constant}\8\\{ASCII\_code}\?\K\bn{8}{5}\?;\C{equivalent to
`\.{not}'}\6
\\{set\_element\_sign}\?:\?\&{constant}\8\\{ASCII\_code}\?\K\bn{8}{6}\?;%
\C{equivalent to `\.{in}'}\6
\\{tab\_mark}\?:\?\&{constant}\8\\{ASCII\_code}\?\K\bn{8}{11}\?;\C{ASCII code
used as tab-skip}\6
\\{line\_feed}\?:\?\&{constant}\8\\{ASCII\_code}\?\K\bn{8}{12}\?;\C{ASCII code
thrown away at end of line}\6
\\{form\_feed}\?:\?\&{constant}\8\\{ASCII\_code}\?\K\bn{8}{14}\?;\C{ASCII code
used at end of page}\6
\\{carriage\_return}\?:\?\&{constant}\8\\{ASCII\_code}\?\K\bn{8}{15}\?;\C{ASCII
code used at end of line}\6
\\{right\_arrow}\?:\?\&{constant}\8\\{ASCII\_code}\?\K\bn{8}{21}\?;%
\C{equivalent to `\.{=>}'}\6
\\{box\_sign}\?:\?\&{constant}\8\\{ASCII\_code}\?\K\bn{8}{22}\?;\C{equivalent
to `\.{<>}'}\6
\\{double\_star}\?:\?\&{constant}\8\\{ASCII\_code}\?\K\bn{8}{23}\?;%
\C{equivalent to `\.{**}'}\6
\\{left\_label\_bracket}\?:\?\&{constant}\8\\{ASCII\_code}\?\K\bn{8}{24}\?;%
\C{equivalent to `\.{<<}'}\6
\\{right\_label\_bracket}\?:\?\&{constant}\8\\{ASCII\_code}\?\K\bn{8}{25}\?;%
\C{equivalent to `\.{>>}'}\6
\\{left\_arrow}\?:\?\&{constant}\8\\{ASCII\_code}\?\K\bn{8}{30}\?;\C{equivalent
to `\.{:=}'}\6
\\{not\_equal}\?:\?\&{constant}\8\\{ASCII\_code}\?\K\bn{8}{32}\?;\C{equivalent
to `\.{/=}'}\6
\\{less\_or\_equal}\?:\?\&{constant}\8\\{ASCII\_code}\?\K\bn{8}{34}\?;%
\C{equivalent to `\.{<=}'}\6
\\{greater\_or\_equal}\?:\?\&{constant}\8\\{ASCII\_code}\?\K\bn{8}{35}\?;%
\C{equivalent to `\.{>=}'}\6
\\{equivalence\_sign}\?:\?\&{constant}\8\\{ASCII\_code}\?\K\bn{8}{36}\?;%
\C{equivalent to `\.{==}'}\6
\\{or\_sign}\?:\?\&{constant}\8\\{ASCII\_code}\?\K\bn{8}{37}\?;\C{equivalent to
`\.{or}'}\par
\fi

\M17. Here now is the system-dependent part of the character set.
If \.{AWEB} is being implemented on a garden-variety \ADA\ for which
only standard ASCII codes will appear in the input and output files, you
don't need to make any changes here.

Changes to the present module will make \.{AWEB} more friendly on computers
that have an extended character set, so that one can type things like
\.^^Z\ instead of \.{/=}. If you have an extended set of characters that
are easily incorporated into text files, you can assign codes arbitrarily
here, giving an \\{xchr} equivalent to whatever characters the users of
\.{AWEB} are allowed to have in their input files, provided that unsuitable
characters do not correspond to special codes like \\{carriage\_return}
that are listed above.

(The present file \.{ATANGLE.AWEB} does not contain any of the non-ASCII
characters, because it is intended to be used with all implementations of
\.{AWEB}.)

\Y\P\?\4\X15:Set initial values\X\mathrel{+}\S\?\6
\\{xchr}\?\1(\?1\?\to\bn{8}{37}\2)\K\1(\?\&{others}\?\8\KR\?\.{\'\ \'}\?\2)\?;%
\par
\fi

\M18. The following system-independent code makes the \\{xord} array contain a
suitable inverse to the information in \\{xchr}.

\Y\P\?\4\X15:Set initial values\X\mathrel{+}\S\?\6
\\{xord}\?\1(\?\\{text\_char}\?\2)\K\1(\?\&{others}\?\8\KR\bn{8}{40}\2)\?;\6
\1\&{for}\8\|i\?\in\?1\?\to\bn{8}{176}\?\8\&{loop}\6
\\{xord}\?\1(\?\\{xchr}\?\1(\?\|i\?\2)\2)\K\?\|i;\2\6
\&{end}\8\&{loop};\par
\fi

\N19.  Input and output.
The input conventions of this program are identical to those of \.{AWEAVE}.
Therefore people who need to make modifications to both systems
should be able to do so without too many headaches.

\fi

\M20. Terminal output is done by writing on file \p{STANDARD\_OUTPUT}, which is
assumed to consist of characters of type \\{text\_char}:

\Y\P\4\D\\{new\_ln}\?\S\?\p{TEXT\_IO}\?.\?\p{NEW\_LINE}\?\6
\?\par
\P\4\D\\{print}\?\1(\?\#\?\2)\S\?\p{TEXT\_IO}\?.\?\p{PUT}\?\1(\?\#\?\2)\C{`%
\\{print}' means write on the terminal}\6
\?\par
\P\4\D\\{print\_ln}\?\1(\?\#\?\2)\S\?\\{print}\?\1(\?\#\?\2)\?; \\{new\_ln}\?%
\C{`\\{print}' and then start new line}\6
\?\par
\P\4\D\\{print\_nl}\?\1(\?\#\?\2)\S\?\\{new\_ln}; \\{print}\?\1(\?\#\?\2)%
\C{print information starting on a
new line}\6
\?\par
\P\4\D\\{print\_int}\?\1(\?\#\?\2)\S\?\\{int\_number\_io}\?.\?\p{PUT}\?\1(\?\#%
\?\2)\C{print out an integer value}\6
\?\par
\fi

\M21. The \\{update\_terminal} procedure is called when we want
to make sure that everything we have output to the terminal so far has
actually left the computer's internal buffers and been sent.
(In the \ADA-system used to develop the initial program this is not
required.)

\Y\P\4\D\\{update\_terminal}\?\S\?\&{null}\?\C{empty the terminal output
buffer}\6
\?\par
\fi

\M22. The main input comes from \\{web\_file}; this input may be overridden
by changes in \\{change\_file}.
(If \\{change\_file} is empty, there are no changes.)

\Y\P\?\4\X10:Types and objects visible from \\{data\_structures}\X\mathrel{+}\S%
\?\6
\\{web\_file}\?:\?\\{text\_file};\C{primary input}\6
\\{change\_file}\?:\?\\{text\_file};\C{updates}\par
\fi

\M23. \P\?\X23:Specification of procedures visible from \\{input\_output}\X\S\?%
\6
\&{procedure}\8\1\\{open\_a\_file}\?(\?\|f\?:\&{in}\?\8\&{out}\8\\{text\_file};%
\?\,\35\?\\{mode\_of\_file}\?:\?\p{TEXT\_IO}\?.\?\p{FILE\_MODE};\?\,\35\?%
\\{default\_name}\?:\?\p{STRING};\?\,\35\1\?\\{success}\?:\?\&{out}\8%
\p{BOOLEAN}\?\2)\?\2;\6
\&{function}\8\1\\{create\_ada\_file}\?(\1\?\\{ada\_file\_name}\?:\?\p{STRING}%
\?\2)\&{return}\?\8\p{BOOLEAN}\2;\par
\A sections~28 and~31.
\U section~4\*.\fi

\M24. The following code opens a file and returns \p{TRUE} if succesfully.
If the file is already open the procedure has no effect; this gives
the ability to open the files by means of the change file, without
changing the \.{AWEB}-file of \.{ATANGLE}.


\Y\P\?\4\X24:Procedures in \\{input\_output}\X\S\?\6
\&{procedure}\8\1\\{open\_a\_file}\?(\?\|f\?:\&{in}\?\8\&{out}\8\\{text\_file};%
\?\,\35\?\\{mode\_of\_file}\?:\?\p{TEXT\_IO}\?.\?\p{FILE\_MODE};\?\,\35\?%
\\{default\_name}\?:\?\p{STRING};\?\,\35\1\?\\{success}\?:\?\&{out}\8%
\p{BOOLEAN}\?\2)\?\2\8\&{is}\8\1\6
\4\&{begin}\6
\&{if}\1\1\8\p{TEXT\_IO}\?.\?\p{IS\_OPEN}\?\1(\?\|f\?\2)\?\2\8\&{then}\6
\\{success}\?\K\?\p{TRUE};\6
\4\&{else}\6
\p{TEXT\_IO}\?.\?\p{OPEN}\?\1(\?\p{FILE}\?\KR\?\|f\?,\39\?\p{MODE}\?\KR\?%
\\{mode\_of\_file}\?,\39\?\p{NAME}\?\KR\?\\{default\_name}\?\2)\?;\5
\\{success}\?\K\?\p{TRUE};\2\6
\&{end}\8\2\&{if}\1;\6
\1\4\&{exception}\6
\4\&{when}\8\p{TEXT\_IO}\?.\?\p{STATUS\_ERROR}\?\mathrel{\.|}\?\p{TEXT\_IO}\?.%
\?\p{NAME\_ERROR}\?\mathrel{\.|}\?\p{TEXT\_IO}\?.\?\p{USE\_ERROR}\KR\6
\\{success}\?\K\?\p{FALSE};\5
\&{return};\2\6
\&{end}\8\\{open\_a\_file};\par
\A sections~26, 29, and~32.
\U section~4\*.\fi

\M25. The main output goes to \\{ada\_file}.

\Y\P\?\4\X10:Types and objects visible from \\{data\_structures}\X\mathrel{+}\S%
\?\6
\\{ada\_file}\?:\?\\{text\_file};\par
\fi

\M26. The following code creates \\{ada\_file}.

\Y\P\?\4\X24:Procedures in \\{input\_output}\X\mathrel{+}\S\?\6
\&{function}\8\1\\{create\_ada\_file}\?(\1\?\\{ada\_file\_name}\?:\?\p{STRING}%
\?\2)\&{return}\?\8\p{BOOLEAN}\2\8\&{is}\8\1\6
\4\&{begin}\6
\&{if}\1\1\8\p{TEXT\_IO}\?.\?\p{IS\_OPEN}\?\1(\?\\{ada\_file}\?\2)\?\2\8%
\&{then}\6
\&{if}\1\1\8\p{TEXT\_IO}\?.\?\p{NAME}\?\1(\?\\{ada\_file}\?\2)=\?\\{ada\_file%
\_name}\2\8\&{then}\6
\?\&{return}\?\8\p{TRUE};\6
\4\&{else}\6
\?\&{return}\?\8\p{FALSE};\2\6
\&{end}\8\2\&{if}\1;\6
\4\&{else}\6
\p{TEXT\_IO}\?.\?\p{CREATE}\?\1(\?\p{FILE}\?\KR\?\\{ada\_file}\?,\39\?\p{MODE}%
\?\KR\?\p{TEXT\_IO}\?.\?\p{OUT\_FILE}\?,\39\?\p{NAME}\?\KR\?\\{ada\_file\_name}%
\?\2)\?;\5
\?\&{return}\?\8\p{TRUE};\2\6
\&{end}\8\2\&{if}\1;\6
\1\4\&{exception}\6
\4\&{when}\8\p{TEXT\_IO}\?.\?\p{STATUS\_ERROR}\?\mathrel{\.|}\?\p{TEXT\_IO}\?.%
\?\p{NAME\_ERROR}\?\mathrel{\.|}\?\p{TEXT\_IO}\?.\?\p{USE\_ERROR}\KR\6
\?\&{return}\?\8\p{FALSE};\2\6
\&{end}\8\\{create\_ada\_file};\par
\fi

\M27. Input goes into an array called \\{buffer}.
\Y\P\?\4\X10:Types and objects visible from \\{data\_structures}\X\mathrel{+}\S%
\?\6
\&{type}\8\\{buffer\_type}\8\&{is}\?\8\1\&{array}\8\1(\?\p{INTEGER}\8\&{range}%
\?\8\BOX\2)\2\8\&{of}\?\8\\{ASCII\_code};\6
\&{subtype}\8\\{buffer\_index}\8\&{is}\8\p{INTEGER}\8\&{range}\80\?\to\?\\{buf%
\_size};\6
\\{buffer}\?:\?\\{buffer\_type}\?\1(\?\\{buffer\_index}\?\2)\?;\par
\fi

\M28. The \\{input\_ln} procedure brings the next line of input from the
specified
file into the \\{buffer} array and returns the value \p{TRUE}, unless the file
has already been entirely read, in which case it returns \p{FALSE}. The
conventions of \TeX\ are followed; i.e., \\{ASCII\_code} numbers representing
the next line of the file are input into \\{buffer}\?(\?0\?)\?, \\{buffer}\?(%
\?1\?)\?, \dots,
\\{buffer}\?(\?\\{limit}\?-\?1\?)\?; trailing blanks are ignored;
and the global variable \\{limit} is set to the length of the
line. The value of \\{limit} must be strictly less than \\{buf\_size}.

We assume that none of the \\{ASCII\_code} values
of \\{buffer}\?(\?\|j\?)\? for 0\?\L\?\|j\?<\?\\{limit} is equal to 0, %
\bn{8}{177}, \\{line\_feed},
\\{form\_feed}, or \\{carriage\_return}.

\Y\P\?\4\X23:Specification of procedures visible from \\{input\_output}\X%
\mathrel{+}\S\?\6
\&{function}\8\1\\{input\_ln}\?(\1\?\|f\?:\?\\{text\_file}\?\2)\&{return}\?\8%
\p{BOOLEAN}\2;\par
\fi

\M29. \P\?\X24:Procedures in \\{input\_output}\X\mathrel{+}\S\?\6
\&{function}\8\1\\{input\_ln}\?(\1\?\|f\?:\?\\{text\_file}\?\2)\&{return}\?\8%
\p{BOOLEAN}\2\8\&{is}\8\C{inputs a line or returns \p{FALSE}}\1\6
\\{final\_limit}\?:\?\\{buffer\_index}\?\K\?0;\C{\\{limit} without trailing
blanks}\6
\|c\?:\?\\{text\_char};\6
\4\&{begin}\6
\\{limit}\?\K\?0;\6
\&{if}\1\1\?\8\R\?\p{TEXT\_IO}\?.\?\p{IS\_OPEN}\?\1(\?\|f\?\2)\ORE\?\p{TEXT%
\_IO}\?.\?\p{END\_OF\_FILE}\?\1(\?\|f\?\2)\?\2\8\&{then}\6
\?\&{return}\?\8\p{FALSE};\6
\4\&{else}\6
\1\&{while}\?\8\R\?\p{TEXT\_IO}\?.\?\p{END\_OF\_LINE}\?\1(\?\|f\?\2)\?\8%
\&{loop}\6
\p{TEXT\_IO}\?.\?\p{GET}\?\1(\?\|f\?,\39\?\|c\?\2)\?;\5
\\{buffer}\?\1(\?\\{limit}\?\2)\K\?\\{xord}\?\1(\?\|c\?\2)\?;\5
\\{incr}\?\1(\?\\{limit}\?\2)\?;\6
\&{if}\1\1\8\\{buffer}\?\1(\?\\{limit}\?-\?1\?\2)\I\?\H{\ }\2\8\&{then}\6
\\{final\_limit}\?\K\?\\{limit};\2\6
\&{end}\8\2\&{if}\1;\6
\&{if}\1\1\8\\{limit}\?=\?\\{buf\_size}\2\8\&{then}\6
\p{TEXT\_IO}\?.\?\p{SKIP\_LINE}\?\1(\?\|f\?\2)\?;\5
\\{decr}\?\1(\?\\{limit}\?\2)\?;\C{keep \\{buffer}\?(\?\\{buf\_size}\?)\?
empty}\6
\\{print\_nl}\?\1(\?\.{"!\ Input\ line\ too\ long"}\?\2)\?;\5
\\{loc}\?\K\?0;\5
\\{error};\5
\?\&{return}\?\8\p{TRUE};\2\6
\&{end}\8\2\&{if}\1;\2\6
\&{end}\8\&{loop};\6
\p{TEXT\_IO}\?.\?\p{SKIP\_LINE}\?\1(\?\|f\?\2)\?;\5
\\{limit}\?\K\?\\{final\_limit};\5
\?\&{return}\?\8\p{TRUE};\2\6
\&{end}\8\2\&{if}\1;\2\6
\&{end}\8\\{input\_ln};\par
\fi

\N30.  Reporting errors to the user.
The \.{ATANGLE} processor operates in two phases: first it inputs the source
file and stores a compressed representation of the program, then it produces
the \ADA\ output from the compressed representation.

The global variable \\{phase\_one} tells whether we are in Phase I or not.

\Y\P\?\4\X10:Types and objects visible from \\{data\_structures}\X\mathrel{+}\S%
\?\6
\\{phase\_one}\?:\?\p{BOOLEAN};\C{\p{TRUE} in Phase I, \p{FALSE} in Phase II}%
\par
\fi

\M31. During the first phase, syntax errors are reported to the user by saying
$$\hbox{`\\{err\_print}\?(\?\.{"!\ Error\ message"}\?)\?'},$$
followed by `\\{jump\_out}' if no recovery from the error is provided.
This will print the error message followed by an indication of where the
error was spotted in the source file. Note that no period follows the error
message, since the error routine will automatically supply a period.

Errors that are noticed during the second phase are reported to the user
in the same fashion, but the error message will be
followed by an indication of where the error was spotted in the output file.

The actual error indications are provided by a procedure called \\{error}.

\Y\P\4\D\\{err\_print}\?\1(\?\#\?\2)\S\?\\{print\_nl}\?\1(\?\#\?\2)\?;\5
\\{error}\par
\Y\P\?\4\X23:Specification of procedures visible from \\{input\_output}\X%
\mathrel{+}\S\?\6
\&{procedure}\8\1\\{error}\2;\par
\fi

\M32. \P\?\X24:Procedures in \\{input\_output}\X\mathrel{+}\S\?\6
\&{procedure}\8\1\\{error}\2\8\&{is}\8\C{prints '\..' and location of error
message}\1\6
\|l\?:\?\\{buffer\_index};\C{index into \\{buffer}}\6
\4\&{begin}\6
\&{if}\1\1\8\\{phase\_one}\2\8\&{then}\6
\X33:Print error location based on input buffer\X;\6
\4\&{else}\6
\X34:Print error location based on output buffer\X;\2\6
\&{end}\8\2\&{if}\1;\6
\\{update\_terminal};\5
\\{mark\_error};\2\6
\&{end}\8\\{error};\par
\fi

\M33. The error locations during Phase I can be indicated by using the global
variables \\{loc}, \\{line}, and \\{changing}, which tell respectively the
first
unlooked-at position in \\{buffer}, the current line number, and whether or not
the current line is from \\{change\_file} or \\{web\_file}.
This routine should be modified on systems whose standard text editor
has special line-numbering conventions.

\Y\P\?\4\X33:Print error location based on input buffer\X\S\?\6
\&{if}\1\1\8\\{changing}\2\8\&{then}\6
\\{print}\?\1(\?\.{".\ (change\ file\ "}\?\2)\?;\6
\4\&{else}\6
\\{print}\?\1(\?\.{".\ ("}\?\2)\?;\2\6
\&{end}\8\2\&{if}\1;\6
\\{print}\?\1(\?\.{"l."}\?\2)\?;\5
\\{print\_int}\?\1(\?\\{line}\?,\39\?1\?\2)\?;\5
\\{print\_ln}\?\1(\?\.{")"}\?\2)\?;\6
\&{if}\1\1\8\\{loc}\?\G\?\\{limit}\2\8\&{then}\6
\|l\?\K\?\\{limit};\6
\4\&{else}\6
\|l\?\K\?\\{loc};\2\6
\&{end}\8\2\&{if}\1;\6
\1\&{for}\8\|k\?\in\?1\?\to\?\|l\8\&{loop}\6
\&{if}\1\1\8\\{buffer}\?\1(\?\|k\?-\?1\?\2)=\?\\{tab\_mark}\2\8\&{then}\6
\\{print}\?\1(\?\.{"\ "}\?\2)\?;\6
\4\&{else}\6
\\{print}\?\1(\?\\{xchr}\?\1(\?\\{buffer}\?\1(\?\|k\?-\?1\?\2)\2)\2)\?;\C{print
the characters already read}\2\6
\&{end}\8\2\&{if}\1;\2\6
\&{end}\8\&{loop};\6
\\{new\_ln};\6
\1\&{for}\8\|k\?\in\?1\?\to\?\|l\8\&{loop}\6
\\{print}\?\1(\?\.{"\ "}\?\2)\?;\2\6
\&{end}\8\&{loop};\C{space out the next line}\6
\1\&{for}\8\|k\?\in\?\|l\?+\?1\?\to\?\\{limit}\8\&{loop}\6
\\{print}\?\1(\?\\{xchr}\?\1(\?\\{buffer}\?\1(\?\|k\?-\?1\?\2)\2)\2)\?;\2\6
\&{end}\8\&{loop};\C{print the part not yet read}\6
\\{print}\?\1(\?\.{"\ "}\?\2)\?\C{this space separates the message from future
asterisks}\par
\U section~32.\fi

\M34. The position of errors detected during the second phase can be indicated
by outputting the partially-filled output buffer, which contains \\{out\_ptr}
entries. The current output file name is determined by \\{file\_name} and
\\{fn\_length}, whereas the line number comes from \\{f\_line}.

\Y\P\?\4\X34:Print error location based on output buffer\X\S\?\6
\\{print}\?\1(\?\.{".("}\?\mathbin{\.\&}\?\\{file\_name}\?\1(\?1\?\to\?\\{fn%
\_length}\?\2)\mathbin{\.\&}\?\.{"\ l."}\?\2)\?;\5
\\{print\_int}\?\1(\?\\{f\_line}\?,\39\?1\?\2)\?;\5
\\{print\_ln}\?\1(\?\.{")"}\?\2)\?;\6
\1\&{for}\8\|j\?\in\?1\?\to\?\\{out\_ptr}\8\&{loop}\6
\\{print}\?\1(\?\\{xchr}\?\1(\?\\{out\_buf}\?\1(\?\|j\?-\?1\?\2)\2)\2)\?;%
\C{print current partial line}\2\6
\&{end}\8\&{loop};\6
\\{print}\?\1(\?\.{"...\ "}\?\2)\?\C{indicate that this information is partial}%
\par
\U section~32.\fi

\M35. The \\{jump\_out} macro just raises the \\{end\_of\_ATANGLE} exception.
It is used when no recovery from a particular error has
been provided.

\Y\P\4\D\\{jump\_out}\?\S\?\&{raise}\8\\{end\_of\_ATANGLE}\par
\P\4\D\\{fatal\_error}\?\1(\?\#\?\2)\S\?\\{print\_nl}\?\1(\?\#\?\2)\?;\5
\\{error};\5
\\{mark\_fatal};\5
\\{jump\_out}\par
\Y\P\?\4\X10:Types and objects visible from \\{data\_structures}\X\mathrel{+}\S%
\?\6
\\{end\_of\_ATANGLE}\?:\&{exception}\?;\par
\fi

\M36. Sometimes the program's behavior is far different from what it should be,
and \.{ATANGLE} prints an error message that is really for the \.{ATANGLE}
maintenance person, not the user. In such cases the program says
\\{confusion}\?(\?\.{"indication\ of\ where\ we\ are"}\?)\?.

\Y\P\4\D\\{confusion}\?\1(\?\#\?\2)\S\?\\{fatal\_error}\?\1(\?\.{"!\ This\ can%
\'t\ happen\ ("}\?\mathbin{\.\&}\?\#\?\mathbin{\.\&}\?\.{")"}\?\2)\6
\?\par
\fi

\M37. An overflow stop occurs if \.{ATANGLE}'s tables aren't large enough.

\Y\P\4\D\\{overflow}\?\1(\?\#\?\2)\S\?\\{fatal\_error}\?\1(\?\.{"!\ Sorry,\ "}%
\?\mathbin{\.\&}\?\#\?\mathbin{\.\&}\?\.{"\ capacity\ exceeded"}\?\2)\6
\?\par
\fi

\N38.  Data structures.
Most of the user's \ADA\ code is packed into seven- or eight-bit integers
in two large arrays called \\{byte\_mem} and \\{tok\_mem}.
The \\{byte\_mem} array holds the names of identifiers, strings, and modules;
the \\{tok\_mem} array holds the replacement texts
for macros and modules. Allocation is sequential, since things are deleted
only during Phase II, and only in a last-in-first-out manner.

Auxiliary arrays \\{byte\_start} and \\{tok\_start} are used as directories to
\\{byte\_mem} and \\{tok\_mem}, and the \\{link}, \\{ilk}, \\{equiv}, and %
\\{text\_link}
arrays give further information about names. These auxiliary arrays
consist of sixteen-bit items.

\Y\P\?\4\X10:Types and objects visible from \\{data\_structures}\X\mathrel{+}\S%
\?\6
\&{subtype}\8\\{eight\_bits}\8\&{is}\8\p{INTEGER}\8\&{range}\80\?\to\?255;%
\C{unsigned one-byte quantity}\6
\&{subtype}\8\\{sixteen\_bits}\8\&{is}\8\p{INTEGER}\8\&{range}\80\?\to\?65%
\_535;\C{unsigned two-byte
quantity}\par
\fi

\M39. \.{ATANGLE} has been designed to avoid the need for indices that are more
than sixteen bits wide, so that it can be used on most computers. But
there are programs that need more than 65536 tokens, and some programs
even need more than 65536 bytes; \TeX\ is one of these.  To get around
this problem, a slight complication has been added to the data structures:
\\{byte\_mem} and \\{tok\_mem} are two-dimensional arrays, whose first index is
either 0 or 1. (For generality, the first index is actually allowed to run
between 0 and \\{ww}\?-\?1 in \\{byte\_mem}, or between 0 and \\{zz}\?-\?1 in %
\\{tok\_mem},
where \\{ww} and \\{zz} are set to 2 and~3; the program will work for any
positive values of \\{ww} and \\{zz}, and it can be simplified in obvious ways
if \\{ww}\?=\?1 or \\{zz}\?=\?1.)

\Y\P\?\4\X10:Types and objects visible from \\{data\_structures}\X\mathrel{+}\S%
\?\6
\\{ww}\?:\?\&{constant}\8\p{INTEGER}\?\K\?2;\6
\\{zz}\?:\?\&{constant}\8\p{INTEGER}\?\K\?3;\7
\&{subtype}\8\\{w\_range}\8\&{is}\8\p{INTEGER}\8\&{range}\80\?\to\?\\{ww}\?-%
\?1;\6
\&{subtype}\8\\{z\_range}\8\&{is}\8\p{INTEGER}\8\&{range}\80\?\to\?\\{zz}\?-%
\?1;\6
\&{subtype}\8\\{byte\_index}\8\&{is}\8\p{INTEGER}\8\&{range}\80\?\to\?\\{max%
\_bytes};\6
\&{subtype}\8\\{token\_index}\8\&{is}\8\p{INTEGER}\8\&{range}\80\?\to\?\\{max%
\_toks};\6
\&{subtype}\8\\{name\_index}\8\&{is}\8\p{INTEGER}\8\&{range}\80\?\to\?\\{max%
\_names};\6
\&{subtype}\8\\{text\_index}\8\&{is}\8\p{INTEGER}\8\&{range}\80\?\to\?\\{max%
\_texts};\7
\&{type}\8\\{byte\_memory}\8\&{is}\?\8\1\&{array}\8\1(\?\\{w\_range}\?,\39\?%
\\{byte\_index}\?\2)\2\8\&{of}\?\8\\{ASCII\_code};\6
\&{optim}\1\6
\&{pragma}\8\1\p{PACK}\?\1(\?\\{byte\_memory}\?\2)\?\2;\2\6
\&{mitpo}\7
\&{type}\8\\{token\_memory}\8\&{is}\?\8\1\&{array}\8\1(\?\\{z\_range}\?,\39\?%
\\{token\_index}\?\2)\2\8\&{of}\?\8\\{eight\_bits};\6
\&{optim}\1\6
\&{pragma}\8\1\p{PACK}\?\1(\?\\{token\_memory}\?\2)\?\2;\2\6
\&{mitpo}\7
\&{type}\8\\{name\_info}\8\&{is}\?\8\1\&{array}\8\1(\?\\{name\_index}\?\2)\2\8%
\&{of}\?\8\\{sixteen\_bits};\6
\&{type}\8\\{token\_info}\8\&{is}\?\8\1\&{array}\8\1(\?\\{text\_index}\?\2)\2\8%
\&{of}\?\8\\{sixteen\_bits};\7
\\{byte\_mem}\?:\?\\{byte\_memory};\C{characters of names}\6
\\{tok\_mem}\?:\?\\{token\_memory};\C{tokens}\6
\\{byte\_start}\?:\?\\{name\_info};\C{directory into \\{byte\_mem}}\6
\\{tok\_start}\?:\?\\{token\_info};\C{directory into \\{tok\_mem}}\6
\\{link}\?:\?\\{name\_info};\C{hash table or tree links}\6
\\{ilk}\?:\?\\{name\_info};\C{type codes or tree links}\6
\\{equiv}\?:\?\\{name\_info};\C{info corresponding to names}\6
\\{text\_link}\?:\?\\{token\_info};\C{relates replacement texts}\par
\fi

\M40. The names of identifiers are found by computing a hash address \|h and
then looking at strings of bytes signified by \\{hash}\?(\?\|h\?)\?, \\{link}%
\?(\?\\{hash}\?(\?\|h\?))\?,
\\{link}\?(\?\\{link}\?(\?\\{hash}\?(\?\|h\?)))\?, \dots, until either finding
the desired name
or encountering a zero.

A `\\{name\_pointer}' variable, which signifies a name, is an index into
\\{byte\_start}. The actual sequence of characters in the name pointed to by
\|p appears in positions \\{byte\_start}\?(\?\|p\?)\? to \\{byte\_start}\?(\?%
\|p\?+\?\\{ww}\?)-\?1, inclusive,
in the segment of \\{byte\_mem} whose first index is \|p\?\mathbin{\&{mod}}\?%
\\{ww}. Thus, when
\\{ww}\?=\?2 the even-numbered name bytes appear in \\{byte\_mem}\?(\?0\?,\?%
\hbox{$*$}\?)\?
and the odd-numbered ones appear in \\{byte\_mem}\?(\?1\?,\?\hbox{$*$}\?)\?.
The pointer 0 is used for undefined module names; we don't
want to use it for the names of identifiers, since 0 stands for a null
pointer in a linked list.

The total number of names in \\{byte\_mem}
is called \\{name\_ptr} and the total number of bytes occupied in
\\{byte\_mem}\?(\?\|w\?,\?\hbox{$*$}\?)\? is called \\{byte\_ptr}\?(\?\|w\?)\?.

We usually have \\{byte\_start}\?(\?\\{name\_ptr}\?+\?\|w\?)=\?\\{byte\_ptr}%
\?((\?\\{name\_ptr}\?+\?\|w\?)\mathbin{\&{mod}}\?\\{ww}\?)\?
for 0\?\L\?\|w\?<\?\\{ww}, since these are the starting positions for the next %
\\{ww}
names to be stored in \\{byte\_mem}.

\Y\P\4\D\\{length}\?\1(\?\#\?\2)\S\?\\{byte\_start}\?\1(\?\#\?+\?\\{ww}\?\2)-\?%
\\{byte\_start}\?\1(\?\#\?\2)\C{the length of a name}\6
\?\par
\P\4\D\\{null\_array}\?\S\1(\?\&{others}\?\8\KR\?0\?\2)\6
\?\par
\Y\P\?\4\X10:Types and objects visible from \\{data\_structures}\X\mathrel{+}\S%
\?\6
\&{subtype}\8\\{name\_pointer}\8\&{is}\8\\{name\_index};\C{identifies a name}\7
\\{name\_ptr}\?:\?\\{name\_pointer}\?\K\?1;\C{first unused position in \\{byte%
\_start}}\6
\\{byte\_ptr}\?:\1\&{array}\8\1(\?\\{w\_range}\?\2)\2\8\&{of}\?\8\\{byte%
\_index}\?\K\?\\{null\_array};\C{first unused position in \\{byte\_mem}}\par
\fi

\M41. \P\?\X15:Set initial values\X\mathrel{+}\S\?\6
\\{byte\_start}\?\1(\?0\?\to\?\\{ww}\?\2)\K\?\\{null\_array};\C{this makes name
0 of length zero}\par
\fi

\M42. Replacement texts are stored in \\{tok\_mem}, using similar conventions.
A `\\{text\_pointer}' variable is an index into \\{tok\_start}, and the
replacement text that corresponds to \|p runs from positions
\\{tok\_start}\?(\?\|p\?)\? to \\{tok\_start}\?(\?\|p\?+\?\\{zz}\?)-\?1,
inclusive, in the segment of
\\{tok\_mem} whose first index is \|p\?\mathbin{\&{mod}}\?\\{zz}. Thus, when %
\\{zz}\?=\?2 the
even-numbered replacement texts appear in \\{tok\_mem}\?(\?0\?,\?\hbox{$*$}\?)%
\? and the
odd-numbered ones appear in \\{tok\_mem}\?(\?1\?,\?\hbox{$*$}\?)\?.
Furthermore,
\\{text\_link}\?(\?\|p\?)\? is used to connect pieces of text that have the
same name,
as we shall see later. The pointer 0 is used for undefined replacement
texts.

The first position of \\{tok\_mem}\?(\?\|z\?,\?\hbox{$*$}\?)\? that is
unoccupied by
replacement text is called \\{tok\_ptr}\?(\?\|z\?)\?, and the first unused
location of
\\{tok\_start} is called \\{text\_ptr}.  We usually have the identity
\\{tok\_start}\?(\?\\{text\_ptr}\?+\?\|z\?)=\?\\{tok\_ptr}\?((\?\\{text\_ptr}%
\?+\?\|z\?)\mathbin{\&{mod}}\?\\{zz}\?)\?, for 0\?\L\?\|z\?<\?\\{zz}, since
these are the starting positions for the next \\{zz} replacement texts to
be stored in \\{tok\_mem}.

It is convenient to maintain a variable \|z that is equal to \\{text\_ptr}\?%
\mathbin{\&{mod}}\?\\{zz}, so that we always insert tokens into segment \|z of %
\\{tok\_mem}.

\Y\P\?\4\X10:Types and objects visible from \\{data\_structures}\X\mathrel{+}\S%
\?\6
\&{subtype}\8\\{text\_pointer}\8\&{is}\8\\{text\_index};\C{identifies a
replacement text}\7
\\{text\_ptr}\?:\?\\{text\_pointer}\?\K\?1;\C{first unused position in \\{tok%
\_start}}\6
\\{tok\_ptr}\?:\1\&{array}\8\1(\?\\{z\_range}\?\2)\2\8\&{of}\?\8\\{token%
\_index}\?\K\?\\{null\_array};\C{first unused position in a given segment of %
\\{tok\_mem}}\6
\|z\?:\?\\{z\_range}\?\K\?1\?\mathbin{\&{mod}}\?\\{zz};\C{current segment of %
\\{tok\_mem}}\6
\&{stat}\1\6
\\{max\_tok\_ptr}\?:\1\&{array}\8\1(\?\\{z\_range}\?\2)\2\8\&{of}\?\8\\{token%
\_index};\C{largest values assumed by \\{tok\_ptr}}\2\6
\&{tats}\par
\fi

\M43. \P\?\X15:Set initial values\X\mathrel{+}\S\?\6
\\{tok\_start}\?\1(\?0\?\to\?\\{zz}\?\2)\K\?\\{null\_array};\C{this makes
replacement text 0 of length zero}\par
\fi

\M44. Three types of identifiers are distinguished by their \\{ilk}:

\yskip\hang \\{normal} identifiers will appear in the \ADA\ program as
ordinary identifiers since they have not been defined to be macros; the
corresponding value in the \\{equiv} array for such identifiers is a link in a
secondary hash table that is used to check whether any two of them agree in
their first \\{unambig\_length} characters after letters are changed to
uppercase.

\yskip\hang \\{simple} identifiers have been defined to be simple macros;
their \\{equiv} value points to the corresponding replacement text.

\yskip\hang \\{parametric} identifiers have been defined to be parametric
macros; like simple identifiers, their \\{equiv} value points to the
replacement text.

\Y\P\?\4\X10:Types and objects visible from \\{data\_structures}\X\mathrel{+}\S%
\?\6
\\{normal}\?:\?\&{constant}\8\\{sixteen\_bits}\?\K\?0;\C{ordinary identifiers
have \\{normal} ilk}\6
\\{simple}\?:\?\&{constant}\8\\{sixteen\_bits}\?\K\?1;\C{simple macros have %
\\{simple} ilk}\6
\\{parametric}\?:\?\&{constant}\8\\{sixteen\_bits}\?\K\?2;\C{parametric macros
have \\{parametric} ilk}\par
\fi

\M45. The names of modules are stored in \\{byte\_mem} together with the
identifier
names, but a hash table is not used for them because \.{ATANGLE} needs to be
able to recognize a module name when given a prefix of that name. A
conventional binary seach tree is used to retrieve module names, with fields
called \\{llink} and \\{rlink} in place of \\{link} and \\{ilk}. The root of
this
tree is \\{rlink}\?(\?0\?)\?. If \|p is a pointer to a module name, \\{equiv}%
\?(\?\|p\?)\? points
to its replacement text, just as in simple and parametric macros, unless this
replacement text has not yet been defined (in which case \\{equiv}\?(\?\|p\?)=%
\?0).

\Y\P\?\4\X10:Types and objects visible from \\{data\_structures}\X\mathrel{+}\S%
\?\6
\\{llink}\?:\?\\{name\_info}\8\&{renames}\8\\{link};\C{left link in binary
search tree for module names}\6
\\{rlink}\?:\?\\{name\_info}\8\&{renames}\8\\{ilk};\C{right link in binary
search tree for module names}\par
\fi

\M46. \P\?\X15:Set initial values\X\mathrel{+}\S\?\6
\\{rlink}\?\1(\?0\?\2)\K\?0;\C{the binary search tree starts out with nothing
in it}\6
\\{equiv}\?\1(\?0\?\2)\K\?0;\C{the undefined module has no replacement text}\par
\fi

\M47. Here is a little procedure that prints the text of a given name.
\Y\P\?\4\X47:Specification of procedures visible from \\{data\_structures}\X\S%
\?\6
\&{procedure}\8\1\\{print\_id}\?(\1\?\|p\?:\?\\{name\_pointer}\?\2)\?\2;\par
\A section~73.
\U section~3\*.\fi

\M48. \P\?\X48:Procedures in \\{data\_structures}\X\S\?\6
\&{procedure}\8\1\\{print\_id}\?(\1\?\|p\?:\?\\{name\_pointer}\?\2)\?\2\8\&{is}%
\8\C{print identifier or module name}\1\6
\|w\?:\?\\{w\_range};\C{segment of \\{byte\_mem}}\6
\4\&{begin}\6
\&{if}\1\1\8\|p\?\G\?\\{name\_ptr}\2\8\&{then}\6
\\{print}\?\1(\?\.{"IMPOSSIBLE"}\?\2)\?;\6
\4\&{else}\6
\|w\?\K\?\|p\?\mathbin{\&{mod}}\?\\{ww};\6
\1\&{for}\8\|k\?\in\?\\{byte\_start}\?\1(\?\|p\?\2)\to\?\\{byte\_start}\?\1(\?%
\|p\?+\?\\{ww}\?\2)-\?1\8\&{loop}\6
\\{print}\?\1(\?\\{xchr}\?\1(\?\\{byte\_mem}\?\1(\?\|w\?,\39\?\|k\?\2)\2)\2)\?;%
\2\6
\&{end}\8\&{loop};\2\6
\&{end}\8\2\&{if}\1;\2\6
\&{end}\8\\{print\_id};\par
\A section~74.
\U section~3\*.\fi

\N49.  Searching for identifiers.
The hash table described above is updated by the \\{id\_lookup} procedure,
which finds a given identifier and returns a pointer to its index in
\\{byte\_start}. If the identifier was not already present, it is inserted with
a given \\{ilk} code; and an error message is printed if the identifier is
being doubly defined.

Because of the way \.{ATANGLE}'s scanning mechanism works, it is most
convenient to let \\{id\_lookup} search for an identifier that is present in
the
\\{buffer} array. Two other global variables specify its position in the
buffer: the first character is \\{buffer}\?(\?\\{id\_first}\?)\?, and the last
is
\\{buffer}\?(\?\\{id\_loc}\?-\?1\?)\?.

We have mentioned that \\{normal} identifiers belong to two hash tables, one
for their true names as they appear in the \.{AWEB} file and the other when
they have been reduced to their first \\{unambig\_length} characters. The hash
tables are kept by the method of simple chaining, where the heads of the
individual lists appear in the \\{hash} and \\{chop\_hash} arrays. If \|h is a
hash code, the primary hash table list starts at \\{hash}\?(\?\|h\?)\? and
proceeds
through \\{link} pointers; the secondary hash table list starts at
\\{chop\_hash}\?(\?\|h\?)\? and proceeds through \\{equiv} pointers. Of course,
the same
identifier will probably have two different values of \|h.

The \\{id\_lookup} procedure uses an auxiliary array called \\{chopped\_id} to
contain up to \\{unambig\_length} characters of the current identifier, if
it is necessary to compute the secondary hash code.

\Y\P\?\4\X10:Types and objects visible from \\{data\_structures}\X\mathrel{+}\S%
\?\6
\\{id\_first}\?:\?\\{buffer\_index};\C{where the current identifier begins in
the buffer}\6
\\{id\_loc}\?:\?\\{buffer\_index};\C{just after the current identifier in the
buffer}\7
\&{subtype}\8\\{hash\_index}\8\&{is}\8\p{INTEGER}\8\&{range}\80\?\to\?\\{hash%
\_size};\6
\&{subtype}\8\\{chop\_hash\_index}\8\&{is}\8\p{INTEGER}\8\&{range}\80\?\to\?%
\\{unambig\_length};\6
\\{hash}\?,\39\?\\{chop\_hash}\?:\1\&{array}\8\1(\?\\{hash\_index}\?\2)\2\8%
\&{of}\?\8\\{sixteen\_bits}\?\K\?\\{null\_array};\C{heads of hash
lists---initially empty}\6
\\{chopped\_id}\?:\1\&{array}\8\1(\?\\{chop\_hash\_index}\?\2)\2\8\&{of}\?\8%
\\{ASCII\_code};\C{chopped identifier}\par
\fi

\M50. Here now is the main procedure for finding identifiers (and strings).
The parameter \|t is set to \\{normal} except when the identifier is
a macro name that is just being defined; in the latter case, \|t will be
\\{simple} or \\{parametric}.

\Y\P\?\4\X50:Specification of procedures visible from \\{hashing}\X\S\?\6
\&{function}\8\1\\{id\_lookup}\?(\1\?\|t\?:\?\\{eight\_bits}\?\2)\&{return}\?\8%
\\{name\_pointer}\2;\par
\A sections~64 and~68.
\U section~5\*.\fi

\M51. \P\?\X51:Procedures in \\{hashing}\X\S\?\6
\&{function}\8\1\\{id\_lookup}\?(\1\?\|t\?:\?\\{eight\_bits}\?\2)\&{return}\?\8%
\\{name\_pointer}\2\8\&{is}\8\C{finds current            identifier}\1\6
\|c\?:\?\\{eight\_bits};\C{byte being chopped}\6
\|i\?:\?\\{buffer\_index};\C{index into \\{buffer}}\6
\|h\?:\?\\{hash\_index};\C{hash code}\6
\|k\?:\?\\{byte\_index};\C{index into \\{byte\_mem}}\6
\|w\?:\?\\{w\_range};\C{segment of \\{byte\_mem}}\6
\|l\?:\?\\{buffer\_index}\?\K\?\\{id\_loc}\?-\?\\{id\_first};\C{length of the
given identifier}\6
\|p\?,\39\?\|q\?:\?\\{name\_pointer};\C{where the identifier is being sought}\6
\|s\?:\?\\{chop\_hash\_index};\C{index into \\{chopped\_id}}\6
\4\&{begin}\6
\X52:Compute the hash code \|h\X;\6
\X53:Compute the name location \|p\X;\6
\&{if}\1\1\?\8\1(\?\|p\?=\?\\{name\_ptr}\?\2)\V\1(\?\|t\?\I\?\\{normal}\?\2)\?%
\2\8\&{then}\6
\X55:Update the tables and check for possible errors\X;\2\6
\&{end}\8\2\&{if}\1;\6
\?\&{return}\?\8\|p;\2\6
\&{end}\8\\{id\_lookup};\par
\A sections~65 and~69.
\U section~5\*.\fi

\M52. A simple hash code is used: If the sequence of
ASCII codes is $c_1c_2\ldots c_n$, its hash value will be
$$(2^{n-1}c_1+2^{n-2}c_2+\cdots+c_n)\,\bmod\,\\{hash\_size}.$$

\Y\P\?\4\X52:Compute the hash code \|h\X\S\?\6
\|h\?\K\?\\{buffer}\?\1(\?\\{id\_first}\?\2)\?;\5
\|i\?\K\?\\{id\_first}\?+\?1;\6
\1\&{while}\8\|i\?<\?\\{id\_loc}\8\&{loop}\6
\|h\?\K\1(\?\|h\?+\?\|h\?+\?\\{buffer}\?\1(\?\|i\?\2)\2)\mathbin{\&{mod}}\?%
\\{hash\_size};\5
\\{incr}\?\1(\?\|i\?\2)\?;\2\6
\&{end}\8\&{loop}\par
\U section~51.\fi

\M53. If the identifier is new, it will be placed in position \|p\?=\?\\{name%
\_ptr},
otherwise \|p will point to its existing location.

\Y\P\?\4\X53:Compute the name location \|p\X\S\?\6
\|p\?\K\?\\{hash}\?\1(\?\|h\?\2)\?;\6
\1\&{while}\8\|p\?\I\?0\8\&{loop}\6
\&{if}\1\1\8\\{length}\?\1(\?\|p\?\2)=\?\|l\2\8\&{then}\6
\X54:Compare name \|p with current identifier, \&{goto}\8\\{found} if equal\X;%
\2\6
\&{end}\8\2\&{if}\1;\6
\|p\?\K\?\\{link}\?\1(\?\|p\?\2)\?;\2\6
\&{end}\8\&{loop};\6
\|p\?\K\?\\{name\_ptr};\C{the current identifier is new}\6
\\{link}\?\1(\?\|p\?\2)\K\?\\{hash}\?\1(\?\|h\?\2)\?;\5
\\{hash}\?\1(\?\|h\?\2)\K\?\|p;\C{insert \|p at beginning of hash list}\6
\4\BL\\{found}\EL\8\&{null}\par
\U section~51.\fi

\M54. \P\?\X54:Compare name \|p with current identifier, \&{goto}\8\\{found} if
equal\X\S\?\6
\|i\?\K\?\\{id\_first};\5
\|k\?\K\?\\{byte\_start}\?\1(\?\|p\?\2)\?;\5
\|w\?\K\?\|p\?\mathbin{\&{mod}}\?\\{ww};\6
\1\&{while}\?\8\1(\?\|i\?<\?\\{id\_loc}\?\2)\W\1(\?\\{buffer}\?\1(\?\|i\?\2)=\?%
\\{byte\_mem}\?\1(\?\|w\?,\39\?\|k\?\2)\2)\?\8\&{loop}\6
\\{incr}\?\1(\?\|i\?\2)\?;\5
\\{incr}\?\1(\?\|k\?\2)\?;\2\6
\&{end}\8\&{loop};\6
\&{if}\1\1\8\|i\?=\?\\{id\_loc}\2\8\&{then}\6
\&{goto}\8\\{found};\C{all characters agree}\2\6
\&{end}\8\2\&{if}\1\par
\U section~53.\fi

\M55. \P\?\X55:Update the tables and check for possible errors\X\S\?\6
\&{if}\1\1\?\8\1(\1(\?\|p\?\I\?\\{name\_ptr}\?\2)\W\1(\?\|t\?\I\?\\{normal}\?%
\2)\W\1(\?\\{ilk}\?\1(\?\|p\?\2)=\?\\{normal}\?\2)\2)\V\1(\1(\?\|p\?=\?\\{name%
\_ptr}\?\2)\W\1(\?\|t\?=\?\\{normal}\?\2)\2)\?\2\8\&{then}\6
\X56:Compute the secondary hash code \|h and put the first characters into the
auxiliary array \\{chopped\_id}\X;\2\6
\&{end}\8\2\&{if}\1;\6
\&{if}\1\1\8\|p\?\I\?\\{name\_ptr}\2\8\&{then}\6
\X57:Give double-definition error and change \|p to type \|t\X;\6
\4\&{else}\6
\X59:Enter a new identifier into the table at position \|p\X;\2\6
\&{end}\8\2\&{if}\1\par
\U section~51.\fi

\M56. The following routine, which is called into play when it is necessary to
look at the secondary hash table, computes the same hash function as before
(but on the chopped data), and places a zero after the chopped identifier
in \\{chopped\_id} to serve as a convenient sentinel.

\Y\P\?\4\X56:Compute the secondary hash code \|h and put the first characters
into the auxiliary array \\{chopped\_id}\X\S\?\6
\|i\?\K\?\\{id\_first};\5
\|s\?\K\?0;\5
\|h\?\K\?0;\6
\1\&{while}\?\8\1(\?\|i\?<\?\\{id\_loc}\?\2)\W\1(\?\|s\?<\?\\{unambig\_length}%
\?\2)\?\8\&{loop}\6
\&{if}\1\1\8\\{buffer}\?\1(\?\|i\?\2)\G\?\H{a}\2\8\&{then}\6
\\{chopped\_id}\?\1(\?\|s\?\2)\K\?\\{buffer}\?\1(\?\|i\?\2)-\bn{8}{40}\?;\6
\4\&{else}\6
\\{chopped\_id}\?\1(\?\|s\?\2)\K\?\\{buffer}\?\1(\?\|i\?\2)\?;\2\6
\&{end}\8\2\&{if}\1;\6
\|h\?\K\1(\?\|h\?+\?\|h\?+\?\\{chopped\_id}\?\1(\?\|s\?\2)\2)\mathbin{\&{mod}}%
\?\\{hash\_size};\5
\\{incr}\?\1(\?\|s\?\2)\?;\5
\\{incr}\?\1(\?\|i\?\2)\?;\2\6
\&{end}\8\&{loop};\6
\\{chopped\_id}\?\1(\?\|s\?\2)\K\?0\par
\U section~55.\fi

\M57. If a macro has appeared before it was defined, \.{ATANGLE} will still
work
all right; after all, such behavior is typical of the replacement texts for
modules, which act very much like macros. However, \.{ATANGLE} finds it
simplest to make a blanket rule that macros should be defined before they are
used. The following routine gives an error message and also fixes up any
damage that may have been caused.

\Y\P\?\4\X57:Give double-definition error and change \|p to type \|t\X\S\?\6
\C{now \|p\?\I\?\\{name\_ptr} and \|t\?\I\?\\{normal}}\6
\&{if}\1\1\8\\{ilk}\?\1(\?\|p\?\2)=\?\\{normal}\2\8\&{then}\6
\\{err\_print}\?\1(\?\.{"!\ This\ identifier\ has\ already\ appeared"}\?\2)\?;\6
\X58:Remove \|p from secondary hash table\X;\6
\4\&{else}\6
\\{err\_print}\?\1(\?\.{"!\ This\ identifier\ was\ defined\ before"}\?\2)\?;\2\6
\&{end}\8\2\&{if}\1;\6
\\{ilk}\?\1(\?\|p\?\2)\K\?\|t\par
\U section~55.\fi

\M58. When we have to remove a secondary hash entry, because a \\{normal}
identifier is changing to another \\{ilk}, the hash code \|h and chopped
identifier have already been computed.

\Y\P\?\4\X58:Remove \|p from secondary hash table\X\S\?\6
\|q\?\K\?\\{chop\_hash}\?\1(\?\|h\?\2)\?;\6
\&{if}\1\1\8\|q\?=\?\|p\2\8\&{then}\6
\\{chop\_hash}\?\1(\?\|h\?\2)\K\?\\{equiv}\?\1(\?\|p\?\2)\?;\6
\4\&{else}\6
\1\&{while}\8\\{equiv}\?\1(\?\|q\?\2)\I\?\|p\8\&{loop}\6
\|q\?\K\?\\{equiv}\?\1(\?\|q\?\2)\?;\2\6
\&{end}\8\&{loop};\6
\\{equiv}\?\1(\?\|q\?\2)\K\?\\{equiv}\?\1(\?\|p\?\2)\?;\2\6
\&{end}\8\2\&{if}\1\par
\U section~57.\fi

\M59. \P\?\X59:Enter a new identifier into the table at position \|p\X\S\?\6
\&{if}\1\1\8\|t\?=\?\\{normal}\2\8\&{then}\6
\X60:Check for ambiguity and update secondary hash\X;\2\6
\&{end}\8\2\&{if}\1;\6
\|w\?\K\?\\{name\_ptr}\?\mathbin{\&{mod}}\?\\{ww};\5
\|k\?\K\?\\{byte\_ptr}\?\1(\?\|w\?\2)\?;\6
\&{if}\1\1\8\|k\?+\?\|l\?>\?\\{max\_bytes}\2\8\&{then}\6
\\{overflow}\?\1(\?\.{"byte\ memory"}\?\2)\?;\2\6
\&{end}\8\2\&{if}\1;\6
\&{if}\1\1\8\\{name\_ptr}\?>\?\\{max\_names}\?-\?\\{ww}\2\8\&{then}\6
\\{overflow}\?\1(\?\.{"name"}\?\2)\?;\2\6
\&{end}\8\2\&{if}\1;\6
\|i\?\K\?\\{id\_first};\C{get ready to move the identifier into \\{byte\_mem}}\6
\1\&{while}\8\|i\?<\?\\{id\_loc}\8\&{loop}\6
\\{byte\_mem}\?\1(\?\|w\?,\39\?\|k\?\2)\K\?\\{buffer}\?\1(\?\|i\?\2)\?;\5
\\{incr}\?\1(\?\|k\?\2)\?;\5
\\{incr}\?\1(\?\|i\?\2)\?;\2\6
\&{end}\8\&{loop};\6
\\{byte\_ptr}\?\1(\?\|w\?\2)\K\?\|k;\5
\\{byte\_start}\?\1(\?\\{name\_ptr}\?+\?\\{ww}\?\2)\K\?\|k;\5
\\{incr}\?\1(\?\\{name\_ptr}\?\2)\?;\5
\\{ilk}\?\1(\?\|p\?\2)\K\?\|t\par
\U section~55.\fi

\M60. \P\?\X60:Check for ambiguity and update secondary hash\X\S\?\6
\|q\?\K\?\\{chop\_hash}\?\1(\?\|h\?\2)\?;\6
\1\&{while}\8\|q\?\I\?0\8\&{loop}\6
\X61:Check if \|q conflicts with \|p\X;\6
\|q\?\K\?\\{equiv}\?\1(\?\|q\?\2)\?;\2\6
\&{end}\8\&{loop};\6
\\{equiv}\?\1(\?\|p\?\2)\K\?\\{chop\_hash}\?\1(\?\|h\?\2)\?;\5
\\{chop\_hash}\?\1(\?\|h\?\2)\K\?\|p\C{put \|p at front of secondary list}\par
\U section~59.\fi

\M61. \P\?\X61:Check if \|q conflicts with \|p\X\S\?\6
\|k\?\K\?\\{byte\_start}\?\1(\?\|q\?\2)\?;\5
\|s\?\K\?0;\5
\|w\?\K\?\|q\?\mathbin{\&{mod}}\?\\{ww};\6
\1\&{while}\?\8\1(\?\|k\?<\?\\{byte\_start}\?\1(\?\|q\?+\?\\{ww}\?\2)\2)\W\1(\?%
\|s\?<\?\\{unambig\_length}\?\2)\?\8\&{loop}\6
\|c\?\K\?\\{byte\_mem}\?\1(\?\|w\?,\39\?\|k\?\2)\?;\6
\&{if}\1\1\8\|c\?\G\?\H{a}\2\8\&{then}\6
\|c\?\K\?\|c\?-\bn{8}{40}\?;\C{convert to uppercase}\2\6
\&{end}\8\2\&{if}\1;\6
\&{if}\1\1\8\\{chopped\_id}\?\1(\?\|s\?\2)\I\?\|c\2\8\&{then}\6
\&{goto}\8\\{not\_found};\2\6
\&{end}\8\2\&{if}\1;\6
\\{incr}\?\1(\?\|s\?\2)\?;\5
\\{incr}\?\1(\?\|k\?\2)\?;\2\6
\&{end}\8\&{loop};\6
\&{if}\1\1\?\8\1(\?\|k\?=\?\\{byte\_start}\?\1(\?\|q\?+\?\\{ww}\?\2)\2)\W\1(\?%
\\{chopped\_id}\?\1(\?\|s\?\2)\I\?0\?\2)\?\2\8\&{then}\6
\&{goto}\8\\{not\_found};\2\6
\&{end}\8\2\&{if}\1;\6
\\{print\_nl}\?\1(\?\.{"!\ Identifier\ conflict\ with\ "}\?\2)\?;\6
\1\&{for}\8\\{kl}\?\in\?\\{byte\_start}\?\1(\?\|q\?\2)\to\?\\{byte\_start}\?\1(%
\?\|q\?+\?\\{ww}\?\2)-\?1\8\&{loop}\6
\\{print}\?\1(\?\\{xchr}\?\1(\?\\{byte\_mem}\?\1(\?\|w\?,\39\?\\{kl}\?\2)\2)\2)%
\?;\2\6
\&{end}\8\&{loop};\6
\\{error};\5
\|q\?\K\?0;\C{only one conflict will be printed, since \\{equiv}\?(\?0\?)=\?0}\6
\4\BL\\{not\_found}\EL\8\&{null}\par
\U section~60.\fi

\N62.  Searching for module names.
The \\{mod\_lookup} procedure finds the module name \\{mod\_text}\?(\?1\?\to\?%
\|l\?)\? in the
search tree, after inserting it if necessary, and returns a pointer to
where it was found.

\Y\P\?\4\X10:Types and objects visible from \\{data\_structures}\X\mathrel{+}\S%
\?\6
\&{subtype}\8\\{long\_name\_index}\8\&{is}\8\p{INTEGER}\8\&{range}\80\?\to\?%
\\{longest\_name};\6
\\{mod\_text}\?:\1\&{array}\8\1(\?\\{long\_name\_index}\?\2)\2\8\&{of}\?\8%
\\{ASCII\_code};\C{name being sought for}\par
\fi

\M63. According to the rules of \.{AWEB}, no module name should be a proper
prefix
of another, so a ``clean'' comparison should occur between any two names. The
result of \\{mod\_lookup} is 0 if this prefix condition is violated. An error
message is printed when such violations are detected during phase two of
\.{AWEAVE}.

\Y\P\?\4\X10:Types and objects visible from \\{data\_structures}\X\mathrel{+}\S%
\?\6
\&{type}\8\\{compare}\8\&{is}\?\8\1(\?\\{less}\?,\39\?\\{equal}\?,\39\?%
\\{greater}\?,\39\?\\{prefix}\?,\39\?\\{extension}\?\2)\?;\par
\fi

\M64. \P\?\X50:Specification of procedures visible from \\{hashing}\X%
\mathrel{+}\S\?\6
\&{function}\8\1\\{mod\_lookup}\?(\1\?\|l\?:\?\\{sixteen\_bits}\?\2)\&{return}%
\?\8\\{name\_pointer}\2;\par
\fi

\M65. \P\?\X51:Procedures in \\{hashing}\X\mathrel{+}\S\?\6
\&{function}\8\1\\{mod\_lookup}\?(\1\?\|l\?:\?\\{sixteen\_bits}\?\2)\&{return}%
\?\8\\{name\_pointer}\2\8\&{is}\8\C{finds module name}\1\6
\|c\?:\?\\{compare}\?\K\?\\{greater};\C{comparison between two names}\6
\|j\?:\?\\{long\_name\_index};\C{index into \\{mod\_text}}\6
\|k\?:\?\\{byte\_index};\C{index into \\{byte\_mem}}\6
\|w\?:\?\\{w\_range};\C{segment of \\{byte\_mem}}\6
\|p\?:\?\\{name\_pointer};\C{current node of the search tree}\6
\|q\?:\?\\{name\_pointer}\?\K\?0;\C{father of node \|p}\6
\4\&{begin}\6
\|p\?\K\?\\{rlink}\?\1(\?0\?\2)\?;\C{\\{rlink}\?(\?0\?)\? is the root of the
tree}\6
\1\&{while}\8\|p\?\I\?0\8\&{loop}\6
\X67:Set \(\|c to the result of comparing the given name to name \|p\X;\6
\|q\?\K\?\|p;\6
\&{if}\1\1\8\|c\?=\?\\{less}\2\8\&{then}\6
\|p\?\K\?\\{llink}\?\1(\?\|q\?\2)\?;\6
\4\&{elsif}\1\8\|c\?=\?\\{greater}\2\8\&{then}\6
\|p\?\K\?\\{rlink}\?\1(\?\|q\?\2)\?;\6
\4\&{else}\6
\&{goto}\8\\{found};\2\6
\&{end}\8\2\&{if}\1;\2\6
\&{end}\8\&{loop};\6
\X66:Enter a new module name into the tree\X;\6
\4\BL\\{found}\EL\8\&{if}\1\1\8\|c\?\I\?\\{equal}\2\8\&{then}\6
\\{err\_print}\?\1(\?\.{"!\ Incompatible\ section\ names"}\?\2)\?;\6
\|p\?\K\?0;\2\6
\&{end}\8\2\&{if}\1;\6
\?\&{return}\?\8\|p;\2\6
\&{end}\8\\{mod\_lookup};\par
\fi

\M66. \P\?\X66:Enter a new module name into the tree\X\S\?\6
\|w\?\K\?\\{name\_ptr}\?\mathbin{\&{mod}}\?\\{ww};\5
\|k\?\K\?\\{byte\_ptr}\?\1(\?\|w\?\2)\?;\6
\&{if}\1\1\8\|k\?+\?\|l\?>\?\\{max\_bytes}\2\8\&{then}\6
\\{overflow}\?\1(\?\.{"byte\ memory"}\?\2)\?;\2\6
\&{end}\8\2\&{if}\1;\6
\&{if}\1\1\8\\{name\_ptr}\?>\?\\{max\_names}\?-\?\\{ww}\2\8\&{then}\6
\\{overflow}\?\1(\?\.{"name"}\?\2)\?;\2\6
\&{end}\8\2\&{if}\1;\6
\|p\?\K\?\\{name\_ptr};\6
\&{if}\1\1\8\|c\?=\?\\{less}\2\8\&{then}\6
\\{llink}\?\1(\?\|q\?\2)\K\?\|p;\6
\4\&{else}\6
\\{rlink}\?\1(\?\|q\?\2)\K\?\|p;\2\6
\&{end}\8\2\&{if}\1;\6
\\{llink}\?\1(\?\|p\?\2)\K\?0;\5
\\{rlink}\?\1(\?\|p\?\2)\K\?0;\5
\|c\?\K\?\\{equal};\5
\\{equiv}\?\1(\?\|p\?\2)\K\?0;\6
\1\&{for}\8\\{kj}\?\in\?1\?\to\?\|l\8\&{loop}\6
\\{byte\_mem}\?\1(\?\|w\?,\39\?\|k\?+\?\\{kj}\?-\?1\?\2)\K\?\\{mod\_text}\?\1(%
\?\\{kj}\?\2)\?;\2\6
\&{end}\8\&{loop};\6
\\{byte\_ptr}\?\1(\?\|w\?\2)\K\?\|k\?+\?\|l;\5
\\{byte\_start}\?\1(\?\\{name\_ptr}\?+\?\\{ww}\?\2)\K\?\|k\?+\?\|l;\5
\\{incr}\?\1(\?\\{name\_ptr}\?\2)\?\par
\U section~65.\fi

\M67. \P\?\X67:Set \(\|c to the result of comparing the given name to name \|p%
\X\S\?\6
\|k\?\K\?\\{byte\_start}\?\1(\?\|p\?\2)\?;\5
\|w\?\K\?\|p\?\mathbin{\&{mod}}\?\\{ww};\5
\|c\?\K\?\\{equal};\5
\|j\?\K\?1;\6
\1\&{while}\?\8\1(\?\|k\?<\?\\{byte\_start}\?\1(\?\|p\?+\?\\{ww}\?\2)\2)\W\1(\?%
\|j\?\L\?\|l\?\2)\W\1(\?\\{mod\_text}\?\1(\?\|j\?\2)=\?\\{byte\_mem}\?\1(\?\|w%
\?,\39\?\|k\?\2)\2)\?\8\&{loop}\6
\\{incr}\?\1(\?\|k\?\2)\?;\5
\\{incr}\?\1(\?\|j\?\2)\?;\2\6
\&{end}\8\&{loop};\6
\&{if}\1\1\8\|k\?=\?\\{byte\_start}\?\1(\?\|p\?+\?\\{ww}\?\2)\?\2\8\&{then}\6
\&{if}\1\1\8\|j\?>\?\|l\2\8\&{then}\6
\|c\?\K\?\\{equal};\6
\4\&{else}\6
\|c\?\K\?\\{extension};\2\6
\&{end}\8\2\&{if}\1;\6
\4\&{elsif}\1\8\|j\?>\?\|l\2\8\&{then}\6
\|c\?\K\?\\{prefix};\6
\4\&{elsif}\1\8\\{mod\_text}\?\1(\?\|j\?\2)<\?\\{byte\_mem}\?\1(\?\|w\?,\39\?%
\|k\?\2)\?\2\8\&{then}\6
\|c\?\K\?\\{less};\6
\4\&{else}\6
\|c\?\K\?\\{greater};\2\6
\&{end}\8\2\&{if}\1\par
\U sections~65 and~69.\fi

\M68. The \\{prefix\_lookup} procedure is supposed to find exactly one module
name that has \\{mod\_text}\?(\?1\?\to\?\|l\?)\? as a prefix. Actually the
algorithm silently
accepts also the situation that some module name is a prefix of
\\{mod\_text}\?(\?1\?\to\?\|l\?)\?, because the user who painstakingly typed in
more than
necessary probably doesn't want to be told about the wasted effort.

\Y\P\?\4\X50:Specification of procedures visible from \\{hashing}\X\mathrel{+}%
\S\?\6
\&{function}\8\1\\{prefix\_lookup}\?(\1\?\|l\?:\?\\{sixteen\_bits}\?\2)%
\&{return}\?\8\\{name\_pointer}\2;\par
\fi

\M69. \P\?\X51:Procedures in \\{hashing}\X\mathrel{+}\S\?\6
\&{function}\8\1\\{prefix\_lookup}\?(\1\?\|l\?:\?\\{sixteen\_bits}\?\2)%
\&{return}\?\8\\{name\_pointer}\2\8\&{is}\8\C{finds name extension}\1\6
\|c\?:\?\\{compare};\C{comparison between two names}\6
\\{count}\?:\?\\{name\_index}\?\K\?0;\C{the number of hits}\6
\|j\?:\?\\{long\_name\_index};\C{index into \\{mod\_text}}\6
\|k\?:\?\\{byte\_index};\C{index into \\{byte\_mem}}\6
\|w\?:\?\\{w\_range};\C{segment of \\{byte\_mem}}\6
\|p\?:\?\\{name\_pointer};\C{current node of the search tree}\6
\|q\?:\?\\{name\_pointer}\?\K\?0;\C{another place to resume the search after
one         branch is done}\6
\|r\?:\?\\{name\_pointer}\?\K\?0;\C{extension found}\6
\4\&{begin}\6
\|p\?\K\?\\{rlink}\?\1(\?0\?\2)\?;\C{begin search at root of tree}\6
\1\&{while}\8\|p\?\I\?0\8\&{loop}\6
\X67:Set \(\|c to the result of comparing the given name to name \|p\X;\6
\&{if}\1\1\8\|c\?=\?\\{less}\2\8\&{then}\6
\|p\?\K\?\\{llink}\?\1(\?\|p\?\2)\?;\6
\4\&{elsif}\1\8\|c\?=\?\\{greater}\2\8\&{then}\6
\|p\?\K\?\\{rlink}\?\1(\?\|p\?\2)\?;\6
\4\&{else}\6
\|r\?\K\?\|p;\5
\\{incr}\?\1(\?\\{count}\?\2)\?;\5
\|q\?\K\?\\{rlink}\?\1(\?\|p\?\2)\?;\5
\|p\?\K\?\\{llink}\?\1(\?\|p\?\2)\?;\2\6
\&{end}\8\2\&{if}\1;\6
\&{if}\1\1\8\|p\?=\?0\2\8\&{then}\6
\|p\?\K\?\|q;\5
\|q\?\K\?0;\2\6
\&{end}\8\2\&{if}\1;\2\6
\&{end}\8\&{loop};\6
\&{if}\1\1\8\\{count}\?\I\?1\2\8\&{then}\6
\&{if}\1\1\8\\{count}\?=\?0\2\8\&{then}\6
\\{err\_print}\?\1(\?\.{"!\ Name\ does\ not\ match"}\?\2)\?;\6
\4\&{else}\6
\\{err\_print}\?\1(\?\.{"!\ Ambiguous\ prefix"}\?\2)\?;\2\6
\&{end}\8\2\&{if}\1;\2\6
\&{end}\8\2\&{if}\1;\6
\?\&{return}\?\8\|r;\C{the result will be 0 if there was no match}\2\6
\&{end}\8\\{prefix\_lookup};\par
\fi

\N70.  Tokens.
Replacement texts, which represent \ADA\ code in a compressed format, appear
in \\{tok\_mem} as mentioned above. The codes in these texts are called
`tokens'; some tokens occupy two consecutive eight-bit byte positions, and
the others take just one byte.

If $p>0$ points to a replacement text, \\{tok\_start}\?(\?\|p\?)\? is the %
\\{tok\_mem}
position of the first eight-bit code of that text. If \\{text\_link}\?(\?\|p%
\?)=\?0,
this is the replacement text for a macro, otherwise it is the replacement
text for a module. In the latter case \\{text\_link}\?(\?\|p\?)\? is either
equal to
\\{module\_flag}, which means that there is no further text for this module, or
\\{text\_link}\?(\?\|p\?)\? points to a continuation of this replacement text;
such links
are created when several modules have \ADA\ texts with the same name, and
they also tie together all the \ADA\ texts of unnamed modules. The
replacement text pointer for the first unnamed module appears in
\\{text\_link}\?(\?0\?)\?, and the most recent such pointer is \\{last%
\_unnamed}.

\Y\P\?\4\X10:Types and objects visible from \\{data\_structures}\X\mathrel{+}\S%
\?\6
\\{module\_flag}\?:\?\&{constant}\8\\{sixteen\_bits}\?\K\?\\{max\_texts};%
\C{final \\{link} in module replacement texts}\6
\\{last\_unnamed}\?:\?\\{text\_pointer}\?\K\?0;\C{most recent replacement text
of unnamed module}\par
\fi

\M71. \P\?\X15:Set initial values\X\mathrel{+}\S\?\6
\\{text\_link}\?\1(\?0\?\2)\K\?0;\par
\fi

\M72. If the first byte of a token is less than \?\bn{8}{200}\?, the token
occupies a
single byte. Otherwise we make a sixteen-bit token by combining two
consecutive bytes \|a and \|b. If \?\bn{8}{200}\L\?\|a\?<\bn{8}{250}\?, then
$(a-\bn{8}{200})\times2^8+b$ points to an identifier; if \?\bn{8}{250}\L\?\|a%
\?<\bn{8}{320}\?,
then $(a-\bn{8}{250})\times2^8+b$ points to a module name; otherwise, i.e.,
if \?\bn{8}{320}\L\?\|a\?<\bn{8}{400}\?, then $(a-\bn{8}{320})\times2^8+b$ is
the number of the
module in which the current replacement text appears.

Codes less than \bn{8}{200} are 7-bit ASCII codes that represent themselves.
In particular, a single-character identifier like `\|x' will be a one-byte
token, while all longer identifiers will occupy two bytes.

Some of the 7-bit ASCII codes will not be present, however, so we can
use them for special purposes. The following symbolic names are used:

\yskip\hang \\{param} denotes insertion of a parameter. This occurs only in
the replacement texts of parametric macros, outside of single-quoted strings
in those texts.

\hang \\{begin\_comment} denotes \.{@\{}, which will become \.{--}.

\hang \\{end\_comment} denotes \.{@\}}, which will become a new line.

\hang \\{character\_tok} denotes the \.' that encloses an character constant.

\hang \\{ascii\_tok} denotes the \.{@'} that precedes a ASCII code constant.

\hang \\{output\_tok} denotes the \.{@\~} that precedes an output file-name.

\hang \\{join} denotes the concatenation of adjacent items with no
space or line breaks allowed between them (the \.{@\&} operation of \.{AWEB}).

\hang \\{double\_dot} denotes `\.{..}' in \ADA.

\hang \\{verbatim} denotes the \.{@=} that begins a verbatim \ADA\ string.
It is also used for the end of the string.

\hang \\{force\_line} denotes the \.{@\\} that forces a new line in the
\ADA\ output.

\Y\P\?\4\X10:Types and objects visible from \\{data\_structures}\X\mathrel{+}\S%
\?\6
\\{param}\?:\?\&{constant}\8\\{eight\_bits}\?\K\?0;\C{ASCII null code will not
appear}\6
\\{verbatim}\?:\?\&{constant}\8\\{eight\_bits}\?\K\?2;\C{extended ASCII alpha
should not appear}\6
\\{force\_line}\?:\?\&{constant}\8\\{eight\_bits}\?\K\?3;\C{extended ASCII beta
should not appear}\6
\\{begin\_comment}\?:\?\&{constant}\8\\{eight\_bits}\?\K\bn{8}{11}\?;\C{ASCII
tab mark will not appear}\6
\\{end\_comment}\?:\?\&{constant}\8\\{eight\_bits}\?\K\bn{8}{12}\?;\C{ASCII
line feed will not appear}\6
\\{character\_tok}\?:\?\&{constant}\8\\{eight\_bits}\?\K\bn{8}{14}\?;\C{ASCII
form feed will not appear}\6
\\{ascii\_tok}\?:\?\&{constant}\8\\{eight\_bits}\?\K\bn{8}{15}\?;\C{ASCII
carriage return will not appear}\6
\\{output\_tok}\?:\?\&{constant}\8\\{eight\_bits}\?\K\bn{8}{16}\?;\C{ASCII %
\.{DLE} will not appear}\6
\\{double\_dot}\?:\?\&{constant}\8\\{eight\_bits}\?\K\bn{8}{40}\?;\C{ASCII
space will not appear except in strings}\6
\\{join}\?:\?\&{constant}\8\\{eight\_bits}\?\K\bn{8}{177}\?;\C{ASCII delete
will not appear}\par
\fi

\M73. \P\?\X47:Specification of procedures visible from \\{data\_structures}\X%
\mathrel{+}\S\?\6
\&{procedure}\8\1\\{store\_two\_bytes}\?(\1\?\|x\?:\?\\{sixteen\_bits}\?\2)\?%
\2;\par
\fi

\M74. The following procedure is used to enter a two-byte value into \\{tok%
\_mem}
when a replacement text is being generated.

\Y\P\?\4\X48:Procedures in \\{data\_structures}\X\mathrel{+}\S\?\6
\&{procedure}\8\1\\{store\_two\_bytes}\?(\1\?\|x\?:\?\\{sixteen\_bits}\?\2)\?\2%
\8\&{is}\8\C{stores high byte, then low byte}\1\6
\4\&{begin}\6
\&{if}\1\1\8\\{tok\_ptr}\?\1(\?\|z\?\2)+\?2\?>\?\\{max\_toks}\2\8\&{then}\6
\\{overflow}\?\1(\?\.{"token"}\?\2)\?;\2\6
\&{end}\8\2\&{if}\1;\6
\\{tok\_mem}\?\1(\?\|z\?,\39\?\\{tok\_ptr}\?\1(\?\|z\?\2)\2)\K\?\|x\?/%
\bn{8}{400}\?;\C{this could be done by a shift command}\6
\\{tok\_mem}\?\1(\?\|z\?,\39\?\\{tok\_ptr}\?\1(\?\|z\?\2)+\?1\?\2)\K\?\|x\?%
\mathbin{\&{mod}}\bn{8}{400}\?;\C{this could be done by a logical and}\6
\\{tok\_ptr}\?\1(\?\|z\?\2)\K\?\\{tok\_ptr}\?\1(\?\|z\?\2)+\?2;\2\6
\&{end}\8\\{store\_two\_bytes};\par
\fi

\N75.  Stacks for output.
Let's make sure that our data structures contain enough information to
produce the entire \ADA\ program as desired, by working next on the
algorithms that actually do produce that program.

\fi

\M76. The output process uses a stack to keep track of what is going on at
different ``levels'' as the macros are being expanded.
Entries on this stack have five parts:

\yskip\hang \\{end\_field} is the \\{tok\_mem} location where the replacement
text of a particular level will end;

\hang \\{byte\_field} is the \\{tok\_mem} location from which the next token
on a particular level will be read;

\hang \\{name\_field} points to the name corresponding to a particular level;

\hang \\{repl\_field} points to the replacement text currently being read
at a particular level.

\hang \\{mod\_field} is the module number, or zero if this is a macro.

\yskip\noindent The current values of these five quantities are referred to
quite frequently, so they are stored in a separate place instead of in
the \\{stack} array. We call the current values \\{cur\_end}, \\{cur\_byte},
\\{cur\_name}, \\{cur\_repl}, and \\{cur\_mod}.

The global variable \\{stack\_ptr} tells how many levels of output are
currently in progress. The end of all output occurs when the stack is
empty, i.e., when \\{stack\_ptr}\?=\?0.

\Y\P\?\4\X10:Types and objects visible from \\{data\_structures}\X\mathrel{+}\S%
\?\6
\&{subtype}\8\\{mod\_range}\8\&{is}\8\p{INTEGER}\8\&{range}\80\?\to%
\bn{8}{27777}\?;\6
\&{type}\8\\{output\_state}\8\&{is}\8\1\6
\&{record}\1\6
\\{end\_field}\?:\?\\{sixteen\_bits};\C{ending location of replacement text}\6
\\{byte\_field}\?:\?\\{sixteen\_bits};\C{present location within replacement
text}\6
\\{name\_field}\?:\?\\{name\_pointer};\C{\\{byte\_start} index for text being
output}\6
\\{repl\_field}\?:\?\\{text\_pointer};\C{\\{tok\_start} index for text being
output}\6
\\{mod\_field}\?:\?\\{mod\_range};\C{module number or zero if not a module}\2\6
\2\&{end}\8\&{record};\7
\\{cur\_state}\?:\?\\{output\_state};\C{\\{cur\_end}, \\{cur\_byte}, \\{cur%
\_name}, \\{cur\_repl}}\6
\\{stack}\?:\1\&{array}\8\1(\?1\?\to\?\\{stack\_size}\?\2)\2\8\&{of}\?\8%
\\{output\_state};\C{info for non-current levels}\6
\\{stack\_ptr}\?:\?\p{INTEGER}\8\&{range}\80\?\to\?\\{stack\_size};\C{first
unused location in the output state stack}\7
\\{cur\_end}\?:\?\\{sixteen\_bits}\8\&{renames}\8\\{cur\_state}\?.\?\\{end%
\_field};\C{current ending location in \\{tok\_mem}}\6
\\{cur\_byte}\?:\?\\{sixteen\_bits}\8\&{renames}\8\\{cur\_state}\?.\?\\{byte%
\_field};\C{location of next output byte in \\{tok\_mem}}\6
\\{cur\_name}\?:\?\\{name\_pointer}\8\&{renames}\8\\{cur\_state}\?.\?\\{name%
\_field};\C{pointer to current name being expanded}\6
\\{cur\_repl}\?:\?\\{text\_pointer}\8\&{renames}\8\\{cur\_state}\?.\?\\{repl%
\_field};\C{pointer to current replacement text}\6
\\{cur\_mod}\?:\?\\{mod\_range}\8\&{renames}\8\\{cur\_state}\?.\?\\{mod%
\_field};\C{current module number being expanded}\par
\fi

\M77. It is convenient to keep a global variable \\{zo} equal to \\{cur\_repl}%
\?\mathbin{\&{mod}}\?\\{zz}.

\Y\P\?\4\X10:Types and objects visible from \\{data\_structures}\X\mathrel{+}\S%
\?\6
\\{zo}\?:\?\\{z\_range};\C{the segment of \\{tok\_mem} from which output
                      is coming}\par
\fi

\M78. Parameters must also be stacked. They are placed in \\{tok\_mem} just
above
the other replacement texts, and dummy parameter `names' are placed in
\\{byte\_start} just after the other names. The variables \\{text\_ptr} and
\\{tok\_ptr}\?(\?\|z\?)\? essentially serve as parameter stack pointers during
the output
phase, so there is no need for a separate data structure to handle this
problem.

\fi

\M79. There is an implicit stack corresponding to meta-comments that are output
via \.{@\{} and \.{@\}}. But this stack need not be represented in detail,
because we only need to know whether it is empty or not. A global variable
\\{brace\_level} tells how many items would be on this stack if it were
present.

\Y\P\?\4\X10:Types and objects visible from \\{data\_structures}\X\mathrel{+}\S%
\?\6
\\{brace\_level}\?:\?\\{eight\_bits};\C{current depth of $\.{@\{}\ldots\.{@\}}$
nesting}\par
\fi

\M80. To get the output process started, we will perform the following
initialization steps. We may assume that \\{text\_link}\?(\?0\?)\? is nonzero,
since it
points to the \ADA\ text in the first unnamed module that generates code; if
there are no such modules, there is nothing to output, and an error message
will have been generated before we do any of the initialization.

\Y\P\?\4\X80:Initialize the output stacks\X\S\?\6
\\{stack\_ptr}\?\K\?1;\5
\\{brace\_level}\?\K\?0;\5
\\{cur\_name}\?\K\?0;\5
\\{cur\_repl}\?\K\?\\{text\_link}\?\1(\?0\?\2)\?;\5
\\{zo}\?\K\?\\{cur\_repl}\?\mathbin{\&{mod}}\?\\{zz};\5
\\{cur\_byte}\?\K\?\\{tok\_start}\?\1(\?\\{cur\_repl}\?\2)\?;\5
\\{cur\_end}\?\K\?\\{tok\_start}\?\1(\?\\{cur\_repl}\?+\?\\{zz}\?\2)\?;\5
\\{cur\_mod}\?\K\?0;\par
\U section~100.\fi

\M81. When the replacement text for name \|p is to be inserted into the output,
the following subroutine is called to save the old level of output and get
the new one going.

\Y\P\?\4\X81:Procedures in \\{output\_phase}\X\S\?\6
\&{procedure}\8\1\\{push\_level}\?(\1\?\|p\?:\?\\{name\_pointer}\?\2)\?\2\8%
\&{is}\8\C{suspends the current level}\1\6
\4\&{begin}\6
\&{if}\1\1\8\\{stack\_ptr}\?=\?\\{stack\_size}\2\8\&{then}\6
\\{overflow}\?\1(\?\.{"stack"}\?\2)\?;\6
\4\&{else}\6
\\{stack}\?\1(\?\\{stack\_ptr}\?\2)\K\?\\{cur\_state};\C{save \\{cur\_end}, %
\\{cur\_byte}, etc.}\6
\\{incr}\?\1(\?\\{stack\_ptr}\?\2)\?;\5
\\{cur\_name}\?\K\?\|p;\5
\\{cur\_repl}\?\K\?\\{equiv}\?\1(\?\|p\?\2)\?;\5
\\{zo}\?\K\?\\{cur\_repl}\?\mathbin{\&{mod}}\?\\{zz};\5
\\{cur\_byte}\?\K\?\\{tok\_start}\?\1(\?\\{cur\_repl}\?\2)\?;\5
\\{cur\_end}\?\K\?\\{tok\_start}\?\1(\?\\{cur\_repl}\?+\?\\{zz}\?\2)\?;\5
\\{cur\_mod}\?\K\?0;\2\6
\&{end}\8\2\&{if}\1;\2\6
\&{end}\8\\{push\_level};\par
\A sections~82, 84, 94, 96, 98, 101, and~114.
\U section~6\*.\fi

\M82. When we come to the end of a replacement text, the \\{pop\_level}
subroutine
does the right thing: It either moves to the continuation of this replacement
text or returns the state to the most recently stacked level. Part of this
subroutine, which updates the parameter stack, will be given later when we
study the parameter stack in more detail.

\Y\P\?\4\X81:Procedures in \\{output\_phase}\X\mathrel{+}\S\?\6
\&{procedure}\8\1\\{pop\_level}\2\8\&{is}\8\C{do this when \\{cur\_byte}
reaches \\{cur\_end}}\1\6
\4\&{begin}\6
\&{if}\1\1\8\\{text\_link}\?\1(\?\\{cur\_repl}\?\2)=\?0\2\8\&{then}\C{end of
macro expansion}\6
\&{if}\1\1\8\\{ilk}\?\1(\?\\{cur\_name}\?\2)=\?\\{parametric}\2\8\&{then}\6
\X88:Remove a parameter from the parameter stack\X;\2\6
\&{end}\8\2\&{if}\1;\6
\4\&{elsif}\1\8\\{text\_link}\?\1(\?\\{cur\_repl}\?\2)<\?\\{module\_flag}\2\8%
\&{then}\C{link to a continuation}\6
\\{cur\_repl}\?\K\?\\{text\_link}\?\1(\?\\{cur\_repl}\?\2)\?;\C{we will stay on
the same level}\6
\\{zo}\?\K\?\\{cur\_repl}\?\mathbin{\&{mod}}\?\\{zz};\5
\\{cur\_byte}\?\K\?\\{tok\_start}\?\1(\?\\{cur\_repl}\?\2)\?;\5
\\{cur\_end}\?\K\?\\{tok\_start}\?\1(\?\\{cur\_repl}\?+\?\\{zz}\?\2)\?;\6
\&{return};\2\6
\&{end}\8\2\&{if}\1;\6
\\{decr}\?\1(\?\\{stack\_ptr}\?\2)\?;\C{we will go down to the previous level}\6
\&{if}\1\1\8\\{stack\_ptr}\?>\?0\2\8\&{then}\6
\\{cur\_state}\?\K\?\\{stack}\?\1(\?\\{stack\_ptr}\?\2)\?;\5
\\{zo}\?\K\?\\{cur\_repl}\?\mathbin{\&{mod}}\?\\{zz};\2\6
\&{end}\8\2\&{if}\1;\2\6
\&{end}\8\\{pop\_level};\par
\fi

\M83. The heart of the output procedure is the \\{get\_output} routine, which
produces the next token of output that is not a reference to a macro. This
procedure handles all the stacking and unstacking that is necessary. It
returns the value \\{module\_number} if the next output begins or ends the
replacement text of some module, in which case \\{cur\_val} is that module's
number (if beginning) or the negative of that value (if ending). The
procedure also returns the value \\{identifier} if the next output is an
identifier of length two or more, in which case \\{cur\_val} points to that
identifier name.

\Y\P\?\4\X10:Types and objects visible from \\{data\_structures}\X\mathrel{+}\S%
\?\6
\\{module\_number}\?:\?\&{constant}\8\\{sixteen\_bits}\?\K\bn{8}{201}\?;\C{code
returned by \\{get\_output} for module numbers}\6
\\{identifier}\?:\?\&{constant}\8\\{sixteen\_bits}\?\K\bn{8}{202}\?;\C{code
returned by \\{get\_output} for identifiers}\7
\\{cur\_val}\?:\?\p{INTEGER};\C{additional information corresponding to output
token}\par
\fi

\M84. If \\{get\_output} finds that no more output remains, it returns the
value
zero.

\Y\P\?\4\X81:Procedures in \\{output\_phase}\X\mathrel{+}\S\?\6
\&{function}\8\1\\{get\_output}\?\8\&{return}\?\8\\{sixteen\_bits}\2\8\&{is}\8%
\C{returns next token after macro expansion}\1\6
\|a\?:\?\\{sixteen\_bits};\C{value of current byte}\6
\|b\?:\?\\{eight\_bits};\C{byte being copied}\6
\\{bal}\?:\?\\{sixteen\_bits};\C{excess of \.( versus \.) while copying a
parameter}\6
\|k\?:\?\\{byte\_index};\C{index into \\{byte\_mem}}\6
\|w\?:\?\\{w\_range};\C{segment of \\{byte\_mem}}\6
\4\&{begin}\6
\4\BL\\{restart}\EL\8\&{if}\1\1\8\\{stack\_ptr}\?=\?0\2\8\&{then}\6
\?\&{return}\?\80;\2\6
\&{end}\8\2\&{if}\1;\6
\&{if}\1\1\8\\{cur\_byte}\?=\?\\{cur\_end}\2\8\&{then}\6
\\{cur\_val}\?\K-\?\\{cur\_mod};\5
\\{pop\_level};\6
\&{if}\1\1\8\\{cur\_val}\?=\?0\2\8\&{then}\6
\&{goto}\8\\{restart};\2\6
\&{end}\8\2\&{if}\1;\6
\?\&{return}\?\8\\{module\_number};\2\6
\&{end}\8\2\&{if}\1;\6
\|a\?\K\?\\{tok\_mem}\?\1(\?\\{zo}\?,\39\?\\{cur\_byte}\?\2)\?;\5
\\{incr}\?\1(\?\\{cur\_byte}\?\2)\?;\6
\&{if}\1\1\8\|a\?<\bn{8}{200}\?\2\8\&{then}\C{one-byte token}\6
\&{if}\1\1\8\|a\?=\?\\{param}\2\8\&{then}\6
\X89:Start scanning current macro parameter, \&{goto}\8\\{restart}\X;\6
\4\&{else}\6
\?\&{return}\?\8\|a;\2\6
\&{end}\8\2\&{if}\1;\2\6
\&{end}\8\2\&{if}\1;\6
\|a\?\K\1(\?\|a\?-\bn{8}{200}\2)\ast\bn{8}{400}+\?\\{tok\_mem}\?\1(\?\\{zo}\?,%
\39\?\\{cur\_byte}\?\2)\?;\5
\\{incr}\?\1(\?\\{cur\_byte}\?\2)\?;\6
\&{if}\1\1\8\|a\?<\bn{8}{24000}\?\2\8\&{then}\C{\?\bn{8}{24000}=(\bn{8}{250}-%
\bn{8}{200})\ast\bn{8}{400}\?}\6
\X86:Expand macro \|a and \&{return}, or \&{goto}\8\\{restart} if no output
found\X;\2\6
\&{end}\8\2\&{if}\1;\6
\&{if}\1\1\8\|a\?<\bn{8}{50000}\?\2\8\&{then}\C{\?\bn{8}{50000}=(\bn{8}{320}-%
\bn{8}{200})\ast\bn{8}{400}\?}\6
\X85:Expand module \|a\?-\bn{8}{24000}\?, \&{goto}\8\\{restart}\X;\2\6
\&{end}\8\2\&{if}\1;\6
\\{cur\_val}\?\K\?\|a\?-\bn{8}{50000}\?;\5
\\{cur\_mod}\?\K\?\\{cur\_val};\5
\?\&{return}\?\8\\{module\_number};\2\6
\&{end}\8\\{get\_output};\par
\fi

\M85. The user may have forgotten to give any \ADA\ text for a module name,
or the \ADA\ text may have been associated with a different name by mistake.

\Y\P\?\4\X85:Expand module \|a\?-\bn{8}{24000}\?, \&{goto}\8\\{restart}\X\S\?\6
\|a\?\K\?\|a\?-\bn{8}{24000}\?;\6
\&{if}\1\1\8\\{equiv}\?\1(\?\|a\?\2)\I\?0\2\8\&{then}\6
\\{push\_level}\?\1(\?\|a\?\2)\?;\6
\4\&{elsif}\1\8\|a\?\I\?0\2\8\&{then}\6
\\{print\_nl}\?\1(\?\.{"!\ Not\ present:\ <"}\?\2)\?;\5
\\{print\_id}\?\1(\?\|a\?\2)\?;\5
\\{print}\?\1(\?\.{">"}\?\2)\?;\5
\\{error};\2\6
\&{end}\8\2\&{if}\1;\6
\&{goto}\8\\{restart}\par
\U section~84.\fi

\M86. \P\?\X86:Expand macro \|a and \&{return}, or \&{goto}\8\\{restart} if no
output found\X\S\?\6
\&{case}\1\8\\{ilk}\?\1(\?\|a\?\2)\?\1\8\&{is}\8\6
\4\&{when}\8\\{normal}\KR\6
\\{cur\_val}\?\K\?\|a;\5
\?\&{return}\?\8\\{identifier};\6
\4\&{when}\8\\{simple}\KR\6
\\{push\_level}\?\1(\?\|a\?\2)\?;\5
\&{goto}\8\\{restart};\6
\4\&{when}\8\\{parametric}\KR\6
\X87:Put a parameter on the parameter stack, or \&{goto}\8\\{restart} if error
occurs\X;\6
\\{push\_level}\?\1(\?\|a\?\2)\?;\5
\&{goto}\8\\{restart};\6
\4\&{when}\8\&{others}\8\KR\6
\\{confusion}\?\1(\?\.{"output"}\?\2)\?;\2\6
\4\2\&{end}\8\&{case};\6
\?\&{return}\?\8\|a\par
\U section~84.\fi

\M87. We come now to the interesting part, the job of putting a parameter on
the parameter stack. First we pop the stack if necessary until getting to
a level that hasn't ended. Then the next character must be a `\.(';
and since parentheses are balanced on each level, the entire parameter must
be present, so we can copy it without difficulty.

\Y\P\?\4\X87:Put a parameter on the parameter stack, or \&{goto}\8\\{restart}
if error occurs\X\S\?\6
\1\&{while}\?\8\1(\?\\{cur\_byte}\?=\?\\{cur\_end}\?\2)\W\1(\?\\{stack\_ptr}\?>%
\?0\?\2)\?\8\&{loop}\6
\\{pop\_level};\2\6
\&{end}\8\&{loop};\6
\&{if}\1\1\?\8\1(\?\\{stack\_ptr}\?=\?0\?\2)\V\1(\?\\{tok\_mem}\?\1(\?\\{zo}\?,%
\39\?\\{cur\_byte}\?\2)\I\?\H{(}\?\2)\?\2\8\&{then}\6
\\{print\_nl}\?\1(\?\.{"!\ No\ parameter\ given\ for\ "}\?\2)\?;\5
\\{print\_id}\?\1(\?\|a\?\2)\?;\5
\\{error};\5
\&{goto}\8\\{restart};\2\6
\&{end}\8\2\&{if}\1;\6
\X90:Copy the parameter into \\{tok\_mem}\X;\6
\\{equiv}\?\1(\?\\{name\_ptr}\?\2)\K\?\\{text\_ptr};\5
\\{ilk}\?\1(\?\\{name\_ptr}\?\2)\K\?\\{simple};\5
\|w\?\K\?\\{name\_ptr}\?\mathbin{\&{mod}}\?\\{ww};\5
\|k\?\K\?\\{byte\_ptr}\?\1(\?\|w\?\2)\?;\6
\&{if}\1\1\8\\{name\_ptr}\?>\?\\{max\_names}\?-\?\\{ww}\2\8\&{then}\6
\\{overflow}\?\1(\?\.{"name"}\?\2)\?;\2\6
\&{end}\8\2\&{if}\1;\6
\\{byte\_start}\?\1(\?\\{name\_ptr}\?+\?\\{ww}\?\2)\K\?\|k;\5
\\{incr}\?\1(\?\\{name\_ptr}\?\2)\?;\6
\&{if}\1\1\8\\{text\_ptr}\?>\?\\{max\_texts}\?-\?\\{zz}\2\8\&{then}\6
\\{overflow}\?\1(\?\.{"text"}\?\2)\?;\2\6
\&{end}\8\2\&{if}\1;\6
\\{text\_link}\?\1(\?\\{text\_ptr}\?\2)\K\?0;\5
\\{tok\_start}\?\1(\?\\{text\_ptr}\?+\?\\{zz}\?\2)\K\?\\{tok\_ptr}\?\1(\?\|z\?%
\2)\?;\5
\\{incr}\?\1(\?\\{text\_ptr}\?\2)\?;\5
\|z\?\K\?\\{text\_ptr}\?\mathbin{\&{mod}}\?\\{zz}\par
\U section~86.\fi

\M88. The \\{pop\_level} routine undoes the effect of parameter-pushing when
a parameter macro is finished:

\Y\P\?\4\X88:Remove a parameter from the parameter stack\X\S\?\6
\\{decr}\?\1(\?\\{name\_ptr}\?\2)\?;\5
\\{decr}\?\1(\?\\{text\_ptr}\?\2)\?;\5
\|z\?\K\?\\{text\_ptr}\?\mathbin{\&{mod}}\?\\{zz}; \6
\&{stat}\1\6
\&{if}\1\1\8\\{tok\_ptr}\?\1(\?\|z\?\2)>\?\\{max\_tok\_ptr}\?\1(\?\|z\?\2)\?\2%
\8\&{then}\6
\\{max\_tok\_ptr}\?\1(\?\|z\?\2)\K\?\\{tok\_ptr}\?\1(\?\|z\?\2)\?;\2\6
\&{end}\8\2\&{if}\1;\2\6
\&{tats}\C{the maximum value of \\{tok\_ptr} occurs just before parameter
popping}\6
\\{tok\_ptr}\?\1(\?\|z\?\2)\K\?\\{tok\_start}\?\1(\?\\{text\_ptr}\?\2)\?\par
\U section~82.\fi

\M89. When a parameter occurs in a replacement text, we treat it as a simple
macro in position (\\{name\_ptr}\?-\?1):

\Y\P\?\4\X89:Start scanning current macro parameter, \&{goto}\8\\{restart}\X\S%
\?\6
\\{push\_level}\?\1(\?\\{name\_ptr}\?-\?1\?\2)\?;\5
\&{goto}\8\\{restart}\par
\U section~84.\fi

\M90. Similarly, a \\{param} token encountered as we copy a parameter is
converted
into a simple macro call for \\{name\_ptr}\?-\?1. Some care is needed to handle
cases like \\{macro}\?(\? \#; \\{print}\?(\?\.{"\#)"}\?)\? \?)\?; the \.{\#}
token will have been
changed to \\{param} outside of strings, but we still must distinguish `real'
parentheses from those in strings.

\Y\P\4\D\\{app\_repl}\?\1(\?\#\?\2)\S\?\6
\&{if}\1\1\8\\{tok\_ptr}\?\1(\?\|z\?\2)=\?\\{max\_toks}\2\8\&{then}\6
\\{overflow}\?\1(\?\.{"token"}\?\2)\?;\2\6
\&{end}\8\2\&{if}\1;\6
\\{tok\_mem}\?\1(\?\|z\?,\39\?\\{tok\_ptr}\?\1(\?\|z\?\2)\2)\K\?\#;\5
\\{incr}\?\1(\?\\{tok\_ptr}\?\1(\?\|z\?\2)\2)\?\par
\Y\P\?\4\X90:Copy the parameter into \\{tok\_mem}\X\S\?\6
\\{bal}\?\K\?1;\5
\\{incr}\?\1(\?\\{cur\_byte}\?\2)\?;\C{skip the opening `\.('}\6
\1\&{loop}\6
\|b\?\K\?\\{tok\_mem}\?\1(\?\\{zo}\?,\39\?\\{cur\_byte}\?\2)\?;\5
\\{incr}\?\1(\?\\{cur\_byte}\?\2)\?;\6
\&{if}\1\1\8\|b\?=\?\\{param}\2\8\&{then}\6
\\{store\_two\_bytes}\?\1(\?\\{name\_ptr}\?+\bn{8}{77777}\2)\?;\6
\4\&{else}\6
\&{if}\1\1\8\|b\?\G\bn{8}{200}\?\2\8\&{then}\6
\\{app\_repl}\?\1(\?\|b\?\2)\?;\5
\|b\?\K\?\\{tok\_mem}\?\1(\?\\{zo}\?,\39\?\\{cur\_byte}\?\2)\?;\5
\\{incr}\?\1(\?\\{cur\_byte}\?\2)\?;\6
\4\&{else}\6
\&{case}\1\8\|b\1\8\&{is}\8\6
\4\&{when}\8\H{(}\KR\6
\\{incr}\?\1(\?\\{bal}\?\2)\?;\6
\4\&{when}\8\H{)}\KR\6
\\{decr}\?\1(\?\\{bal}\?\2)\?;\5
\?\&{exit}\8\&{when}\?\8\\{bal}\?=\?0;\6
\4\&{when}\8\\{character\_tok}\?\mathrel{\.|}\?\\{ascii\_tok}\KR\6
\\{app\_repl}\?\1(\?\|b\?\2)\?;\5
\|b\?\K\?\\{tok\_mem}\?\1(\?\\{zo}\?,\39\?\\{cur\_byte}\?\2)\?;\5
\\{incr}\?\1(\?\\{cur\_byte}\?\2)\?;\6
\4\&{when}\8\H{"}\KR\6
\1\&{repeat}\6
\\{app\_repl}\?\1(\?\|b\?\2)\?;\5
\|b\?\K\?\\{tok\_mem}\?\1(\?\\{zo}\?,\39\?\\{cur\_byte}\?\2)\?;\5
\\{incr}\?\1(\?\\{cur\_byte}\?\2)\?;\6
\&{if}\1\1\8\|b\?=\?\H{"}\?\W\?\\{tok\_mem}\?\1(\?\\{zo}\?,\39\?\\{cur\_byte}\?%
\2)=\?\H{"}\2\8\&{then}\6
\\{app\_repl}\?\1(\?\H{"}\?\2)\?;\5
\\{app\_repl}\?\1(\?\H{"}\?\2)\?;\5
\|b\?\K\?\\{tok\_mem}\?\1(\?\\{zo}\?,\39\?\\{cur\_byte}\?+\?1\?\2)\?;\5
\\{cur\_byte}\?\K\?\\{cur\_byte}\?+\?2;\2\6
\&{end}\8\2\&{if}\1;\2\6
\&{until}\?\8\1(\?\|b\?=\?\H{"}\?\2)\?;\C{copy string, don't change \\{bal}}\6
\4\&{when}\8\&{others}\8\KR\6
\&{null};\2\6
\4\2\&{end}\8\&{case};\2\6
\&{end}\8\2\&{if}\1;\6
\\{app\_repl}\?\1(\?\|b\?\2)\?;\2\6
\&{end}\8\2\&{if}\1;\2\6
\&{end}\8\&{loop}\par
\U section~87.\fi

\N91.  Producing the output.
The \\{get\_output} routine above handles most of the complexity of output
generation, but there are two further considerations that have a nontrivial
effect on \.{ATANGLE}'s algorithms.

First, we want to make sure that the output is broken into lines not
exceeding \\{line\_length} characters per line, where these breaks occur at
valid places (e.g., not in the middle of a string or a constant or an
identifier, not between `\./' and `\.=', not at a `\.{@\&}' position
where quantities are being joined together). Therefore we assemble the
output into a buffer before deciding where the line breaks will appear.
However, we make very little attempt to make ``logical'' line breaks that
would enhance the readability of the output; people are supposed to read
the input of \.{ATANGLE} or the \TeX ed output of \.{AWEAVE}, but not the
tangled-up output. The only concession to readability is that a break after
a semicolon will be made if possible, since commonly used ``pretty
printing'' routines give better results in such cases.

The following solution to these problems has been adopted: An array
\\{out\_buf} contains characters that have been generated but not yet output,
and there are three pointers into this array. One of these, \\{out\_ptr}, is
the number of characters currently in the buffer, and we will have
1\?\L\?\\{out\_ptr}\?\L\?\\{line\_length} most of the time. The second is %
\\{break\_ptr},
which is the largest value \?\L\?\\{out\_ptr} such that we are definitely
entitled
to end a line by outputting the characters \\{out\_buf}\?(\?1\?\to(\?\\{break%
\_ptr}\?-\?1\?))\?;
we will always have \\{break\_ptr}\?\L\?\\{line\_length}. Finally, \\{semi%
\_ptr} is either
zero or the largest known value of a legal break after a semicolon
on the current line; we will always have \\{semi\_ptr}\?\L\?\\{break\_ptr}.

\Y\P\?\4\X10:Types and objects visible from \\{data\_structures}\X\mathrel{+}\S%
\?\6
\&{subtype}\8\\{out\_buffer\_index}\8\&{is}\8\p{INTEGER}\8\&{range}\80\?\to\?%
\\{out\_buf\_size};\6
\\{out\_buf}\?:\?\\{buffer\_type}\?\1(\?\\{out\_buffer\_index}\?\2)\?;%
\C{assembled characters}\6
\\{out\_ptr}\?:\?\\{out\_buffer\_index};\C{first available place in \\{out%
\_buf}}\6
\\{break\_ptr}\?:\?\\{out\_buffer\_index};\C{last breaking place in \\{out%
\_buf}}\6
\\{semi\_ptr}\?:\?\\{out\_buffer\_index};\C{last semicolon breaking place in %
\\{out\_buf}}\par
\fi

\M92. Besides having those three pointers,
the output process is in one of several states:

\yskip\hang \\{num\_or\_id} means that the last item in the buffer is a number
or
identifier, hence a blank space or line break must be inserted if the next
item is also a number or identifier.

\yskip\hang \\{unbreakable} means that the last item in the buffer was followed
by the \.{@\&} operation that inhibits spaces between it and the next item.

\yskip\hang \\{misc} means none of the above.

\yskip\noindent
In all states the last item in the buffer lies between \\{break\_ptr} and
\\{out\_ptr}\?-\?1, inclusive.

Beside the states there is a boolean variable, called \\{comment\_output},
which is set true, if the material to be output is part of the meta-comments
between `\.{@\{}' and `\.{@\}}' and therefore has to be preceded by
`\.{--}'.

\Y\P\?\4\X10:Types and objects visible from \\{data\_structures}\X\mathrel{+}\S%
\?\6
\\{misc}\?:\?\&{constant}\8\\{eight\_bits}\?\K\?0;\C{state associated with
special characters}\6
\\{num\_or\_id}\?:\?\&{constant}\8\\{eight\_bits}\?\K\?1;\C{state associated
with numbers and identifiers}\6
\\{unbreakable}\?:\?\&{constant}\8\\{eight\_bits}\?\K\?2;\C{state associated
with \.{@\&}}\7
\\{out\_state}\?:\?\\{eight\_bits};\C{current status of partial output}\6
\\{comment\_output}\?:\?\p{BOOLEAN};\C{outputting meta-comment?}\par
\fi

\M93. During the output process, \\{line} will equal the number of the next
line
to be output.

\Y\P\?\4\X93:Initialize the output buffer\X\S\?\6
\\{out\_state}\?\K\?\\{misc};\5
\\{out\_ptr}\?\K\?0;\5
\\{break\_ptr}\?\K\?0;\5
\\{semi\_ptr}\?\K\?0;\5
\\{out\_buf}\?\1(\?0\?\2)\K\?0;\5
\\{line}\?\K\?1;\5
\\{comment\_output}\?\K\?\p{FALSE};\par
\U section~100.\fi

\M94. Here is a routine that is invoked when \\{out\_ptr}\?>\?\\{line\_length}
or when it is time to flush out the final line. The \\{flush\_buffer} procedure
often writes out the line up to the current \\{break\_ptr} position, then moves
the
remaining information to the front of \\{out\_buf}. However, it prefers to
write only up to \\{semi\_ptr}, if the residual line won't be too long.

\Y\P\4\D\\{check\_break}\?\S\?\6
\&{if}\1\1\8\\{out\_ptr}\?>\?\\{line\_length}\2\8\&{then}\6
\\{flush\_buffer}; \&{end} \&{if}\1\par
\Y\P\?\4\X81:Procedures in \\{output\_phase}\X\mathrel{+}\S\?\6
\&{procedure}\8\1\\{flush\_buffer}\2\8\&{is}\8\C{writes one line to output
file}\1\6
\|b\?:\?\\{out\_buffer\_index}\?\K\?\\{break\_ptr};\C{value of \\{break\_ptr}
upon entry}\6
\4\&{begin}\6
\&{if}\1\1\8\\{comment\_output}\2\8\&{then}\6
\p{TEXT\_IO}\?.\?\p{PUT}\?\1(\?\\{ada\_file}\?,\39\?\.{"--"}\?\2)\?;\2\6
\&{end}\8\2\&{if}\1;\6
\&{if}\1\1\?\8\1(\?\\{semi\_ptr}\?\I\?0\?\2)\W\1(\?\\{out\_ptr}\?-\?\\{semi%
\_ptr}\?\L\?\\{line\_length}\?\2)\?\2\8\&{then}\6
\\{break\_ptr}\?\K\?\\{semi\_ptr};\2\6
\&{end}\8\2\&{if}\1;\6
\1\&{for}\8\|k\?\in\?1\?\to\?\\{break\_ptr}\8\&{loop}\6
\p{TEXT\_IO}\?.\?\p{PUT}\?\1(\?\\{ada\_file}\?,\39\?\\{xchr}\?\1(\?\\{out\_buf}%
\?\1(\?\|k\?-\?1\?\2)\2)\2)\?;\2\6
\&{end}\8\&{loop};\6
\\{new\_ln}\?\1(\?\\{ada\_file}\?\2)\?;\5
\\{incr}\?\1(\?\\{line}\?\2)\?;\5
\\{incr}\?\1(\?\\{f\_line}\?\2)\?;\6
\&{if}\1\1\8\\{line}\?\mathbin{\&{mod}}\?100\?=\?0\2\8\&{then}\6
\\{print}\?\1(\?\.{"."}\?\2)\?;\6
\&{if}\1\1\8\\{line}\?\mathbin{\&{mod}}\?500\?=\?0\2\8\&{then}\6
\\{print\_int}\?\1(\?\\{line}\?,\39\?1\?\2)\?;\5
\\{update\_terminal};\C{progress report}\2\6
\&{end}\8\2\&{if}\1;\2\6
\&{end}\8\2\&{if}\1;\6
\&{if}\1\1\8\\{break\_ptr}\?<\?\\{out\_ptr}\2\8\&{then}\6
\&{if}\1\1\8\\{out\_buf}\?\1(\?\\{break\_ptr}\?\2)=\?\H{\ }\2\8\&{then}\6
\\{incr}\?\1(\?\\{break\_ptr}\?\2)\?;\C{drop space at break}\6
\&{if}\1\1\8\\{break\_ptr}\?>\?\|b\2\8\&{then}\6
\|b\?\K\?\\{break\_ptr};\2\6
\&{end}\8\2\&{if}\1;\2\6
\&{end}\8\2\&{if}\1;\6
\\{out\_buf}\?\1(\?0\?\to\?\\{out\_ptr}\?-\?\\{break\_ptr}\?-\?1\?\2)\K\?\\{out%
\_buf}\?\1(\?\\{break\_ptr}\?\to\?\\{out\_ptr}\?-\?1\?\2)\?;\2\6
\&{end}\8\2\&{if}\1;\6
\\{out\_ptr}\?\K\?\\{out\_ptr}\?-\?\\{break\_ptr};\5
\\{break\_ptr}\?\K\?\|b\?-\?\\{break\_ptr};\5
\\{semi\_ptr}\?\K\?0;\6
\&{if}\1\1\8\\{out\_ptr}\?>\?\\{line\_length}\2\8\&{then}\6
\\{err\_print}\?\1(\?\.{"!\ Long\ line\ must\ be\ truncated"}\?\2)\?;\6
\\{out\_ptr}\?\K\?\\{line\_length};\2\6
\&{end}\8\2\&{if}\1;\2\6
\&{end}\8\\{flush\_buffer};\par
\fi

\M95. \P\?\X95:Empty the last line from the buffer\X\S\?\6
\\{break\_ptr}\?\K\?\\{out\_ptr};\5
\\{semi\_ptr}\?\K\?0;\5
\\{flush\_buffer};\6
\&{if}\1\1\8\\{brace\_level}\?\I\?0\2\8\&{then}\6
\\{new\_ln};\5
\\{print}\?\1(\?\.{"!\ Program\ ended\ at\ brace\ level\ "}\?\2)\?;\5
\\{print\_int}\?\1(\?\\{brace\_level}\?,\39\?1\?\2)\?;\5
\\{error};\2\6
\&{end}\8\2\&{if}\1;\par
\U section~100.\fi

\M96. Another simple and useful routine appends the decimal equivalent of
a nonnegative integer to the output buffer.

\Y\P\4\D\\{app}\?\1(\?\#\?\2)\S\?\\{out\_buf}\?\1(\?\\{out\_ptr}\?\2)\K\?\#; %
\\{incr}\?\1(\?\\{out\_ptr}\?\2)\C{append a single character}\6
\?\par
\Y\P\?\4\X81:Procedures in \\{output\_phase}\X\mathrel{+}\S\?\6
\&{procedure}\8\1\\{send\_val}\?(\1\?\|v\?:\?\p{INTEGER}\?\2)\?\2\8\&{is}\8%
\C{puts \|v into buffer, assumes \|v\?\G\?0}\1\6
\|k\?:\?\\{out\_buffer\_index}\?\K\?\\{out\_buf\_size};\C{index into \\{out%
\_buf}}\6
\|u\?:\?\p{INTEGER}\?\K\?\|v;\6
\4\&{begin}\C{first we put the digits at the very end of \\{out\_buf}}\6
\1\&{repeat}\6
\\{out\_buf}\?\1(\?\|k\?\2)\K\?\|u\?\mathbin{\&{mod}}\?10;\5
\|u\?\K\?\|u\?/\?10;\5
\\{decr}\?\1(\?\|k\?\2)\?;\2\6
\&{until}\?\8\1(\?\|u\?=\?0\?\2)\?;\6
\1\&{repeat}\6
\\{incr}\?\1(\?\|k\?\2)\?;\5
\\{app}\?\1(\?\\{out\_buf}\?\1(\?\|k\?\2)+\?\H{0}\?\2)\?;\2\6
\&{until}\?\8\1(\?\|k\?=\?\\{out\_buf\_size}\?\2)\?;\C{then we append them,
most significant first}\2\6
\&{end}\8\\{send\_val};\par
\fi

\M97. The output states are kept up to date by the \\{send\_out} procedure,
which
has two parameters: \|t tells the type of information being sent and
\|v contains the information proper. Some information may also be passed
in the array \\{out\_contrib}.

\yskip\hang If \|t\?=\?\\{misc} then \|v is a character to be output.

\hang If \|t\?=\?\\{str} then \|v is the length of a string or something like
`\.{/=}' in \\{out\_contrib}.

\hang If \|t\?=\?\\{ident} then \|v is the length of an identifier in \\{out%
\_contrib}.

\hang If \|t\?=\?\\{frac} then \|v is the length of a fraction and/or exponent
in
\\{out\_contrib}.

\Y\P\?\4\X10:Types and objects visible from \\{data\_structures}\X\mathrel{+}\S%
\?\6
\\{str}\?:\?\&{constant}\8\\{eight\_bits}\?\K\?1;\C{\\{send\_out} code for a
string}\6
\\{ident}\?:\?\&{constant}\8\\{eight\_bits}\?\K\?2;\C{\\{send\_out} code for an
identifier}\6
\\{frac}\?:\?\&{constant}\8\\{eight\_bits}\?\K\?3;\C{\\{send\_out} code for a
fraction}\7
\\{out\_contrib}\?:\?\\{buffer\_type}\?\1(\?1\?\to\?\\{line\_length}\?\2)\?;%
\C{a contribution to \\{out\_buf}}\par
\fi

\M98. A slightly subtle point in the following code is that the user may ask
for a \\{join} operation (i.e., \.{@\&}) following whatever is being sent
out.  We will see later that \\{join} is implemented in part by calling
\\{send\_out}\?(\?\\{frac}\?,\?0\?)\?.

\Y\P\?\4\X81:Procedures in \\{output\_phase}\X\mathrel{+}\S\?\6
\&{procedure}\8\1\\{send\_out}\?(\?\|t\?:\?\\{eight\_bits};\?\,\35\1\?\|v\?:\?%
\\{sixteen\_bits}\?\2)\?\2\8\&{is}\8\C{outputs \|v of type \|t}\1\6
\4\&{begin}\6
\X99:Get the buffer ready for appending the new information\X;\6
\&{if}\1\1\8\|t\?\I\?\\{misc}\2\8\&{then}\6
\1\&{for}\8\|k\?\in\?1\?\to\?\|v\8\&{loop}\6
\\{app}\?\1(\?\\{out\_contrib}\?\1(\?\|k\?\2)\2)\?;\2\6
\&{end}\8\&{loop};\6
\4\&{else}\6
\\{app}\?\1(\?\|v\?\2)\?;\2\6
\&{end}\8\2\&{if}\1;\6
\\{check\_break};\6
\&{if}\1\1\?\8\1(\?\|t\?=\?\\{misc}\?\2)\W\1(\?\|v\?=\?\H{;}\?\2)\?\2\8\&{then}%
\6
\\{semi\_ptr}\?\K\?\\{out\_ptr};\5
\\{break\_ptr}\?\K\?\\{out\_ptr};\2\6
\&{end}\8\2\&{if}\1;\6
\&{if}\1\1\8\|t\?\G\?\\{ident}\2\8\&{then}\6
\\{out\_state}\?\K\?\\{num\_or\_id};\C{\|t\?=\?\\{ident} or \\{frac}}\6
\4\&{else}\6
\\{out\_state}\?\K\?\\{misc};\C{\|t\?=\?\\{str} or \\{misc}}\2\6
\&{end}\8\2\&{if}\1;\2\6
\&{end}\8\\{send\_out};\par
\fi

\M99. \P\?\X99:Get the buffer ready for appending the new information\X\S\?\6
\4\BL\\{restart}\EL\8\&{case}\1\8\\{out\_state}\1\8\&{is}\8\6
\4\&{when}\8\\{num\_or\_id}\KR\6
\&{if}\1\1\8\|t\?\I\?\\{frac}\2\8\&{then}\6
\\{break\_ptr}\?\K\?\\{out\_ptr};\6
\&{if}\1\1\8\|t\?=\?\\{ident}\2\8\&{then}\6
\\{app}\?\1(\?\H{\ }\?\2)\?;\2\6
\&{end}\8\2\&{if}\1;\2\6
\&{end}\8\2\&{if}\1;\6
\4\&{when}\8\\{misc}\KR\6
\&{if}\1\1\8\|t\?\I\?\\{frac}\2\8\&{then}\6
\\{break\_ptr}\?\K\?\\{out\_ptr};\2\6
\&{end}\8\2\&{if}\1;\6
\4\&{when}\8\&{others}\8\KR\6
\&{null};\C{this is for \\{unbreakable} state}\2\6
\4\2\&{end}\8\&{case}\par
\U section~98.\fi

\N100.  The big output switch.
To complete the output process, we need a routine that takes the results
of \\{get\_output} and feeds them to \\{send\_out} or \\{send\_val}.
This procedure `\\{send\_the\_output}' will be invoked just once, as follows:

\Y\P\?\4\X100:Phase II: Output the contents of the compressed tables\X\S\?\6
\&{if}\1\1\8\\{text\_link}\?\1(\?0\?\2)=\?0\2\8\&{then}\6
\\{print\_nl}\?\1(\?\.{"!\ No\ output\ was\ specified."}\?\2)\?;\5
\\{mark\_harmless};\6
\4\&{else}\6
\\{print\_nl}\?\1(\?\.{"Writing\ the\ output\ file"}\?\2)\?;\5
\\{update\_terminal};\6
\X80:Initialize the output stacks\X\6
\X93:Initialize the output buffer\X\6
\\{send\_the\_output};\6
\X95:Empty the last line from the buffer\X\6
\\{print\_nl}\?\1(\?\.{"Done."}\?\2)\?;\2\6
\&{end}\8\2\&{if}\1\par
\U section~114.\fi

\M101. A many-way switch is used to send the output:

\Y\P\?\4\X81:Procedures in \\{output\_phase}\X\mathrel{+}\S\?\6
\&{procedure}\8\1\\{send\_the\_output}\2\8\&{is}\8\1\6
\\{cur\_char}\?:\?\\{eight\_bits};\C{the latest character received}\6
\|k\?:\?\p{INTEGER}\8\&{range}\80\?\to\?\\{line\_length};\C{index into \\{out%
\_contrib}}\6
\|j\?:\?\\{byte\_index};\C{index into \\{byte\_mem}}\6
\|w\?:\?\\{w\_range};\C{segment of \\{byte\_mem}}\6
\4\&{begin}\6
\1\&{while}\8\\{stack\_ptr}\?>\?0\8\&{loop}\6
\\{cur\_char}\?\K\?\\{get\_output};\6
\4\BL\\{reswitch}\EL\8\&{case}\1\8\\{cur\_char}\1\8\&{is}\8\6
\4\&{when}\80\KR\8\&{null};\C{this case might arise if output ends
unexpectedly}\6
\hbox{\4}\X104:Cases related to identifiers\X\6
\hbox{\4}\X110:Cases related to constants, possibly leading to \\{get%
\_fraction} or \\{reswitch}\X\6
\hbox{\4}\X102:Cases like \.{/=} and \.{:=}\X\6
\4\&{when}\8\H{"}\KR\8\X105:Send a string, \&{goto}\8\\{reswitch}\X;\6
\4\&{when}\?\8\X103:Other printable characters\X\?\KR\8\\{send\_out}\?\1(\?%
\\{misc}\?,\39\?\\{cur\_char}\?\2)\?;\6
\hbox{\4}\X112:Cases involving \.{@\{} and \.{@\}}\X\6
\4\&{when}\8\\{join}\KR\8\\{send\_out}\?\1(\?\\{frac}\?,\39\?0\?\2)\?;\5
\\{out\_state}\?\K\?\\{unbreakable};\6
\4\&{when}\8\\{verbatim}\KR\8\X106:Send verbatim string\X;\6
\4\&{when}\8\\{output\_tok}\KR\8\X108:Get the \\{file\_name} and redirect
further output\X;\6
\4\&{when}\8\\{force\_line}\KR\8\X113:Force a line break\X;\6
\4\&{when}\8\&{others}\8\KR\6
\\{new\_ln};\5
\\{print}\?\1(\?\.{"!\ Can\'t\ output\ ASCII\ code\ "}\?\2)\?;\5
\\{print\_int}\?\1(\?\\{cur\_char}\?,\39\?1\?\2)\?;\5
\\{error};\2\6
\4\2\&{end}\8\&{case};\6
\&{goto}\8\\{continue};\6
\4\BL\\{get\_fraction}\EL\8\X111:Special code to finish real constants\X;\6
\4\BL\\{continue}\EL\8\&{null};\2\6
\&{end}\8\&{loop};\2\6
\&{end}\8\\{send\_the\_output};\par
\fi

\M102. \P\?\X102:Cases like \.{/=} and \.{:=}\X\S\?\6
\4\&{when}\8\\{and\_sign}\KR\8\\{out\_contrib}\?\1(\?1\?\to\?3\?\2)\K\1(\?\H{A}%
\?,\39\?\H{N}\?,\39\?\H{D}\?\2)\?;\5
\\{send\_out}\?\1(\?\\{ident}\?,\39\?3\?\2)\?;\6
\4\&{when}\8\\{not\_sign}\KR\8\\{out\_contrib}\?\1(\?1\?\to\?3\?\2)\K\1(\?\H{N}%
\?,\39\?\H{O}\?,\39\?\H{T}\?\2)\?;\5
\\{send\_out}\?\1(\?\\{ident}\?,\39\?3\?\2)\?;\6
\4\&{when}\8\\{set\_element\_sign}\KR\8\\{out\_contrib}\?\1(\?1\?\to\?2\?\2)\K%
\1(\?\H{I}\?,\39\?\H{N}\?\2)\?;\5
\\{send\_out}\?\1(\?\\{ident}\?,\39\?2\?\2)\?;\6
\4\&{when}\8\\{or\_sign}\KR\8\\{out\_contrib}\?\1(\?1\?\to\?2\?\2)\K\1(\?\H{O}%
\?,\39\?\H{R}\?\2)\?;\5
\\{send\_out}\?\1(\?\\{ident}\?,\39\?2\?\2)\?;\6
\4\&{when}\8\\{left\_arrow}\KR\8\\{out\_contrib}\?\1(\?1\?\to\?2\?\2)\K\1(\?%
\H{:}\?,\39\?\H{=}\?\2)\?;\5
\\{send\_out}\?\1(\?\\{str}\?,\39\?2\?\2)\?;\6
\4\&{when}\8\\{not\_equal}\KR\8\\{out\_contrib}\?\1(\?1\?\to\?2\?\2)\K\1(\?%
\H{/}\?,\39\?\H{=}\?\2)\?;\5
\\{send\_out}\?\1(\?\\{str}\?,\39\?2\?\2)\?;\6
\4\&{when}\8\\{less\_or\_equal}\KR\8\\{out\_contrib}\?\1(\?1\?\to\?2\?\2)\K\1(%
\?\H{<}\?,\39\?\H{=}\?\2)\?;\5
\\{send\_out}\?\1(\?\\{str}\?,\39\?2\?\2)\?;\6
\4\&{when}\8\\{greater\_or\_equal}\KR\8\\{out\_contrib}\?\1(\?1\?\to\?2\?\2)\K%
\1(\?\H{>}\?,\39\?\H{=}\?\2)\?;\5
\\{send\_out}\?\1(\?\\{str}\?,\39\?2\?\2)\?;\6
\4\&{when}\8\\{equivalence\_sign}\KR\8\\{out\_contrib}\?\1(\?1\?\to\?2\?\2)\K%
\1(\?\H{=}\?,\39\?\H{=}\?\2)\?;\5
\\{send\_out}\?\1(\?\\{str}\?,\39\?2\?\2)\?;\6
\4\&{when}\8\\{double\_dot}\KR\8\\{out\_contrib}\?\1(\?1\?\to\?2\?\2)\K\1(\?%
\H{.}\?,\39\?\H{.}\?\2)\?;\5
\\{send\_out}\?\1(\?\\{str}\?,\39\?2\?\2)\?;\6
\4\&{when}\8\\{right\_arrow}\KR\8\\{out\_contrib}\?\1(\?1\?\to\?2\?\2)\K\1(\?%
\H{=}\?,\39\?\H{>}\?\2)\?;\5
\\{send\_out}\?\1(\?\\{str}\?,\39\?2\?\2)\?;\6
\4\&{when}\8\\{box\_sign}\KR\8\\{out\_contrib}\?\1(\?1\?\to\?2\?\2)\K\1(\?\H{<}%
\?,\39\?\H{>}\?\2)\?;\5
\\{send\_out}\?\1(\?\\{str}\?,\39\?2\?\2)\?;\6
\4\&{when}\8\\{left\_label\_bracket}\KR\8\\{out\_contrib}\?\1(\?1\?\to\?2\?\2)%
\K\1(\?\H{<}\?,\39\?\H{<}\?\2)\?;\5
\\{send\_out}\?\1(\?\\{str}\?,\39\?2\?\2)\?;\6
\4\&{when}\8\\{right\_label\_bracket}\KR\8\\{out\_contrib}\?\1(\?1\?\to\?2\?\2)%
\K\1(\?\H{>}\?,\39\?\H{>}\?\2)\?;\5
\\{send\_out}\?\1(\?\\{str}\?,\39\?2\?\2)\?;\6
\4\&{when}\8\\{double\_star}\KR\8\\{out\_contrib}\?\1(\?1\?\to\?2\?\2)\K\1(\?%
\H{*}\?,\39\?\H{*}\?\2)\?;\5
\\{send\_out}\?\1(\?\\{str}\?,\39\?2\?\2)\?;\6
\4\&{when}\8\\{character\_tok}\KR\8\\{out\_contrib}\?\1(\?1\?\to\?3\?\2)\K\1(\?%
\H{\'}\?,\39\?\\{get\_output}\?,\39\?\H{\'}\?\2)\?;\5
\\{send\_out}\?\1(\?\\{str}\?,\39\?3\?\2)\?;\6
\4\&{when}\8\\{ascii\_tok}\KR\6
\&{if}\1\1\8\\{out\_state}\?=\?\\{num\_or\_id}\2\8\&{then}\6
\\{send\_out}\?\1(\?\\{misc}\?,\39\?\H{\ }\?\2)\?;\2\6
\&{end}\8\2\&{if}\1;\6
\\{send\_val}\?\1(\?\\{get\_output}\?\2)\?;\5
\\{out\_state}\?\K\?\\{num\_or\_id};\par
\U section~101.\fi

\M103. Please don't ask how all of the following characters can actually get
through \.{ATANGLE} outside of strings. It seems that \H{\{}
cannot actually occur at this point of the program, but it has
been included just in case \.{ATANGLE} changes.

\Y\P\?\4\X103:Other printable characters\X\S\?\6
\H{!}\?\mathrel{\.|}\?\H{\#}\?\mathrel{\.|}\?\H{\$}\?\mathrel{\.|}\?\H{\%}\?%
\mathrel{\.|}\?\H{\&}\?\mathrel{\.|}\?\H{\'}\?\mathrel{\.|}\?\H{(}\?\mathrel{%
\.|}\?\H{)}\?\mathrel{\.|}\?\H{*}\?\mathrel{\.|}\?\H{,}\?\mathrel{\.|}\?\H{/}\?%
\mathrel{\.|}\?\H{:}\?\mathrel{\.|}\?\H{;}\?\mathrel{\.|}\?\H{<}\?\mathrel{\.|}%
\?\H{=}\?\mathrel{\.|}\?\H{>}\?\mathrel{\.|}\?\H{?}\?\mathrel{\.|}\?\H{@}\?%
\mathrel{\.|}\?\H{[}\?\mathrel{\.|}\?\H{\\}\?\mathrel{\.|}\?\H{]}\?\mathrel{%
\.|}\?\H{\^}\?\mathrel{\.|}\?\H{\_}\?\mathrel{\.|}\?\H{\`}\?\mathrel{\.|}\?\H{%
\{}\?\mathrel{\.|}\?\H{|}\?\mathrel{\.|}\?\H{+}\?\mathrel{\.|}\?\H{-}\?\6
\?\par
\U section~101.\fi

\M104. Single-character identifiers represent themselves, while longer ones
appear in \\{byte\_mem}. All must be converted to uppercase,
with underlines removed. Extremely long identifiers must be chopped.

\Y\P\?\4\X104:Cases related to identifiers\X\S\?\6
\4\&{when}\8\H{A}\?\to\?\H{Z}\KR\8\\{out\_contrib}\?\1(\?1\?\2)\K\?\\{cur%
\_char};\5
\\{send\_out}\?\1(\?\\{ident}\?,\39\?1\?\2)\?;\6
\4\&{when}\8\H{a}\?\to\?\H{z}\KR\8\\{out\_contrib}\?\1(\?1\?\2)\K\?\\{cur%
\_char}\?-\bn{8}{40}\?;\5
\\{send\_out}\?\1(\?\\{ident}\?,\39\?1\?\2)\?;\6
\4\&{when}\8\\{identifier}\KR\8\|k\?\K\?0;\5
\|j\?\K\?\\{byte\_start}\?\1(\?\\{cur\_val}\?\2)\?;\5
\|w\?\K\?\\{cur\_val}\?\mathbin{\&{mod}}\?\\{ww};\6
\1\&{while}\?\8\1(\?\|k\?<\?\\{max\_id\_length}\?\2)\W\1(\?\|j\?<\?\\{byte%
\_start}\?\1(\?\\{cur\_val}\?+\?\\{ww}\?\2)\2)\?\8\&{loop}\6
\\{incr}\?\1(\?\|k\?\2)\?;\5
\\{out\_contrib}\?\1(\?\|k\?\2)\K\?\\{byte\_mem}\?\1(\?\|w\?,\39\?\|j\?\2)\?;\5
\\{incr}\?\1(\?\|j\?\2)\?;\6
\&{if}\1\1\8\\{out\_contrib}\?\1(\?\|k\?\2)\G\?\H{a}\2\8\&{then}\6
\\{out\_contrib}\?\1(\?\|k\?\2)\K\?\\{out\_contrib}\?\1(\?\|k\?\2)-\bn{8}{40}%
\?;\2\6
\&{end}\8\2\&{if}\1;\2\6
\&{end}\8\&{loop};\6
\\{send\_out}\?\1(\?\\{ident}\?,\39\?\|k\?\2)\?;\par
\U section~101.\fi

\M105. After sending a string, we need to look ahead at the next character, in
order
to see if there were two consecutive single-quote marks. Afterwards we go to
\\{reswitch} to process the next character.

\Y\P\?\4\X105:Send a string, \&{goto}\8\\{reswitch}\X\S\?\6
\|k\?\K\?1;\5
\\{out\_contrib}\?\1(\?1\?\2)\K\?\H{"};\6
\1\&{repeat}\6
\&{if}\1\1\8\|k\?<\?\\{line\_length}\2\8\&{then}\6
\\{incr}\?\1(\?\|k\?\2)\?;\2\6
\&{end}\8\2\&{if}\1;\6
\\{out\_contrib}\?\1(\?\|k\?\2)\K\?\\{get\_output};\2\6
\&{until}\?\8\1(\1(\?\\{out\_contrib}\?\1(\?\|k\?\2)=\?\H{"}\?\2)\V\1(\?%
\\{stack\_ptr}\?=\?0\?\2)\2)\?;\6
\&{if}\1\1\8\|k\?=\?\\{line\_length}\2\8\&{then}\6
\\{err\_print}\?\1(\?\.{"!\ String\ too\ long"}\?\2)\?;\2\6
\&{end}\8\2\&{if}\1;\6
\\{send\_out}\?\1(\?\\{str}\?,\39\?\|k\?\2)\?;\5
\\{cur\_char}\?\K\?\\{get\_output};\6
\&{if}\1\1\8\\{cur\_char}\?=\?\H{"}\2\8\&{then}\6
\\{out\_state}\?\K\?\\{unbreakable};\2\6
\&{end}\8\2\&{if}\1;\6
\&{goto}\8\\{reswitch}\par
\U section~101.\fi

\M106. Sending a verbatim string is similar, but we don't have to look ahead.

\Y\P\?\4\X106:Send verbatim string\X\S\?\6
\|k\?\K\?0;\6
\1\&{repeat}\6
\&{if}\1\1\8\|k\?<\?\\{line\_length}\2\8\&{then}\6
\\{incr}\?\1(\?\|k\?\2)\?;\2\6
\&{end}\8\2\&{if}\1;\6
\\{out\_contrib}\?\1(\?\|k\?\2)\K\?\\{get\_output};\2\6
\&{until}\?\8\1(\1(\?\\{out\_contrib}\?\1(\?\|k\?\2)=\?\\{verbatim}\?\2)\V\1(\?%
\\{stack\_ptr}\?=\?0\?\2)\2)\?;\6
\&{if}\1\1\8\|k\?=\?\\{line\_length}\2\8\&{then}\6
\\{err\_print}\?\1(\?\.{"!\ Verbatim\ string\ too\ long"}\?\2)\?;\2\6
\&{end}\8\2\&{if}\1;\6
\\{send\_out}\?\1(\?\\{str}\?,\39\?\|k\?-\?1\?\2)\?\par
\U section~101.\fi

\M107. Getting the new output file name is similar to sending a verbatim
string;
if there is a legal name, the file will be created and the variables
\\{file\_name} and \\{fn\_length} will be set accordingly and \\{f\_line} is
set
to~1.

\Y\P\?\4\X10:Types and objects visible from \\{data\_structures}\X\mathrel{+}\S%
\?\6
\\{file\_name}\?:\?\p{STRING}\?\1(\?1\?\to\?80\?\2)\?;\6
\\{fn\_length}\?:\?\p{INTEGER}\8\&{range}\81\?\to\?80;\6
\\{f\_line}\?:\?\p{INTEGER};\par
\fi

\M108. Before we redirect output we will empty the output buffer.

\Y\P\?\4\X108:Get the \\{file\_name} and redirect further output\X\S\?\6
\X113:Force a line break\X;\6
\|k\?\K\?0;\6
\1\&{loop}\6
\\{cur\_char}\?\K\?\\{get\_output};\5
\?\&{exit}\8\&{when}\8\1(\?\\{cur\_char}\?=\?\\{output\_tok}\?\2)\V\1(\?%
\\{stack\_ptr}\?=\?0\?\2)\?;\6
\\{incr}\?\1(\?\|k\?\2)\?;\5
\\{file\_name}\?\1(\?\|k\?\2)\K\?\\{xchr}\?\1(\?\\{cur\_char}\?\2)\?;\2\6
\&{end}\8\&{loop};\6
\&{if}\1\1\8\|k\?>\?0\2\8\&{then}\6
\p{TEXT\_IO}\?.\?\p{CLOSE}\?\1(\?\\{ada\_file}\?\2)\?;\5
\\{fn\_length}\?\K\?\|k;\6
\&{if}\1\1\?\8\R\?\\{create\_ada\_file}\?\1(\?\\{file\_name}\?\1(\?1\?\to\?\|k%
\?\2)\2)\?\2\8\&{then}\6
\\{fatal\_error}\?\1(\?\.{"!\ Could\ not\ open\ file\ \`"}\?\mathbin{\.\&}\?%
\\{file\_name}\?\1(\?1\?\to\?\|k\?\2)\mathbin{\.\&}\?\.{"\'."}\?\2)\?;\2\6
\&{end}\8\2\&{if}\1;\6
\\{f\_line}\?\K\?1;\6
\4\&{else}\6
\\{err\_print}\?\1(\?\.{"!\ File\ name\ should\ have\ length\ >\ 0"}\?\2)\?;\2\6
\&{end}\8\2\&{if}\1\par
\U section~101.\fi

\M109. In the cases related to numeric literals there is some care needed,
because of the ``based literal'', which \ADA\ supports; a global
variable \\{scanning\_based\_number} tells whether there is more to do.

\Y\P\?\4\X10:Types and objects visible from \\{data\_structures}\X\mathrel{+}\S%
\?\6
\\{scanning\_based\_number}\?:\?\p{BOOLEAN}\?\K\?\p{FALSE};\par
\fi

\M110. \P\?\X110:Cases related to constants, possibly leading to \\{get%
\_fraction} or \\{reswitch}\X\S\?\6
\4\&{when}\8\H{0}\?\to\?\H{9}\KR\6
\\{scanning\_based\_number}\?\K\?\p{FALSE};\5
\|k\?\K\?0;\6
\1\&{loop}\6
\&{if}\1\1\8\\{cur\_char}\?\I\?\H{\_}\2\8\&{then}\6
\\{incr}\?\1(\?\|k\?\2)\?;\5
\\{out\_contrib}\?\1(\?\|k\?\2)\K\?\\{cur\_char};\5
\\{cur\_char}\?\K\?\\{get\_output};\6
\&{if}\1\1\8\\{cur\_char}\?=\?\H{\#}\2\8\&{then}\6
\\{scanning\_based\_number}\?\K\R\?\\{scanning\_based\_number};\6
\4\&{else}\6
\?\&{exit}\8\&{when}\8\1(\1(\R\?\\{scanning\_based\_number}\?\2)\W\1(\?\\{cur%
\_char}\?\notin\?\H{0}\?\to\?\H{9}\?\2)\W\1(\?\\{cur\_char}\?\I\?\H{\_}\?\2)\2)%
\?;\2\6
\&{end}\8\2\&{if}\1;\6
\4\&{else}\6
\\{cur\_char}\?\K\?\\{get\_output};\2\6
\&{end}\8\2\&{if}\1;\2\6
\&{end}\8\&{loop};\6
\\{send\_out}\?\1(\?\\{ident}\?,\39\?\|k\?\2)\?;\5
\|k\?\K\?0;\6
\&{if}\1\1\8\\{cur\_char}\?=\?\H{e}\2\8\&{then}\6
\\{cur\_char}\?\K\?\H{E};\2\6
\&{end}\8\2\&{if}\1;\6
\&{if}\1\1\8\\{cur\_char}\?=\?\H{E}\2\8\&{then}\6
\&{goto}\8\\{get\_fraction};\6
\4\&{else}\6
\&{goto}\8\\{reswitch};\2\6
\&{end}\8\2\&{if}\1;\6
\4\&{when}\8\H{.}\KR\6
\|k\?\K\?1;\5
\\{out\_contrib}\?\1(\?1\?\2)\K\?\H{.};\5
\\{cur\_char}\?\K\?\\{get\_output};\6
\&{if}\1\1\8\\{cur\_char}\?=\?\H{.}\2\8\&{then}\6
\\{out\_contrib}\?\1(\?2\?\2)\K\?\H{.};\5
\\{send\_out}\?\1(\?\\{str}\?,\39\?2\?\2)\?;\6
\4\&{elsif}\1\8\\{cur\_char}\?\in\?\H{0}\?\to\?\H{9}\2\8\&{then}\6
\&{goto}\8\\{get\_fraction};\6
\4\&{else}\6
\\{send\_out}\?\1(\?\\{misc}\?,\39\?\H{.}\?\2)\?;\5
\&{goto}\8\\{reswitch};\2\6
\&{end}\8\2\&{if}\1;\par
\U section~101.\fi

\M111. The following code appears at label `\\{get\_fraction}', when we want to
scan to the end of a real constant. The first \|k characters of a fraction
have already been placed in \\{out\_contrib}, and \\{cur\_char} is the next
character.

\Y\P\?\4\X111:Special code to finish real constants\X\S\?\6
\1\&{repeat}\6
\&{if}\1\1\8\|k\?<\?\\{line\_length}\2\8\&{then}\6
\\{incr}\?\1(\?\|k\?\2)\?;\2\6
\&{end}\8\2\&{if}\1;\6
\\{out\_contrib}\?\1(\?\|k\?\2)\K\?\\{cur\_char};\5
\\{cur\_char}\?\K\?\\{get\_output};\6
\&{if}\1\1\?\8\1(\?\\{out\_contrib}\?\1(\?\|k\?\2)=\?\H{E}\?\2)\W\1(\1(\?\\{cur%
\_char}\?=\?\H{+}\?\2)\V\1(\?\\{cur\_char}\?=\?\H{-}\?\2)\2)\?\2\8\&{then}\6
\&{if}\1\1\8\|k\?<\?\\{line\_length}\2\8\&{then}\6
\\{incr}\?\1(\?\|k\?\2)\?;\2\6
\&{end}\8\2\&{if}\1;\6
\\{out\_contrib}\?\1(\?\|k\?\2)\K\?\\{cur\_char};\5
\\{cur\_char}\?\K\?\\{get\_output};\6
\4\&{elsif}\1\8\\{cur\_char}\?=\?\H{e}\2\8\&{then}\6
\\{cur\_char}\?\K\?\H{E};\2\6
\&{end}\8\2\&{if}\1;\2\6
\&{until}\?\8\1(\1(\?\\{cur\_char}\?\I\?\H{E}\?\2)\W\1(\?\\{cur\_char}\?\notin%
\?\H{0}\?\to\?\H{9}\?\2)\2)\?;\6
\&{if}\1\1\8\|k\?=\?\\{line\_length}\2\8\&{then}\6
\\{err\_print}\?\1(\?\.{"!\ Fraction\ too\ long"}\?\2)\?;\2\6
\&{end}\8\2\&{if}\1;\6
\\{send\_out}\?\1(\?\\{frac}\?,\39\?\|k\?\2)\?;\5
\&{goto}\8\\{reswitch}\par
\U section~101.\fi

\M112. Here the variable \\{comment\_output} is set corresponding to the value
of
\\{brace\_level}; further after each comment a line break is inserted.

\Y\P\?\4\X112:Cases involving \.{@\{} and \.{@\}}\X\S\?\6
\4\&{when}\8\\{begin\_comment}\KR\6
\X113:Force a line break\X;\6
\\{incr}\?\1(\?\\{brace\_level}\?\2)\?;\5
\\{comment\_output}\?\K\?\p{TRUE};\6
\4\&{when}\8\\{end\_comment}\KR\6
\&{if}\1\1\8\\{brace\_level}\?>\?0\2\8\&{then}\6
\\{decr}\?\1(\?\\{brace\_level}\?\2)\?;\6
\&{if}\1\1\8\\{brace\_level}\?=\?0\2\8\&{then}\6
\X113:Force a line break\X;\6
\\{comment\_output}\?\K\?\p{FALSE};\2\6
\&{end}\8\2\&{if}\1;\6
\4\&{else}\6
\\{err\_print}\?\1(\?\.{"!\ Extra\ @\}"}\?\2)\?;\2\6
\&{end}\8\2\&{if}\1;\6
\4\&{when}\8\\{module\_number}\KR\6
\\{semi\_ptr}\?\K\?0;\5
\\{out\_contrib}\?\1(\?1\?\to\?2\?\2)\K\1(\?\H{-}\?,\39\?\H{-}\?\2)\?;\5
\\{send\_out}\?\1(\?\\{str}\?,\39\?2\?\2)\?;\6
\1\&{while}\8\\{cur\_char}\?=\?\\{module\_number}\8\&{loop}\6
\\{out\_state}\?\K\?\\{unbreakable};\5
\\{send\_out}\?\1(\?\\{misc}\?,\39\?\H{\{}\?\2)\?;\5
\\{out\_state}\?\K\?\\{unbreakable};\6
\&{if}\1\1\8\\{cur\_val}\?<\?0\2\8\&{then}\6
\\{send\_out}\?\1(\?\\{misc}\?,\39\?\H{:}\?\2)\?;\5
\\{out\_state}\?\K\?\\{unbreakable};\5
\\{send\_val}\?\1(-\?\\{cur\_val}\?\2)\?;\6
\4\&{else}\6
\\{send\_val}\?\1(\?\\{cur\_val}\?\2)\?;\5
\\{out\_state}\?\K\?\\{unbreakable};\5
\\{send\_out}\?\1(\?\\{misc}\?,\39\?\H{:}\?\2)\?;\2\6
\&{end}\8\2\&{if}\1;\6
\\{out\_state}\?\K\?\\{unbreakable};\5
\\{send\_out}\?\1(\?\\{misc}\?,\39\?\H{\}}\?\2)\?;\5
\\{cur\_char}\?\K\?\\{get\_output};\2\6
\&{end}\8\&{loop};\6
\X113:Force a line break\X;\6
\&{goto}\8\\{reswitch};\par
\U section~101.\fi

\M113. \P\?\X113:Force a line break\X\S\?\6
\\{send\_out}\?\1(\?\\{str}\?,\39\?0\?\2)\?;\C{normalize the buffer}\6
\1\&{while}\8\\{out\_ptr}\?>\?0\8\&{loop}\6
\&{if}\1\1\8\\{out\_ptr}\?\L\?\\{line\_length}\2\8\&{then}\6
\\{break\_ptr}\?\K\?\\{out\_ptr};\2\6
\&{end}\8\2\&{if}\1;\6
\\{flush\_buffer};\2\6
\&{end}\8\&{loop};\6
\\{out\_state}\?\K\?\\{misc}\par
\U sections~101, 108, and~112(3).\fi

\M114. \P\?\X81:Procedures in \\{output\_phase}\X\mathrel{+}\S\?\6
\&{procedure}\8\1\\{phase\_ii}\2\8\&{is}\8\1\6
\4\&{begin}\6
\X100:Phase II: Output the contents of the compressed tables\X;\2\6
\&{end}\8\\{phase\_ii};\par
\fi

\N115.  Introduction to the input phase.
We have now seen that \.{ATANGLE} will be able to output the full
\ADA\ program, if we can only get that program into the byte memory in
the proper format. The input process is something like the output process
in reverse, since we compress the text as we read it in and we expand it
as we write it out.

There are three main input routines. The most interesting is the one that
gets the next token of an \ADA\ text; the other two are used to scan rapidly
past \TeX\ text in the \.{AWEB} source code. One of the latter routines will
jump to the next token that starts with `\.{@}', and the other skips to the
end of an \ADA\ comment.

\fi

\M116. But first we need to consider the low-level routine \\{get\_line}
that takes care of merging \\{change\_file} into \\{web\_file}. The \\{get%
\_line}
procedure also updates the line numbers for error messages.

\Y\P\?\4\X10:Types and objects visible from \\{data\_structures}\X\mathrel{+}\S%
\?\6
\\{line}\?:\?\p{INTEGER};\C{the number of the current line in the current file}%
\6
\\{other\_line}\?:\?\p{INTEGER};\C{the number of the current line in the input
file that   is not currently being read}\6
\\{temp\_line}\?:\?\p{INTEGER};\C{used when interchanging \\{line} with %
\\{other\_line}}\6
\\{limit}\?:\?\\{buffer\_index};\C{the last character position occupied in the
buffer}\6
\\{loc}\?:\?\\{buffer\_index};\C{the next character position to be read from
the buffer}\6
\\{input\_has\_ended}\?:\?\p{BOOLEAN};\C{if \p{TRUE}, there is no more input}\6
\\{changing}\?:\?\p{BOOLEAN};\C{if \p{TRUE}, the current line is from \\{change%
\_file}}\par
\fi

\M117. As we change \\{changing} from \p{TRUE} to \p{FALSE} and back again, we
must
remember to swap the values of \\{line} and \\{other\_line} so that the
\\{err\_print} routine will be sure to report the correct line number.

\Y\P\4\D\\{change\_changing}\?\S\?\\{changing}\?\K\R\?\\{changing};\5
\\{temp\_line}\?\K\?\\{other\_line};\5
\\{other\_line}\?\K\?\\{line};\5
\\{line}\?\K\?\\{temp\_line}\C{\\{line}\hbox{$\null\BA\null$}\\{other\_line}}%
\par
\fi

\M118. When \\{changing} is \p{FALSE}, the next line of \\{change\_file} is
kept in
\\{change\_buffer}\?(\?0\?\to\?\\{change\_limit}\?)\?, for purposes of
comparison with the next
line of \\{web\_file}. After the change file has been completely input, we
set \\{change\_limit}\?\K\?0, so that no further matches will be made.

\Y\P\?\4\X10:Types and objects visible from \\{data\_structures}\X\mathrel{+}\S%
\?\6
\\{change\_buffer}\?:\?\\{buffer\_type}\?\1(\?\\{buffer\_index}\?\2)\?;\6
\\{change\_limit}\?:\?\\{buffer\_index};\C{the last position occupied in %
\\{change\_buffer}}\par
\fi

\M119. Here's a simple function that checks if the two buffers are different.

\Y\P\?\4\X119:Procedures in \\{input\_phase}\X\S\?\6
\&{function}\8\1\\{lines\_dont\_match}\?\8\&{return}\?\8\p{BOOLEAN}\2\8\&{is}\8%
\1\6
\4\&{begin}\6
\&{if}\1\1\8\\{change\_limit}\?\I\?\\{limit}\2\8\&{then}\6
\?\&{return}\?\8\p{TRUE};\2\6
\&{end}\8\2\&{if}\1;\6
\&{if}\1\1\8\\{limit}\?>\?0\?\ANT\?\\{change\_buffer}\?\1(\?0\?\to\?\\{limit}%
\?-\?1\?\2)\I\?\\{buffer}\?\1(\?0\?\to\?\\{limit}\?-\?1\?\2)\?\2\8\&{then}\6
\?\&{return}\?\8\p{TRUE};\6
\4\&{else}\6
\?\&{return}\?\8\p{FALSE};\2\6
\&{end}\8\2\&{if}\1;\2\6
\&{end}\8\\{lines\_dont\_match};\par
\A sections~120, 124, 127, 132, 133, 134, 137, 149, 155, 157, and~164.
\U section~7\*.\fi

\M120. Procedure \\{prime\_the\_change\_buffer} sets \\{change\_buffer} in
preparation
for the next matching operation. Since blank lines in the change file are
not used for matching, we have \?(\?\\{change\_limit}\?=\?0\?)\W\R\?%
\\{changing} if and
only if the change file is exhausted. This procedure is called only
when \\{changing} is true; hence error messages will be reported correctly.

\Y\P\?\4\X119:Procedures in \\{input\_phase}\X\mathrel{+}\S\?\6
\&{procedure}\8\1\\{prime\_the\_change\_buffer}\2\8\&{is}\8\1\6
\4\&{begin}\6
\\{change\_limit}\?\K\?0;\C{this value will be used if the change file ends}\6
\X121:Skip over comment lines in the change file; \&{return} if end of file\X;\6
\X122:Skip to the next nonblank line; \&{return} if end of file\X;\6
\X123:Move \\{buffer} and \\{limit} to \\{change\_buffer} and \\{change\_limit}%
\X;\2\6
\&{end}\8\\{prime\_the\_change\_buffer};\par
\fi

\M121. While looking for a line that begins with \.{@x} in the change file,
we allow lines that begin with \.{@}, as long as they don't begin with
\.{@y} or \.{@z} (which would probably indicate that the change file is
fouled up).

\Y\P\?\4\X121:Skip over comment lines in the change file; \&{return} if end of
file\X\S\?\6
\1\&{loop}\6
\\{incr}\?\1(\?\\{line}\?\2)\?;\6
\&{if}\1\1\?\8\R\?\\{input\_ln}\?\1(\?\\{change\_file}\?\2)\?\2\8\&{then}\6
\&{return};\2\6
\&{end}\8\2\&{if}\1;\6
\&{if}\1\1\8\\{limit}\?<\?2\2\8\&{then}\6
\&{goto}\8\\{continue};\2\6
\&{end}\8\2\&{if}\1;\6
\&{if}\1\1\8\\{buffer}\?\1(\?0\?\2)\I\?\H{@}\2\8\&{then}\6
\&{goto}\8\\{continue};\2\6
\&{end}\8\2\&{if}\1;\6
\&{if}\1\1\8\\{buffer}\?\1(\?1\?\2)\in\?\H{X}\?\to\?\H{Z}\2\8\&{then}\6
\\{buffer}\?\1(\?1\?\2)\K\?\\{buffer}\?\1(\?1\?\2)+\?\H{z}\?-\?\H{Z};%
\C{lowercasify}\2\6
\&{end}\8\2\&{if}\1;\6
\?\&{exit}\8\&{when}\?\8\\{buffer}\?\1(\?1\?\2)=\?\H{x};\6
\&{if}\1\1\8\\{buffer}\?\1(\?1\?\2)\in\?\H{y}\?\to\?\H{z}\2\8\&{then}\6
\\{loc}\?\K\?2;\5
\\{err\_print}\?\1(\?\.{"!\ Where\ is\ the\ matching\ @x?"}\?\2)\?;\2\6
\&{end}\8\2\&{if}\1;\6
\4\BL\\{continue}\EL\8\&{null};\2\6
\&{end}\8\&{loop}\par
\U section~120.\fi

\M122. Here we are looking at lines following the \.{@x}.

\Y\P\?\4\X122:Skip to the next nonblank line; \&{return} if end of file\X\S\?\6
\1\&{repeat}\6
\\{incr}\?\1(\?\\{line}\?\2)\?;\6
\&{if}\1\1\?\8\R\?\\{input\_ln}\?\1(\?\\{change\_file}\?\2)\?\2\8\&{then}\6
\\{err\_print}\?\1(\?\.{"!\ Change\ file\ ended\ after\ @x"}\?\2)\?;\6
\&{return};\2\6
\&{end}\8\2\&{if}\1;\2\6
\&{until}\?\8\1(\?\\{limit}\?>\?0\?\2)\?\par
\U section~120.\fi

\M123. \P\?\X123:Move \\{buffer} and \\{limit} to \\{change\_buffer} and %
\\{change\_limit}\X\S\?\6
\\{change\_limit}\?\K\?\\{limit};\6
\&{if}\1\1\8\\{limit}\?>\?0\2\8\&{then}\6
\\{change\_buffer}\?\1(\?0\?\to\?\\{limit}\?-\?1\?\2)\K\?\\{buffer}\?\1(\?0\?%
\to\?\\{limit}\?-\?1\?\2)\?;\2\6
\&{end}\8\2\&{if}\1\par
\U sections~120 and~124.\fi

\M124. The following procedure is used to see if the next change entry should
go into effect; it is called only when \\{changing} is false.
The idea is to test whether or not the current
contents of \\{buffer} matches the current contents of \\{change\_buffer}.
If not, there's nothing more to do; but if so, a change is called for:
All of the text down to the \.{@y} is supposed to match. An error
message is issued if any discrepancy is found. Then the procedure
prepares to read the next line from \\{change\_file}.

\Y\P\?\4\X119:Procedures in \\{input\_phase}\X\mathrel{+}\S\?\6
\&{procedure}\8\1\\{check\_change}\2\8\&{is}\8\C{switches to \\{change\_file}
if the buffers match}\1\6
\|n\?:\?\p{INTEGER}\?\K\?0;\C{the number of discrepancies found}\6
\4\&{begin}\6
\&{if}\1\1\8\\{lines\_dont\_match}\2\8\&{then}\6
\&{return};\2\6
\&{end}\8\2\&{if}\1;\6
\1\&{loop}\6
\\{change\_changing};\C{now it's \p{TRUE}}\6
\\{incr}\?\1(\?\\{line}\?\2)\?;\6
\&{if}\1\1\?\8\R\?\\{input\_ln}\?\1(\?\\{change\_file}\?\2)\?\2\8\&{then}\6
\\{err\_print}\?\1(\?\.{"!\ Change\ file\ ended\ before\ @y"}\?\2)\?;\6
\\{change\_limit}\?\K\?0;\5
\\{change\_changing};\C{\p{FALSE} again}\6
\&{return};\2\6
\&{end}\8\2\&{if}\1;\6
\X125:If the current line starts with \.{@y}, report any discrepancies and %
\&{return}\X;\6
\X123:Move \\{buffer} and \\{limit} to \\{change\_buffer} and \\{change\_limit}%
\X;\6
\\{change\_changing};\C{now it's \p{FALSE}}\6
\\{incr}\?\1(\?\\{line}\?\2)\?;\6
\&{if}\1\1\?\8\R\?\\{input\_ln}\?\1(\?\\{web\_file}\?\2)\?\2\8\&{then}\6
\\{err\_print}\?\1(\?\.{"!\ AWEB\ file\ ended\ during\ a\ change"}\?\2)\?;\6
\\{input\_has\_ended}\?\K\?\p{TRUE};\5
\&{return};\2\6
\&{end}\8\2\&{if}\1;\6
\&{if}\1\1\8\\{lines\_dont\_match}\2\8\&{then}\6
\\{incr}\?\1(\?\|n\?\2)\?;\2\6
\&{end}\8\2\&{if}\1;\2\6
\&{end}\8\&{loop};\2\6
\&{end}\8\\{check\_change};\par
\fi

\M125. \P\?\X125:If the current line starts with \.{@y}, report any
discrepancies and \&{return}\X\S\?\6
\&{if}\1\1\?\8\1(\?\\{limit}\?>\?1\?\2)\ANT\?\\{buffer}\?\1(\?0\?\2)=\?\H{@}\2%
\8\&{then}\6
\&{if}\1\1\8\\{buffer}\?\1(\?1\?\2)\in\?\H{X}\?\to\?\H{Z}\2\8\&{then}\6
\\{buffer}\?\1(\?1\?\2)\K\?\\{buffer}\?\1(\?1\?\2)+\?\H{z}\?-\?\H{Z};%
\C{lowercasify}\2\6
\&{end}\8\2\&{if}\1;\6
\&{if}\1\1\?\8\1(\?\\{buffer}\?\1(\?1\?\2)=\?\H{x}\?\2)\V\1(\?\\{buffer}\?\1(%
\?1\?\2)=\?\H{z}\?\2)\?\2\8\&{then}\6
\\{loc}\?\K\?2;\5
\\{err\_print}\?\1(\?\.{"!\ Where\ is\ the\ matching\ @y?"}\?\2)\?;\6
\4\&{elsif}\1\8\\{buffer}\?\1(\?1\?\2)=\?\H{y}\2\8\&{then}\6
\&{if}\1\1\8\|n\?>\?0\2\8\&{then}\6
\\{loc}\?\K\?2;\5
\\{new\_ln};\5
\\{print}\?\1(\?\.{"!\ Hmm...\ "}\?\2)\?;\5
\\{print\_int}\?\1(\?\|n\?,\39\?1\?\2)\?;\5
\\{print}\?\1(\?\.{"\ of\ the\ preceding\ lines\ failed\ to\ match"}\?\2)\?;\5
\\{error};\2\6
\&{end}\8\2\&{if}\1;\6
\&{return};\2\6
\&{end}\8\2\&{if}\1;\2\6
\&{end}\8\2\&{if}\1\par
\U section~124.\fi

\M126. \P\?\X126:Initialize the input system\X\S\?\6
\\{line}\?\K\?0;\5
\\{other\_line}\?\K\?0;\6
\\{changing}\?\K\?\p{TRUE};\5
\\{prime\_the\_change\_buffer};\5
\\{change\_changing};\6
\\{limit}\?\K\?0;\5
\\{loc}\?\K\?1;\5
\\{buffer}\?\1(\?0\?\2)\K\?\H{\ };\5
\\{input\_has\_ended}\?\K\?\p{FALSE}\par
\U section~164.\fi

\M127. The \\{get\_line} procedure is called when \\{loc}\?>\?\\{limit}; it
puts the next
line of merged input into the buffer and updates the other variables
appropriately. A space is placed at the right end of the line.

\Y\P\?\4\X119:Procedures in \\{input\_phase}\X\mathrel{+}\S\?\6
\&{procedure}\8\1\\{get\_line}\2\8\&{is}\8\C{inputs the next line}\1\6
\4\&{begin}\6
\4\BL\\{restart}\EL\8\&{if}\1\1\8\\{changing}\2\8\&{then}\6
\X129:Read from \\{change\_file} and maybe turn off \\{changing}\X;\2\6
\&{end}\8\2\&{if}\1;\6
\&{if}\1\1\?\8\R\?\\{changing}\2\8\&{then}\6
\X128:Read from \\{web\_file} and maybe turn on \\{changing}\X;\6
\&{if}\1\1\8\\{changing}\2\8\&{then}\6
\&{goto}\8\\{restart};\2\6
\&{end}\8\2\&{if}\1;\2\6
\&{end}\8\2\&{if}\1;\6
\\{loc}\?\K\?0;\5
\\{buffer}\?\1(\?\\{limit}\?\2)\K\?\H{\ };\2\6
\&{end}\8\\{get\_line};\par
\fi

\M128. \P\?\X128:Read from \\{web\_file} and maybe turn on \\{changing}\X\S\?\6
\\{incr}\?\1(\?\\{line}\?\2)\?;\6
\&{if}\1\1\?\8\R\?\\{input\_ln}\?\1(\?\\{web\_file}\?\2)\?\2\8\&{then}\6
\\{input\_has\_ended}\?\K\?\p{TRUE};\6
\4\&{elsif}\1\?\8\1(\?\\{limit}\?=\?\\{change\_limit}\?\2)\ANT\1(\1(\?%
\\{buffer}\?\1(\?0\?\2)=\?\\{change\_buffer}\?\1(\?0\?\2)\2)\ANT\1(\?\\{change%
\_limit}\?>\?0\?\2)\2)\?\2\8\&{then}\6
\\{check\_change};\2\6
\&{end}\8\2\&{if}\1\par
\U section~127.\fi

\M129. \P\?\X129:Read from \\{change\_file} and maybe turn off \\{changing}\X\S%
\?\6
\\{incr}\?\1(\?\\{line}\?\2)\?;\6
\&{if}\1\1\?\8\R\?\\{input\_ln}\?\1(\?\\{change\_file}\?\2)\?\2\8\&{then}\6
\\{err\_print}\?\1(\?\.{"!\ Change\ file\ ended\ without\ @z"}\?\2)\?;\6
\\{buffer}\?\1(\?0\?\to\?1\?\2)\K\1(\?\H{@}\?,\39\?\H{z}\?\2)\?;\5
\\{limit}\?\K\?2;\2\6
\&{end}\8\2\&{if}\1;\6
\&{if}\1\1\?\8\1(\?\\{limit}\?>\?1\?\2)\ANT\1(\?\\{buffer}\?\1(\?0\?\2)=\?\H{@}%
\?\2)\?\2\8\&{then}\C{check if the change has ended}\6
\&{if}\1\1\8\\{buffer}\?\1(\?1\?\2)\in\?\H{X}\?\to\?\H{Z}\2\8\&{then}\6
\\{buffer}\?\1(\?1\?\2)\K\?\\{buffer}\?\1(\?1\?\2)+\?\H{z}\?-\?\H{Z};%
\C{lowercasify}\2\6
\&{end}\8\2\&{if}\1;\6
\&{if}\1\1\8\\{buffer}\?\1(\?1\?\2)\in\?\H{x}\?\to\?\H{y}\2\8\&{then}\6
\\{loc}\?\K\?2;\5
\\{err\_print}\?\1(\?\.{"!\ Where\ is\ the\ matching\ @z?"}\?\2)\?;\6
\4\&{elsif}\1\8\\{buffer}\?\1(\?1\?\2)=\?\H{z}\2\8\&{then}\6
\\{prime\_the\_change\_buffer};\5
\\{change\_changing};\2\6
\&{end}\8\2\&{if}\1;\2\6
\&{end}\8\2\&{if}\1\par
\U section~127.\fi

\M130. At the end of the program, we will tell the user if the change file
had a line that didn't match any relevant line in \\{web\_file}.

\Y\P\?\4\X130:Check that all changes have been read\X\S\?\6
\&{if}\1\1\8\\{change\_limit}\?\I\?0\2\8\&{then}\C{\\{changing} is false}\6
\\{buffer}\?\1(\?0\?\to\?\\{change\_limit}\?\2)\K\?\\{change\_buffer}\?\1(\?0\?%
\to\?\\{change\_limit}\?\2)\?;\5
\\{limit}\?\K\?\\{change\_limit};\5
\\{changing}\?\K\?\p{TRUE};\5
\\{line}\?\K\?\\{other\_line};\5
\\{loc}\?\K\?\\{change\_limit};\5
\\{err\_print}\?\1(\?\.{"!\ Change\ file\ entry\ did\ not\ match"}\?\2)\?;\2\6
\&{end}\8\2\&{if}\1\par
\U section~167.\fi

\M131. Important milestones are reached during the input phase when certain
control codes are sensed.

Control codes in \.{AWEB} begin with `\.{@}', and the next character
identifies the code. Some of these are of interest only to \.{AWEAVE},
so \.{ATANGLE} ignores them; the others are converted by \.{ATANGLE} into
internal code numbers by the \\{control\_code} function below. The ordering
of these internal code numbers has been chosen to simplify the program logic;
larger numbers are given to the control codes that denote more significant
milestones.

\Y\P\?\4\X10:Types and objects visible from \\{data\_structures}\X\mathrel{+}\S%
\?\6
\\{ignore}\?:\?\&{constant}\8\\{eight\_bits}\?\K\?0;\C{control code of no
interest to \.{ATANGLE}}\6
\\{control\_text}\?:\?\&{constant}\8\\{eight\_bits}\?\K\bn{8}{203}\?;\C{control
code for `\.{@t}', `\.{@\^}', etc.}\6
\\{format}\?:\?\&{constant}\8\\{eight\_bits}\?\K\bn{8}{204}\?;\C{control code
for `\.{@f}'}\6
\\{definition}\?:\?\&{constant}\8\\{eight\_bits}\?\K\bn{8}{205}\?;\C{control
code for `\.{@d}'}\6
\\{begin\_ada}\?:\?\&{constant}\8\\{eight\_bits}\?\K\bn{8}{206}\?;\C{control
code for `\.{@a}'}\6
\\{module\_name}\?:\?\&{constant}\8\\{eight\_bits}\?\K\bn{8}{207}\?;\C{control
code for `\.{@<}'}\6
\\{new\_module}\?:\?\&{constant}\8\\{eight\_bits}\?\K\bn{8}{210}\?;\C{control
code for `\.{@\ }' and `\.{@*}'}\par
\fi

\M132. \P\?\X119:Procedures in \\{input\_phase}\X\mathrel{+}\S\?\6
\&{function}\8\1\\{control\_code}\?(\1\?\|c\?:\?\\{ASCII\_code}\?\2)\&{return}%
\?\8\\{eight\_bits}\2\8\&{is}\8\C{convert \|c after \.{@}}\1\6
\4\&{begin}\6
\&{case}\1\8\|c\1\8\&{is}\8\6
\4\&{when}\8\H{@}\KR\8\?\&{return}\?\8\H{@};\C{`quoted' at sign}\6
\4\&{when}\8\H{\'}\KR\8\?\&{return}\?\8\\{ascii\_tok};\C{precedes ASCII code
constant}\6
\4\&{when}\8\H{\~}\KR\8\?\&{return}\?\8\\{output\_tok};\C{precedes output
file-name}\6
\4\&{when}\8\H{\ }\?\mathrel{\.|}\?\\{tab\_mark}\KR\8\?\&{return}\?\8\\{new%
\_module};\C{beginning of a new module}\6
\4\&{when}\8\H{*}\KR\8\\{print}\?\1(\?\.{"*"}\?\2)\?;\5
\\{print\_int}\?\1(\?\\{module\_count}\?+\?1\?,\39\?1\?\2)\?;\5
\\{update\_terminal};\C{print a progress report}\6
\?\&{return}\?\8\\{new\_module};\C{beginning of a new module}\6
\4\&{when}\8\H{D}\?\mathrel{\.|}\?\H{d}\KR\8\?\&{return}\?\8\\{definition};%
\C{macro definition}\6
\4\&{when}\8\H{F}\?\mathrel{\.|}\?\H{f}\KR\8\?\&{return}\?\8\\{format};%
\C{format definition}\6
\4\&{when}\8\H{\{}\KR\8\?\&{return}\?\8\\{begin\_comment};\C{begin-comment
delimiter}\6
\4\&{when}\8\H{\}}\KR\8\?\&{return}\?\8\\{end\_comment};\C{end-comment
delimiter}\6
\4\&{when}\8\H{A}\?\mathrel{\.|}\?\H{a}\KR\8\?\&{return}\?\8\\{begin\_ada};\C{%
\ADA\ text in unnamed module}\6
\4\&{when}\8\H{T}\?\mathrel{\.|}\?\H{t}\?\mathrel{\.|}\?\H{\^}\?\mathrel{\.|}\?%
\H{i}\?\mathrel{\.|}\?\H{I}\?\mathrel{\.|}\?\H{b}\?\mathrel{\.|}\?\H{B}\?%
\mathrel{\.|}\?\H{r}\?\mathrel{\.|}\?\H{R}\?\mathrel{\.|}\?\H{p}\?\mathrel{\.|}%
\?\H{P}\?\mathrel{\.|}\?\H{s}\?\mathrel{\.|}\?\H{S}\?\mathrel{\.|}\?\H{.}\?%
\mathrel{\.|}\?\H{:}\KR\8\?\&{return}\?\8\\{control\_text};\C{control text to
be ignored}\6
\4\&{when}\8\H{\&}\KR\8\?\&{return}\?\8\\{join};\C{concatenate two tokens}\6
\4\&{when}\8\H{<}\KR\8\?\&{return}\?\8\\{module\_name};\C{beginning of a module
name}\6
\4\&{when}\8\H{=}\KR\8\?\&{return}\?\8\\{verbatim};\C{beginning of \ADA\
verbatim mode}\6
\4\&{when}\8\H{\\}\KR\8\?\&{return}\?\8\\{force\_line};\C{force a new line in %
\ADA\ output}\6
\4\&{when}\8\&{others}\8\KR\8\?\&{return}\?\8\\{ignore};\C{ignore all other
cases}\2\6
\4\2\&{end}\8\&{case};\2\6
\&{end}\8\\{control\_code};\par
\fi

\M133. The \\{skip\_ahead} procedure reads through the input at fairly high
speed
until finding the next non-ignorable control code, which it returns.

\Y\P\?\4\X119:Procedures in \\{input\_phase}\X\mathrel{+}\S\?\6
\&{function}\8\1\\{skip\_ahead}\?\8\&{return}\?\8\\{eight\_bits}\2\8\&{is}\8%
\C{skip to next control code}\1\6
\|c\?:\?\\{eight\_bits};\C{control code found}\6
\4\&{begin}\6
\1\&{loop}\6
\&{if}\1\1\8\\{loc}\?>\?\\{limit}\2\8\&{then}\6
\\{get\_line};\6
\&{if}\1\1\8\\{input\_has\_ended}\2\8\&{then}\6
\?\&{return}\?\8\\{new\_module};\2\6
\&{end}\8\2\&{if}\1;\2\6
\&{end}\8\2\&{if}\1;\6
\\{buffer}\?\1(\?\\{limit}\?+\?1\?\2)\K\?\H{@};\6
\1\&{while}\8\\{buffer}\?\1(\?\\{loc}\?\2)\I\?\H{@}\8\&{loop}\6
\\{incr}\?\1(\?\\{loc}\?\2)\?;\2\6
\&{end}\8\&{loop};\6
\&{if}\1\1\8\\{loc}\?\L\?\\{limit}\2\8\&{then}\6
\\{loc}\?\K\?\\{loc}\?+\?2;\5
\|c\?\K\?\\{control\_code}\?\1(\?\\{buffer}\?\1(\?\\{loc}\?-\?1\?\2)\2)\?;\6
\&{if}\1\1\?\8\1(\?\|c\?\I\?\\{ignore}\?\2)\V\1(\?\\{buffer}\?\1(\?\\{loc}\?-%
\?1\?\2)=\?\H{>}\?\2)\?\2\8\&{then}\6
\?\&{return}\?\8\|c;\2\6
\&{end}\8\2\&{if}\1;\2\6
\&{end}\8\2\&{if}\1;\2\6
\&{end}\8\&{loop};\2\6
\&{end}\8\\{skip\_ahead};\par
\fi

\M134. The \\{skip\_comment} procedure reads through the input at somewhat high
speed
until finding the first unmatched right brace or until coming to the end
of the file. It ignores characters following `\.\\' characters, since all
braces that aren't nested are supposed to be hidden in that way. For
example, consider the process of skipping the first comment below,
where the string containing the right brace has been typed as \.{`\\.\\\}'}
in the \.{AWEB} file.

\Y\P\?\4\X119:Procedures in \\{input\_phase}\X\mathrel{+}\S\?\6
\&{procedure}\8\1\\{skip\_comment}\2\8\&{is}\8\C{skips to next unmatched `\.%
\}'}\1\6
\\{bal}\?:\?\\{eight\_bits}\?\K\?0;\C{excess of left braces}\6
\|c\?:\?\\{ASCII\_code};\C{current character}\6
\4\&{begin}\6
\1\&{loop}\6
\&{if}\1\1\8\\{loc}\?>\?\\{limit}\2\8\&{then}\6
\\{get\_line};\6
\&{if}\1\1\8\\{input\_has\_ended}\2\8\&{then}\6
\\{err\_print}\?\1(\?\.{"!\ Input\ ended\ in\ mid-comment"}\?\2)\?;\6
\&{return};\2\6
\&{end}\8\2\&{if}\1;\2\6
\&{end}\8\2\&{if}\1;\6
\|c\?\K\?\\{buffer}\?\1(\?\\{loc}\?\2)\?;\5
\\{incr}\?\1(\?\\{loc}\?\2)\?;\6
\X135:Do special things when \|c\?=\?\.{"@"}\?,\?\.{"\\"}\?,\?\.{"\{"}\?,\?\.{"%
\}"}; \&{return} at end\X;\2\6
\&{end}\8\&{loop};\2\6
\&{end}\8\\{skip\_comment};\par
\fi

\M135. \P\?\X135:Do special things when \|c\?=\?\.{"@"}\?,\?\.{"\\"}\?,\?\.{"%
\{"}\?,\?\.{"\}"}; \&{return} at end\X\S\?\6
\&{if}\1\1\8\|c\?=\?\H{@}\2\8\&{then}\6
\|c\?\K\?\\{buffer}\?\1(\?\\{loc}\?\2)\?;\6
\&{if}\1\1\?\8\1(\?\|c\?\I\?\H{\ }\?\2)\W\1(\?\|c\?\I\?\\{tab\_mark}\?\2)\W\1(%
\?\|c\?\I\?\H{*}\?\2)\W\1(\?\|c\?\I\?\H{z}\?\2)\W\1(\?\|c\?\I\?\H{Z}\?\2)\?\2\8%
\&{then}\6
\\{incr}\?\1(\?\\{loc}\?\2)\?;\6
\4\&{else}\6
\\{err\_print}\?\1(\?\.{"!\ Section\ ended\ in\ mid-comment"}\?\2)\?;\6
\\{decr}\?\1(\?\\{loc}\?\2)\?;\5
\&{return};\2\6
\&{end}\8\2\&{if}\1;\6
\4\&{elsif}\1\?\8\1(\?\|c\?=\?\H{\\}\?\2)\W\1(\?\\{buffer}\?\1(\?\\{loc}\?\2)\I%
\?\H{@}\?\2)\?\2\8\&{then}\6
\\{incr}\?\1(\?\\{loc}\?\2)\?;\6
\4\&{elsif}\1\8\|c\?=\?\H{\{}\2\8\&{then}\6
\\{incr}\?\1(\?\\{bal}\?\2)\?;\6
\4\&{elsif}\1\8\|c\?=\?\H{\}}\2\8\&{then}\6
\&{if}\1\1\8\\{bal}\?=\?0\2\8\&{then}\6
\&{return};\2\6
\&{end}\8\2\&{if}\1;\6
\\{decr}\?\1(\?\\{bal}\?\2)\?;\2\6
\&{end}\8\2\&{if}\1\par
\U section~134.\fi

\N136.  Inputting the next token.
As stated above, \.{ATANGLE}'s most interesting input procedure is the
\\{get\_next} routine that inputs the next token. However, the procedure
isn't especially difficult.

In most cases the tokens output by \\{get\_next} have the form used in
replacement texts, except that two-byte tokens are not produced.
An identifier that isn't one letter long is represented by the
output `\\{identifier}', and in such a case the global variables
\\{id\_first} and \\{id\_loc} will have been set to the appropriate values
needed by the \\{id\_lookup} procedure.
Control codes produce the corresponding output of the \\{control\_code}
function above; and if that code is \\{module\_name}, the value of \\{cur%
\_module}
will point to the \\{byte\_start} entry for that module name.

Another global variable, \\{scanning\_based\_number}, is \p{TRUE} during the
time
that letters should be treated as if they were digits; \\{scanning\_character}
is set \p{TRUE} if the input \.'$x$\.', where $x$ is a character, or the
ASCII code controlcode `\.{@'} is encoutered.

\Y\P\?\4\X10:Types and objects visible from \\{data\_structures}\X\mathrel{+}\S%
\?\6
\\{cur\_module}\?:\?\\{name\_pointer};\C{name of module just scanned}\6
\\{scanning\_character}\?:\?\p{BOOLEAN}\?\K\?\p{FALSE};\C{are we scanning a
character literal?}\par
\fi

\M137. At the top level, \\{get\_next} is a multi-way switch based on the next
character in the input buffer. A \\{new\_module} code is inserted at the
very end of the input file.

\Y\P\?\4\X119:Procedures in \\{input\_phase}\X\mathrel{+}\S\?\6
\&{function}\8\1\\{get\_next}\?\8\&{return}\?\8\\{eight\_bits}\2\8\&{is}\8%
\C{produces the next input token}\1\6
\|c\?:\?\\{eight\_bits};\C{the current character}\6
\|d\?:\?\\{eight\_bits};\C{the next character}\6
\|k\?:\?\\{long\_name\_index};\C{index into \\{mod\_text}}\6
\4\&{begin}\6
\4\BL\\{restart}\EL\8\&{if}\1\1\8\\{loc}\?>\?\\{limit}\2\8\&{then}\6
\\{get\_line};\6
\&{if}\1\1\8\\{input\_has\_ended}\2\8\&{then}\6
\?\&{return}\?\8\\{new\_module};\6
\4\&{elsif}\1\8\\{scanning\_based\_number}\2\8\&{then}\6
\\{err\_print}\?\1(\?\.{"!\ Improper\ based\ number,\ \#\ sign\ missing"}\?\2)%
\?;\6
\\{scanning\_based\_number}\?\K\?\p{FALSE};\2\6
\&{end}\8\2\&{if}\1;\2\6
\&{end}\8\2\&{if}\1;\6
\|c\?\K\?\\{buffer}\?\1(\?\\{loc}\?\2)\?;\5
\\{incr}\?\1(\?\\{loc}\?\2)\?;\6
\&{if}\1\1\8\\{scanning\_based\_number}\2\8\&{then}\6
\&{if}\1\1\8\|c\?=\?\H{\#}\2\8\&{then}\6
\\{scanning\_based\_number}\?\K\?\p{FALSE};\2\6
\&{end}\8\2\&{if}\1;\6
\?\&{return}\?\8\|c;\6
\4\&{elsif}\1\8\\{scanning\_character}\2\8\&{then}\6
\&{if}\1\1\8\|c\?=\?\H{@}\2\8\&{then}\6
\&{if}\1\1\8\\{buffer}\?\1(\?\\{loc}\?\2)=\?\H{@}\2\8\&{then}\6
\\{incr}\?\1(\?\\{loc}\?\2)\?;\6
\4\&{else}\6
\\{err\_print}\?\1(\?\.{"!\ You\ should\ double\ @\ signs\ in\ character\
literals"}\?\2)\?;\2\6
\&{end}\8\2\&{if}\1;\2\6
\&{end}\8\2\&{if}\1;\6
\\{incr}\?\1(\?\\{loc}\?\2)\?;\5
\\{scanning\_character}\?\K\?\p{FALSE};\5
\?\&{return}\?\8\|c;\2\6
\&{end}\8\2\&{if}\1;\6
\&{case}\1\8\|c\1\8\&{is}\8\6
\4\&{when}\8\H{\#}\KR\6
\&{if}\1\1\?\8\1(\?\\{loc}\?>\?1\?\2)\ANT\1(\?\\{buffer}\?\1(\?\\{loc}\?-\?2\?%
\2)\in\?\H{0}\?\to\?\H{9}\?\2)\?\2\8\&{then}\6
\\{scanning\_based\_number}\?\K\?\p{TRUE};\2\6
\&{end}\8\2\&{if}\1;\6
\?\&{return}\?\8\H{\#};\6
\4\&{when}\8\H{A}\?\to\?\H{Z}\?\mathrel{\.|}\?\H{a}\?\to\?\H{z}\KR\8\X139:Get
an identifier\X;\6
\4\&{when}\8\H{@}\KR\8\X140:Get control code and possible module name\X;\6
\hbox{\4}\X138:Compress two-symbol combinations like `\.{:=}'\X\6
\4\&{when}\8\H{\ }\?\mathrel{\.|}\?\\{tab\_mark}\KR\8\&{goto}\8\\{restart};%
\C{ignore spaces and tabs}\6
\4\&{when}\8\H{\{}\KR\8\\{skip\_comment};\5
\&{goto}\8\\{restart};\6
\4\&{when}\8\H{\'}\KR\8\X146:If this is the beginning of an character token
then set \\{scanning\_character} to \p{TRUE} and \?\&{return}\?\8\\{character%
\_tok}; otherwise do nothing\X;\6
\4\&{when}\8\&{others}\8\KR\8\&{null};\2\6
\4\2\&{end}\8\&{case};\6
\?\&{return}\?\8\|c;\2\6
\&{end}\8\\{get\_next};\par
\fi

\M138.
\Y\P\4\D\\{compress}\?\1(\?\#\?\2)\S\?\6
\&{if}\1\1\8\\{loc}\?\L\?\\{limit}\2\8\&{then}\6
\|c\?\K\?\#;\5
\\{incr}\?\1(\?\\{loc}\?\2)\?;\2\6
\&{end}\8\2\&{if}\1\par
\Y\P\?\4\X138:Compress two-symbol combinations like `\.{:=}'\X\S\?\6
\4\&{when}\8\H{.}\KR\6
\&{if}\1\1\8\\{buffer}\?\1(\?\\{loc}\?\2)=\?\H{.}\2\8\&{then}\6
\\{compress}\?\1(\?\\{double\_dot}\?\2)\?;\2\6
\&{end}\8\2\&{if}\1;\6
\4\&{when}\8\H{:}\KR\6
\&{if}\1\1\8\\{buffer}\?\1(\?\\{loc}\?\2)=\?\H{=}\2\8\&{then}\6
\\{compress}\?\1(\?\\{left\_arrow}\?\2)\?;\2\6
\&{end}\8\2\&{if}\1;\6
\4\&{when}\8\H{=}\KR\6
\&{if}\1\1\8\\{buffer}\?\1(\?\\{loc}\?\2)=\?\H{=}\2\8\&{then}\6
\\{compress}\?\1(\?\\{equivalence\_sign}\?\2)\?;\6
\4\&{elsif}\1\8\\{buffer}\?\1(\?\\{loc}\?\2)=\?\H{>}\2\8\&{then}\6
\\{compress}\?\1(\?\\{right\_arrow}\?\2)\?;\2\6
\&{end}\8\2\&{if}\1;\6
\4\&{when}\8\H{>}\KR\6
\&{if}\1\1\8\\{buffer}\?\1(\?\\{loc}\?\2)=\?\H{=}\2\8\&{then}\6
\\{compress}\?\1(\?\\{greater\_or\_equal}\?\2)\?;\6
\4\&{elsif}\1\8\\{buffer}\?\1(\?\\{loc}\?\2)=\?\H{>}\2\8\&{then}\6
\\{compress}\?\1(\?\\{right\_label\_bracket}\?\2)\?;\2\6
\&{end}\8\2\&{if}\1;\6
\4\&{when}\8\H{<}\KR\6
\&{if}\1\1\8\\{buffer}\?\1(\?\\{loc}\?\2)=\?\H{=}\2\8\&{then}\6
\\{compress}\?\1(\?\\{less\_or\_equal}\?\2)\?;\6
\4\&{elsif}\1\8\\{buffer}\?\1(\?\\{loc}\?\2)=\?\H{<}\2\8\&{then}\6
\\{compress}\?\1(\?\\{left\_label\_bracket}\?\2)\?;\6
\4\&{elsif}\1\8\\{buffer}\?\1(\?\\{loc}\?\2)=\?\H{>}\2\8\&{then}\6
\\{compress}\?\1(\?\\{box\_sign}\?\2)\?;\2\6
\&{end}\8\2\&{if}\1;\6
\4\&{when}\8\H{/}\KR\6
\&{if}\1\1\8\\{buffer}\?\1(\?\\{loc}\?\2)=\?\H{=}\2\8\&{then}\6
\\{compress}\?\1(\?\\{not\_equal}\?\2)\?;\2\6
\&{end}\8\2\&{if}\1;\6
\4\&{when}\8\H{*}\KR\6
\&{if}\1\1\8\\{buffer}\?\1(\?\\{loc}\?\2)=\?\H{*}\2\8\&{then}\6
\\{compress}\?\1(\?\\{double\_star}\?\2)\?;\2\6
\&{end}\8\2\&{if}\1;\par
\U section~137.\fi

\M139. We have to look at the preceding character to make sure this isn't part
of a real constant, before trying to find an identifier starting with
`\.e' or `\.E'.

\Y\P\?\4\X139:Get an identifier\X\S\?\6
\&{if}\1\1\?\8\1(\1(\1(\?\|c\?=\?\H{e}\?\2)\V\1(\?\|c\?=\?\H{E}\?\2)\2)\W\1(\?%
\\{loc}\?>\?1\?\2)\2)\ANT\1(\1(\?\\{buffer}\?\1(\?\\{loc}\?-\?2\?\2)\in\?\H{0}%
\?\to\?\H{9}\?\2)\V\?\\{buffer}\?\1(\?\\{loc}\?-\?2\?\2)=\?\H{\#}\?\2)\?\2\8%
\&{then}\6
\|c\?\K\?0;\2\6
\&{end}\8\2\&{if}\1;\6
\&{if}\1\1\8\|c\?\I\?0\2\8\&{then}\6
\\{decr}\?\1(\?\\{loc}\?\2)\?;\5
\\{id\_first}\?\K\?\\{loc};\6
\1\&{repeat}\6
\\{incr}\?\1(\?\\{loc}\?\2)\?;\5
\|d\?\K\?\\{buffer}\?\1(\?\\{loc}\?\2)\?;\2\6
\&{until}\?\8\1(\1(\1(\?\|d\?<\?\H{0}\?\2)\V\1(\1(\?\|d\?>\?\H{9}\?\2)\W\1(\?%
\|d\?<\?\H{A}\?\2)\2)\V\1(\1(\?\|d\?>\?\H{Z}\?\2)\W\1(\?\|d\?<\?\H{a}\?\2)\2)\V%
\1(\?\|d\?>\?\H{z}\?\2)\2)\W\1(\?\|d\?\I\?\H{\_}\?\2)\2)\?;\6
\&{if}\1\1\8\\{loc}\?>\?\\{id\_first}\?+\?1\2\8\&{then}\6
\\{id\_loc}\?\K\?\\{loc};\5
\?\&{return}\?\8\\{identifier};\2\6
\&{end}\8\2\&{if}\1;\6
\4\&{else}\6
\?\&{return}\?\8\H{E};\C{exponent of a real constant}\2\6
\&{end}\8\2\&{if}\1\par
\U section~137.\fi

\M140. After an \.{@} sign has been scanned, the next character tells us
whether there is more work to do.

\Y\P\?\4\X140:Get control code and possible module name\X\S\?\6
\|c\?\K\?\\{control\_code}\?\1(\?\\{buffer}\?\1(\?\\{loc}\?\2)\2)\?;\5
\\{incr}\?\1(\?\\{loc}\?\2)\?;\6
\&{if}\1\1\8\|c\?=\?\\{ignore}\2\8\&{then}\6
\&{goto}\8\\{restart};\6
\4\&{elsif}\1\8\|c\?=\?\\{module\_name}\2\8\&{then}\6
\X141:Scan the module name and make \\{cur\_module} point to it\X;\6
\4\&{elsif}\1\8\|c\?=\?\\{control\_text}\2\8\&{then}\6
\1\&{repeat}\6
\|c\?\K\?\\{skip\_ahead};\2\6
\&{until}\?\8\1(\?\|c\?\I\?\H{@}\?\2)\?;\6
\&{if}\1\1\8\\{buffer}\?\1(\?\\{loc}\?-\?1\?\2)\I\?\H{>}\2\8\&{then}\6
\\{err\_print}\?\1(\?\.{"!\ Improper\ @\ within\ control\ text"}\?\2)\?;\2\6
\&{end}\8\2\&{if}\1;\6
\&{goto}\8\\{restart};\6
\4\&{elsif}\1\8\|c\?=\?\\{ascii\_tok}\2\8\&{then}\6
\&{if}\1\1\?\8\1(\?\\{buffer}\?\1(\?\\{loc}\?+\?1\?\2)=\?\H{\'}\?\2)\ORE\1(\?%
\\{buffer}\?\1(\?\\{loc}\?\to\?\\{loc}\?+\?2\?\2)=\1(\?\H{@}\?,\39\?\H{@}\?,\39%
\?\H{\'}\?\2)\2)\?\2\8\&{then}\6
\\{scanning\_character}\?\K\?\p{TRUE};\6
\4\&{else}\6
\\{err\_print}\?\1(\?\.{"!\ Improper\ ASCII-code\ sequence,\ \'\ sign\ mising"}%
\?\2)\?;\2\6
\&{end}\8\2\&{if}\1;\2\6
\&{end}\8\2\&{if}\1\par
\U section~137.\fi

\M141. \P\?\X141:Scan the module name and make \\{cur\_module} point to it\X\S%
\?\6
\X143:Put module name into \\{mod\_text}\?(\?1\?\to\?\|k\?)\?\X;\6
\&{if}\1\1\?\8\1(\?\|k\?>\?3\?\2)\ANT\1(\?\\{mod\_text}\?\1(\?\|k\?-\?2\?\to\?%
\|k\?\2)=\1(\?\H{.}\?,\39\?\H{.}\?,\39\?\H{.}\?\2)\2)\?\2\8\&{then}\6
\\{cur\_module}\?\K\?\\{prefix\_lookup}\?\1(\?\|k\?-\?3\?\2)\?;\6
\4\&{else}\6
\\{cur\_module}\?\K\?\\{mod\_lookup}\?\1(\?\|k\?\2)\?;\2\6
\&{end}\8\2\&{if}\1\par
\U section~140.\fi

\M142. Module names are placed into the \\{mod\_text} array with consecutive
spaces,
tabs, and carriage-returns replaced by single spaces. There will be no
spaces at the beginning or the end. (We set \\{mod\_text}\?(\?0\?)\K\?\H{\ } to
facilitate
this, since the \\{mod\_lookup} routine uses \\{mod\_text}\?(\?1\?)\? as the
first
character of the name.)

\Y\P\?\4\X15:Set initial values\X\mathrel{+}\S\?\6
\\{mod\_text}\?\1(\?0\?\2)\K\?\H{\ };\par
\fi

\M143. \P\?\X143:Put module name into \\{mod\_text}\?(\?1\?\to\?\|k\?)\?\X\S\?\6
\|k\?\K\?0;\6
\1\&{loop}\6
\&{if}\1\1\8\\{loc}\?>\?\\{limit}\2\8\&{then}\6
\\{get\_line};\6
\&{if}\1\1\8\\{input\_has\_ended}\2\8\&{then}\6
\\{err\_print}\?\1(\?\.{"!\ Input\ ended\ in\ section\ name"}\?\2)\?;\6
\&{exit};\2\6
\&{end}\8\2\&{if}\1;\2\6
\&{end}\8\2\&{if}\1;\6
\|d\?\K\?\\{buffer}\?\1(\?\\{loc}\?\2)\?;\5
\X144:If end of name, \&{exit}\X;\6
\\{incr}\?\1(\?\\{loc}\?\2)\?;\6
\&{if}\1\1\8\|k\?<\?\\{longest\_name}\?-\?1\2\8\&{then}\6
\\{incr}\?\1(\?\|k\?\2)\?;\2\6
\&{end}\8\2\&{if}\1;\6
\&{if}\1\1\?\8\1(\?\|d\?=\?\H{\ }\?\2)\V\1(\?\|d\?=\?\\{tab\_mark}\?\2)\?\2\8%
\&{then}\6
\|d\?\K\?\H{\ };\6
\&{if}\1\1\8\\{mod\_text}\?\1(\?\|k\?-\?1\?\2)=\?\H{\ }\2\8\&{then}\6
\\{decr}\?\1(\?\|k\?\2)\?;\2\6
\&{end}\8\2\&{if}\1;\2\6
\&{end}\8\2\&{if}\1;\6
\\{mod\_text}\?\1(\?\|k\?\2)\K\?\|d;\2\6
\&{end}\8\&{loop};\6
\X145:Check for overlong name\X;\6
\&{if}\1\1\?\8\1(\?\\{mod\_text}\?\1(\?\|k\?\2)=\?\H{\ }\?\2)\W\1(\?\|k\?>\?0\?%
\2)\?\2\8\&{then}\6
\\{decr}\?\1(\?\|k\?\2)\?;\2\6
\&{end}\8\2\&{if}\1\par
\U section~141.\fi

\M144. \P\?\X144:If end of name, \&{exit}\X\S\?\6
\&{if}\1\1\8\|d\?=\?\H{@}\2\8\&{then}\6
\|d\?\K\?\\{buffer}\?\1(\?\\{loc}\?+\?1\?\2)\?;\6
\&{if}\1\1\8\|d\?=\?\H{>}\2\8\&{then}\6
\\{loc}\?\K\?\\{loc}\?+\?2;\5
\&{exit};\6
\4\&{elsif}\1\?\8\1(\?\|d\?=\?\H{\ }\?\2)\V\1(\?\|d\?=\?\\{tab\_mark}\?\2)\V\1(%
\?\|d\?=\?\H{*}\?\2)\?\2\8\&{then}\6
\\{err\_print}\?\1(\?\.{"!\ Section\ name\ didn\'t\ end"}\?\2)\?;\6
\&{exit};\2\6
\&{end}\8\2\&{if}\1;\6
\\{incr}\?\1(\?\|k\?\2)\?;\5
\\{mod\_text}\?\1(\?\|k\?\2)\K\?\H{@};\5
\\{incr}\?\1(\?\\{loc}\?\2)\?;\C{now \|d\?=\?\\{buffer}\?(\?\\{loc}\?)\? again}%
\2\6
\&{end}\8\2\&{if}\1\par
\U section~143.\fi

\M145. \P\?\X145:Check for overlong name\X\S\?\6
\&{if}\1\1\8\|k\?\G\?\\{longest\_name}\?-\?2\2\8\&{then}\6
\\{print\_nl}\?\1(\?\.{"!\ Section\ name\ too\ long:\ "}\?\2)\?;\6
\1\&{for}\8\|j\?\in\?\p{INTEGER}\8\&{range}\81\?\to\?25\8\&{loop}\6
\\{print}\?\1(\?\\{xchr}\?\1(\?\\{mod\_text}\?\1(\?\|j\?\2)\2)\2)\?;\2\6
\&{end}\8\&{loop};\6
\\{print}\?\1(\?\.{"..."}\?\2)\?;\5
\\{mark\_harmless};\2\6
\&{end}\8\2\&{if}\1\par
\U section~143.\fi

\M146. \P\?\X146:If this is the beginning of an character token then set %
\\{scanning\_character} to \p{TRUE} and \?\&{return}\?\8\\{character\_tok};
otherwise do nothing\X\S\?\6
\&{if}\1\1\?\8\1(\?\\{buffer}\?\1(\?\\{loc}\?+\?1\?\2)=\?\H{\'}\?\2)\ORE\1(\?%
\\{buffer}\?\1(\?\\{loc}\?\to\?\\{loc}\?+\?2\?\2)=\1(\?\H{@}\?,\39\?\H{@}\?,\39%
\?\H{\'}\?\2)\2)\?\2\8\&{then}\6
\\{scanning\_character}\?\K\?\p{TRUE};\5
\?\&{return}\?\8\\{character\_tok};\2\6
\&{end}\8\2\&{if}\1\par
\U section~137.\fi

\N147.  Scanning a macro definition.
The rules for generating the replacement texts corresponding to simple
macros, parametric macros, and \ADA\ texts of a module are almost
identical, so a single procedure is used for all three cases. The
differences are that

\yskip\item{a)} The sign \# denotes a parameter only when it appears
outside of strings in a parametric macro; otherwise it stands for the
ASCII character \#.

\item{b)}Module names are not allowed in simple macros or parametric macros;
in fact, the appearance of a module name terminates such macros and denotes
the name of the current module.

\item{c)}The symbols \.{@d} and \.{@f} and \.{@a} are not allowed after
module names, while they terminate macro definitions.

\fi

\M148. Therefore there is a procedure \\{scan\_repl} whose parameter \|t
specifies
either \\{simple} or \\{parametric} or \\{module\_name}. After \\{scan\_repl}
has
acted, \\{cur\_repl\_text} will point to the replacement text just generated,
and
\\{next\_control} will contain the control code that terminated the activity.

\Y\P\?\4\X10:Types and objects visible from \\{data\_structures}\X\mathrel{+}\S%
\?\6
\\{cur\_repl\_text}\?:\?\\{text\_pointer};\C{replacement text formed by \\{scan%
\_repl}}\6
\\{next\_control}\?:\?\\{sixteen\_bits};\par
\fi

\M149. \P\?\X119:Procedures in \\{input\_phase}\X\mathrel{+}\S\?\6
\&{procedure}\8\1\\{scan\_repl}\?(\1\?\|t\?:\?\\{eight\_bits}\?\2)\?\2\8\&{is}%
\8\C{creates a replacement text}\1\6
\|a\?:\?\\{sixteen\_bits};\C{the current token}\6
\|b\?:\?\\{ASCII\_code};\C{a character from the buffer}\6
\\{bal}\?:\?\\{eight\_bits}\?\K\?0;\C{left parentheses minus right parentheses}%
\6
\4\&{begin}\6
\1\&{loop}\6
\4\BL\\{continue}\EL\8\|a\?\K\?\\{get\_next};\6
\&{case}\1\8\|a\1\8\&{is}\8\6
\4\&{when}\8\H{(}\KR\8\\{incr}\?\1(\?\\{bal}\?\2)\?;\6
\4\&{when}\8\H{)}\KR\8\&{if}\1\1\8\\{bal}\?=\?0\2\8\&{then}\8\\{err\_print}\?%
\1(\?\.{"!\ Extra\ )"}\?\2)\?;\8\&{else}\8\\{decr}\?\1(\?\\{bal}\?\2)\?;\2\6
\&{end}\8\2\&{if}\1;\6
\4\&{when}\8\H{"}\KR\8\X152:Copy a string from the buffer to \\{tok\_mem}\X;\6
\4\&{when}\8\\{character\_tok}\?\mathrel{\.|}\?\\{ascii\_tok}\KR\8\\{app\_repl}%
\?\1(\?\|a\?\2)\?;\5
\|a\?\K\?\\{get\_next};\6
\4\&{when}\8\H{\#}\KR\8\&{if}\1\1\8\|t\?=\?\\{parametric}\2\8\&{then}\8\|a\?\K%
\?\\{param};\2\6
\&{end}\8\2\&{if}\1;\6
\hbox{\4}\X151:In cases that \|a is a non-ASCII token (\\{identifier}, %
\\{module\_name}, etc.), either process it and change \|a to a byte that should
be stored, or \&{goto}\8\\{continue} if \|a should be ignored, or \?\&{exit}\8%
\&{when}\?\8\|a signals the end of this replacement text\X\6
\4\&{when}\8\&{others}\8\KR\8\&{null};\2\6
\4\2\&{end}\8\&{case};\6
\\{app\_repl}\?\1(\?\|a\?\2)\?;\C{store \|a in \\{tok\_mem}}\2\6
\&{end}\8\&{loop};\6
\\{next\_control}\?\K\?\|a;\6
\X150:Make sure the parentheses balance\X;\6
\&{if}\1\1\8\\{text\_ptr}\?>\?\\{max\_texts}\?-\?\\{zz}\2\8\&{then}\6
\\{overflow}\?\1(\?\.{"text"}\?\2)\?;\2\6
\&{end}\8\2\&{if}\1;\6
\\{cur\_repl\_text}\?\K\?\\{text\_ptr};\5
\\{tok\_start}\?\1(\?\\{text\_ptr}\?+\?\\{zz}\?\2)\K\?\\{tok\_ptr}\?\1(\?\|z\?%
\2)\?;\5
\\{incr}\?\1(\?\\{text\_ptr}\?\2)\?;\6
\&{if}\1\1\8\|z\?=\?\\{zz}\?-\?1\2\8\&{then}\6
\|z\?\K\?0;\6
\4\&{else}\6
\\{incr}\?\1(\?\|z\?\2)\?;\2\6
\&{end}\8\2\&{if}\1;\2\6
\&{end}\8\\{scan\_repl};\par
\fi

\M150. \P\?\X150:Make sure the parentheses balance\X\S\?\6
\&{if}\1\1\8\\{bal}\?>\?0\2\8\&{then}\6
\&{if}\1\1\8\\{bal}\?=\?1\2\8\&{then}\6
\\{err\_print}\?\1(\?\.{"!\ Missing\ )"}\?\2)\?;\6
\4\&{else}\6
\\{new\_ln};\5
\\{print}\?\1(\?\.{"!\ Missing\ "}\?\2)\?;\5
\\{print\_int}\?\1(\?\\{bal}\?,\39\?1\?\2)\?;\5
\\{print}\?\1(\?\.{"\ )\'s"}\?\2)\?;\5
\\{error};\2\6
\&{end}\8\2\&{if}\1;\6
\1\&{while}\8\\{bal}\?>\?0\8\&{loop}\6
\\{app\_repl}\?\1(\?\H{)}\?\2)\?;\5
\\{decr}\?\1(\?\\{bal}\?\2)\?;\2\6
\&{end}\8\&{loop};\2\6
\&{end}\8\2\&{if}\1\par
\U section~149.\fi

\M151. \P\?\X151:In cases that \|a is a non-ASCII token (\\{identifier}, %
\\{module\_name}, etc.), either process it and change \|a to a byte that should
be stored, or \&{goto}\8\\{continue} if \|a should be ignored, or \?\&{exit}\8%
\&{when}\?\8\|a signals the end of this replacement text\X\S\?\6
\4\&{when}\8\\{identifier}\KR\8\|a\?\K\?\\{id\_lookup}\?\1(\?\\{normal}\?\2)\?;%
\5
\\{app\_repl}\?\1(\1(\?\|a\?/\bn{8}{400}\2)+\bn{8}{200}\2)\?;\5
\|a\?\K\?\|a\?\mathbin{\&{mod}}\bn{8}{400}\?;\6
\4\&{when}\8\\{module\_name}\KR\6
\?\&{exit}\8\&{when}\?\8\|t\?\I\?\\{module\_name};\6
\\{app\_repl}\?\1(\1(\?\\{cur\_module}\?/\bn{8}{400}\2)+\bn{8}{250}\2)\?;\5
\|a\?\K\?\\{cur\_module}\?\mathbin{\&{mod}}\bn{8}{400}\?;\6
\4\&{when}\8\\{verbatim}\KR\8\X153:Copy verbatim string from the buffer to %
\\{tok\_mem}\X;\6
\4\&{when}\8\\{output\_tok}\KR\8\X154:Copy output file-name from the buffer to %
\\{tok\_mem}\X;\6
\4\&{when}\8\\{definition}\?\mathrel{\.|}\?\\{format}\?\mathrel{\.|}\?\\{begin%
\_ada}\KR\6
\?\&{exit}\8\&{when}\?\8\|t\?\I\?\\{module\_name};\6
\\{err\_print}\?\1(\?\.{"!\ @"}\?\mathbin{\.\&}\?\\{xchr}\?\1(\?\\{buffer}\?\1(%
\?\\{loc}\?-\?1\?\2)\2)\mathbin{\.\&}\?\.{"\ is\ ignored\ in\ ADA\ text"}\?\2)%
\?;\6
\&{goto}\8\\{continue};\6
\4\&{when}\8\\{new\_module}\KR\8\&{exit};\par
\U section~149.\fi

\M152. \P\?\X152:Copy a string from the buffer to \\{tok\_mem}\X\S\?\6
\|b\?\K\?\H{"};\6
\1\&{loop}\6
\\{app\_repl}\?\1(\?\|b\?\2)\?;\6
\&{if}\1\1\8\|b\?=\?\H{@}\2\8\&{then}\6
\&{if}\1\1\8\\{buffer}\?\1(\?\\{loc}\?\2)=\?\H{@}\2\8\&{then}\6
\\{incr}\?\1(\?\\{loc}\?\2)\?;\C{store only one \.{@}}\6
\4\&{else}\6
\\{err\_print}\?\1(\?\.{"!\ You\ should\ double\ @\ signs\ in\ strings"}\?\2)%
\?;\2\6
\&{end}\8\2\&{if}\1;\2\6
\&{end}\8\2\&{if}\1;\6
\&{if}\1\1\8\\{loc}\?=\?\\{limit}\2\8\&{then}\6
\\{err\_print}\?\1(\?\.{"!\ String\ didn\'t\ end"}\?\2)\?;\6
\\{buffer}\?\1(\?\\{loc}\?\to\?\\{loc}\?+\?1\?\2)\K\1(\?\H{"}\?,\39\?0\?\2)\?;%
\2\6
\&{end}\8\2\&{if}\1;\6
\|b\?\K\?\\{buffer}\?\1(\?\\{loc}\?\2)\?;\5
\\{incr}\?\1(\?\\{loc}\?\2)\?;\6
\&{if}\1\1\8\|b\?=\?\H{"}\2\8\&{then}\6
\?\&{exit}\8\&{when}\?\8\\{buffer}\?\1(\?\\{loc}\?\2)\I\?\H{"};\6
\\{incr}\?\1(\?\\{loc}\?\2)\?;\5
\\{app\_repl}\?\1(\?\H{"}\?\2)\?;\2\6
\&{end}\8\2\&{if}\1;\2\6
\&{end}\8\&{loop}\C{now \|a holds the final \H{"} that will be stored}\par
\U section~149.\fi

\M153. \P\?\X153:Copy verbatim string from the buffer to \\{tok\_mem}\X\S\?\6
\\{app\_repl}\?\1(\?\\{verbatim}\?\2)\?;\5
\\{buffer}\?\1(\?\\{limit}\?+\?1\?\2)\K\?\H{@};\6
\1\&{while}\8\\{buffer}\?\1(\?\\{loc}\?\2)\I\?\H{@}\8\&{loop}\6
\\{app\_repl}\?\1(\?\\{buffer}\?\1(\?\\{loc}\?\2)\2)\?;\5
\\{incr}\?\1(\?\\{loc}\?\2)\?;\6
\&{if}\1\1\?\8\1(\?\\{loc}\?<\?\\{limit}\?\2)\ANT\1(\?\\{buffer}\?\1(\?\\{loc}%
\?\to\?\\{loc}\?+\?1\?\2)=\1(\?\H{@}\?,\39\?\H{@}\?\2)\2)\?\2\8\&{then}\6
\\{app\_repl}\?\1(\?\H{@}\?\2)\?;\5
\\{loc}\?\K\?\\{loc}\?+\?2;\2\6
\&{end}\8\2\&{if}\1;\2\6
\&{end}\8\&{loop};\6
\&{if}\1\1\8\\{loc}\?\G\?\\{limit}\2\8\&{then}\6
\\{err\_print}\?\1(\?\.{"!\ Verbatim\ string\ didn\'t\ end"}\?\2)\?;\6
\&{if}\1\1\8\\{loc}\?\G\?\\{buf\_size}\?-\?1\2\8\&{then}\6
\\{loc}\?\K\?\\{buf\_size}\?-\?2;\2\6
\&{end}\8\2\&{if}\1;\6
\4\&{elsif}\1\8\\{buffer}\?\1(\?\\{loc}\?+\?1\?\2)\I\?\H{>}\2\8\&{then}\6
\\{err\_print}\?\1(\?\.{"!\ You\ should\ double\ @\ signs\ in\ verbatim\
strings"}\?\2)\?;\2\6
\&{end}\8\2\&{if}\1;\6
\\{loc}\?\K\?\\{loc}\?+\?2\C{another \\{verbatim} byte will be stored, since %
\|a\?=\?\\{verbatim}}\par
\U section~151.\fi

\M154. \P\?\X154:Copy output file-name from the buffer to \\{tok\_mem}\X\S\?\6
\\{app\_repl}\?\1(\?\\{output\_tok}\?\2)\?;\5
\\{buffer}\?\1(\?\\{limit}\?+\?1\?\2)\K\?\H{@};\6
\1\&{while}\8\\{buffer}\?\1(\?\\{loc}\?\2)\I\?\H{@}\8\&{loop}\6
\\{app\_repl}\?\1(\?\\{buffer}\?\1(\?\\{loc}\?\2)\2)\?;\5
\\{incr}\?\1(\?\\{loc}\?\2)\?;\6
\&{if}\1\1\?\8\1(\?\\{loc}\?<\?\\{limit}\?\2)\ANT\1(\?\\{buffer}\?\1(\?\\{loc}%
\?\to\?\\{loc}\?+\?1\?\2)=\1(\?\H{@}\?,\39\?\H{@}\?\2)\2)\?\2\8\&{then}\6
\\{app\_repl}\?\1(\?\H{@}\?\2)\?;\5
\\{loc}\?\K\?\\{loc}\?+\?2;\2\6
\&{end}\8\2\&{if}\1;\2\6
\&{end}\8\&{loop};\6
\&{if}\1\1\8\\{loc}\?\G\?\\{limit}\2\8\&{then}\6
\\{err\_print}\?\1(\?\.{"!\ Output\ file\ name\ didn\'t\ end"}\?\2)\?;\6
\4\&{elsif}\1\8\\{buffer}\?\1(\?\\{loc}\?+\?1\?\2)\I\?\H{>}\2\8\&{then}\6
\\{err\_print}\?\1(\?\.{"!\ You\ should\ double\ @\ signs\ in\ output\
file-names"}\?\2)\?;\2\6
\&{end}\8\2\&{if}\1;\6
\\{loc}\?\K\?\\{loc}\?+\?2\C{another \\{output\_tok} byte will be stored, since
\|a\?=\?\\{output\_tok}}\par
\U section~151.\fi

\M155. The following procedure is used to define a simple or parametric macro,
just after the `\.{==}' of its definition has been scanned.

\Y\P\?\4\X119:Procedures in \\{input\_phase}\X\mathrel{+}\S\?\6
\&{procedure}\8\1\\{define\_macro}\?(\1\?\|t\?:\?\\{eight\_bits}\?\2)\?\2\8%
\&{is}\8\1\6
\|p\?:\?\\{name\_pointer};\C{the identifier being defined}\6
\4\&{begin}\6
\|p\?\K\?\\{id\_lookup}\?\1(\?\|t\?\2)\?;\5
\\{scan\_repl}\?\1(\?\|t\?\2)\?;\6
\\{equiv}\?\1(\?\|p\?\2)\K\?\\{cur\_repl\_text};\5
\\{text\_link}\?\1(\?\\{cur\_repl\_text}\?\2)\K\?0;\2\6
\&{end}\8\\{define\_macro};\par
\fi

\N156.  Scanning a module.
The \\{scan\_module} procedure starts when `\.{@\ }' or `\.{@*}' has been
sensed in the input, and it proceeds until the end of that module.  It
uses \\{module\_count} to keep track of the current module number; with luck,
\.{AWEAVE} and \.{ATANGLE} will both assign the same numbers to modules.

\Y\P\?\4\X10:Types and objects visible from \\{data\_structures}\X\mathrel{+}\S%
\?\6
\\{module\_count}\?:\?\\{mod\_range};\C{the current module number}\par
\fi

\M157. The top level of \\{scan\_module} is trivial.

\Y\P\?\4\X119:Procedures in \\{input\_phase}\X\mathrel{+}\S\?\6
\&{procedure}\8\1\\{scan\_module}\2\8\&{is}\8\1\6
\|p\?:\?\\{name\_pointer};\C{module name for the current module}\6
\4\&{begin}\6
\\{incr}\?\1(\?\\{module\_count}\?\2)\?;\6
\X158:Scan the definition part of the current module\X;\6
\X160:Scan the \ADA\ part of the current module\X;\2\6
\&{end}\8\\{scan\_module};\par
\fi

\M158. \P\?\X158:Scan the definition part of the current module\X\S\?\6
\\{next\_control}\?\K\?0;\6
\1\&{loop}\6
\4\BL\\{continue}\EL\8\1\&{while}\8\\{next\_control}\?\L\?\\{format}\8\&{loop}\6
\\{next\_control}\?\K\?\\{skip\_ahead};\6
\&{if}\1\1\8\\{next\_control}\?=\?\\{module\_name}\2\8\&{then}\C{we want to
scan the module name too}\6
\\{loc}\?\K\?\\{loc}\?-\?2;\5
\\{next\_control}\?\K\?\\{get\_next};\2\6
\&{end}\8\2\&{if}\1;\2\6
\&{end}\8\&{loop};\6
\?\&{exit}\8\&{when}\?\8\\{next\_control}\?\I\?\\{definition};\6
\\{next\_control}\?\K\?\\{get\_next};\C{get identifier name}\6
\&{if}\1\1\8\\{next\_control}\?\I\?\\{identifier}\2\8\&{then}\6
\\{err\_print}\?\1(\?\.{"!\ Definition\ flushed,\ must\ start\ with\ "}\?%
\mathbin{\.\&}\?\.{"identifier\ of\ length\ >\ 1"}\?\2)\?;\6
\&{goto}\8\\{continue};\2\6
\&{end}\8\2\&{if}\1;\6
\\{next\_control}\?\K\?\\{get\_next};\C{get token after the identifier}\6
\&{if}\1\1\8\\{next\_control}\?=\?\\{equivalence\_sign}\2\8\&{then}\6
\\{define\_macro}\?\1(\?\\{simple}\?\2)\?;\5
\&{goto}\8\\{continue};\6
\4\&{else}\6
\X159:If the next text is `\?(\?\#\?)\S\?', call \\{define\_macro} and \&{goto}%
\8\\{continue}\X;\2\6
\&{end}\8\2\&{if}\1;\6
\\{err\_print}\?\1(\?\.{"!\ Definition\ flushed\ since\ it\ starts\ badly"}\?%
\2)\?;\2\6
\&{end}\8\&{loop}\par
\U section~157.\fi

\M159. \P\?\X159:If the next text is `\?(\?\#\?)\S\?', call \\{define\_macro}
and \&{goto}\8\\{continue}\X\S\?\6
\&{if}\1\1\8\\{next\_control}\?=\?\H{(}\2\8\&{then}\6
\\{next\_control}\?\K\?\\{get\_next};\6
\&{if}\1\1\8\\{next\_control}\?=\?\H{\#}\2\8\&{then}\6
\\{next\_control}\?\K\?\\{get\_next};\6
\&{if}\1\1\8\\{next\_control}\?=\?\H{)}\2\8\&{then}\6
\\{next\_control}\?\K\?\\{get\_next};\6
\&{if}\1\1\8\\{next\_control}\?=\?\H{=}\2\8\&{then}\6
\\{err\_print}\?\1(\?\.{"!\ Use\ ==\ for\ macros"}\?\2)\?;\6
\\{next\_control}\?\K\?\\{equivalence\_sign};\2\6
\&{end}\8\2\&{if}\1;\6
\&{if}\1\1\8\\{next\_control}\?=\?\\{equivalence\_sign}\2\8\&{then}\6
\\{define\_macro}\?\1(\?\\{parametric}\?\2)\?;\5
\&{goto}\8\\{continue};\2\6
\&{end}\8\2\&{if}\1;\2\6
\&{end}\8\2\&{if}\1;\2\6
\&{end}\8\2\&{if}\1;\2\6
\&{end}\8\2\&{if}\1\par
\U section~158.\fi

\M160. \P\?\X160:Scan the \ADA\ part of the current module\X\S\?\6
\&{case}\1\8\\{next\_control}\1\8\&{is}\8\6
\4\&{when}\8\\{begin\_ada}\KR\6
\|p\?\K\?0;\6
\4\&{when}\8\\{module\_name}\KR\6
\|p\?\K\?\\{cur\_module};\6
\X161:Check that \?=\? or \?\S\? follows this module name, otherwise \&{return}%
\X;\6
\4\&{when}\8\&{others}\8\KR\6
\&{return};\2\6
\4\2\&{end}\8\&{case};\6
\X162:Insert the module number into \\{tok\_mem}\X;\6
\\{scan\_repl}\?\1(\?\\{module\_name}\?\2)\?;\C{now \\{cur\_repl\_text} points
to the replacement text}\6
\X163:Update the data structure so that the replacement text is accessible\X\par
\U section~157.\fi

\M161. \P\?\X161:Check that \?=\? or \?\S\? follows this module name, otherwise
\&{return}\X\S\?\6
\1\&{repeat}\6
\\{next\_control}\?\K\?\\{get\_next};\2\6
\&{until}\?\8\1(\?\\{next\_control}\?\I\?\H{+}\?\2)\?;\C{allow optional `%
\.{+=}'}\6
\&{if}\1\1\?\8\1(\?\\{next\_control}\?\I\?\H{=}\?\2)\W\1(\?\\{next\_control}\?%
\I\?\\{equivalence\_sign}\?\2)\?\2\8\&{then}\6
\\{err\_print}\?\1(\?\.{"!\ ADA\ text\ flushed,\ =\ sign\ is\ missing"}\?\2)\?;%
\6
\1\&{repeat}\6
\\{next\_control}\?\K\?\\{skip\_ahead};\2\6
\&{until}\?\8\1(\?\\{next\_control}\?=\?\\{new\_module}\?\2)\?;\6
\&{return};\2\6
\&{end}\8\2\&{if}\1\par
\U section~160.\fi

\M162. \P\?\X162:Insert the module number into \\{tok\_mem}\X\S\?\6
\\{store\_two\_bytes}\?\1(\bn{8}{150000}+\?\\{module\_count}\?\2)\?\C{\?%
\bn{8}{150000}=\bn{8}{320}\ast\bn{8}{400}\?}\par
\U section~160.\fi

\M163. \P\?\X163:Update the data structure so that the replacement text is
accessible\X\S\?\6
\&{if}\1\1\8\|p\?=\?0\2\8\&{then}\C{unnamed module}\6
\\{text\_link}\?\1(\?\\{last\_unnamed}\?\2)\K\?\\{cur\_repl\_text};\5
\\{last\_unnamed}\?\K\?\\{cur\_repl\_text};\6
\4\&{elsif}\1\8\\{equiv}\?\1(\?\|p\?\2)=\?0\2\8\&{then}\6
\\{equiv}\?\1(\?\|p\?\2)\K\?\\{cur\_repl\_text};\C{first module of this name}\6
\4\&{else}\6
\|p\?\K\?\\{equiv}\?\1(\?\|p\?\2)\?;\6
\1\&{while}\8\\{text\_link}\?\1(\?\|p\?\2)<\?\\{module\_flag}\8\&{loop}\6
\|p\?\K\?\\{text\_link}\?\1(\?\|p\?\2)\?;\C{find end of list}\2\6
\&{end}\8\&{loop};\6
\\{text\_link}\?\1(\?\|p\?\2)\K\?\\{cur\_repl\_text};\2\6
\&{end}\8\2\&{if}\1;\6
\\{text\_link}\?\1(\?\\{cur\_repl\_text}\?\2)\K\?\\{module\_flag}\C{mark this
replacement text as a nonmacro}\par
\U section~160.\fi

\M164. \P\?\X119:Procedures in \\{input\_phase}\X\mathrel{+}\S\?\6
\&{procedure}\8\1\\{phase\_i}\2\8\&{is}\8\1\6
\4\&{begin}\6
\X126:Initialize the input system\X;\6
\\{print\_ln}\?\1(\?\\{banner}\?\2)\?;\C{print a ``banner line''}\6
\X167:Phase I: Read all the user's text and compress it into \\{tok\_mem}\X;\2\6
\&{end}\8\\{phase\_i};\par
\fi

\N165.  The main program.
We have defined plenty of procedures, and it is time to put the last
pieces of the puzzle in place. Here is where \.{ATANGLE} starts, and where
it ends.

\Y\P\4\D\\{web\_input\_file\_name}\?\S\?\.{"input.aweb"}\?\6
\?\par
\P\4\D\\{change\_input\_file\_name}\?\S\?\.{"input.ach"}\?\6
\?\par
\P\4\D\\{ada\_output\_file\_name}\?\S\?\.{"web\_output.a"}\?\6
\?\par
\Y\P\?\4\X10:Types and objects visible from \\{data\_structures}\X\mathrel{+}\S%
\?\6
\\{opened}\?:\?\p{BOOLEAN}\?\K\?\p{FALSE};\C{was previous \\{open\_a\_file}
succesfull?}\par
\fi

\M166\*.
\Y\P\O{atangle.adb}\?\6
\?\&{with}\8\\{Ada}\?.\?\\{Command\_Line};\6
\&{use}\8\\{Ada}\?.\?\\{Command\_Line};\6
\&{with}\8\p{TEXT\_IO}\?,\39\?\\{int\_number\_io};\6
\&{with}\8\\{data\_structures}\?,\39\?\\{input\_output}\?,\39\?\\{input\_phase}%
\?,\39\?\\{output\_phase};\6
\&{use}\8\\{data\_structures}\?,\39\?\\{input\_output}\?,\39\?\\{input\_phase}%
\?,\39\?\\{output\_phase};\6
\&{procedure}\8\1\\{ATANGLE}\2\8\&{is}\8\1\6
\\{ac}\?:\?\p{INTEGER}\?\K\?0;\6
\4\&{begin}\6
\\{print\_ln}\?\1(\?\.{"Trying\ to\ open\ input\ file."}\?\2)\?;\5
\\{ac}\?\K\?\\{Argument\_Count};\6
\&{if}\1\1\8\\{ac}\?<\?1\2\8\&{then}\6
\\{fatal\_error}\?\1(\?\.{"No\ file\ on\ command\ line"}\?\2)\?;\2\6
\&{end}\8\2\&{if}\1;\6
\\{open\_a\_file}\?\1(\?\\{web\_file}\?,\39\?\p{TEXT\_IO}\?.\?\p{IN\_FILE}\?,%
\39\?\\{Argument}\?\1(\?1\?\2),\39\?\\{opened}\?\2)\?;\6
\&{if}\1\1\?\8\R\?\\{opened}\2\8\&{then}\6
\\{fatal\_error}\?\1(\?\.{"AWEB\ file\ not\ open"}\?\2)\?;\2\6
\&{end}\8\2\&{if}\1;\6
\\{opened}\?\K\?\p{FALSE};\6
\&{if}\1\1\8\\{ac}\?\G\?2\2\8\&{then}\6
\\{open\_a\_file}\?\1(\?\\{change\_file}\?,\39\?\p{TEXT\_IO}\?.\?\p{IN\_FILE}%
\?,\39\?\\{Argument}\?\1(\?2\?\2),\39\?\\{opened}\?\2)\?;\2\6
\&{end}\8\2\&{if}\1;\6
\&{if}\1\1\?\8\R\?\\{opened}\2\8\&{then}\6
\\{new\_ln};\5
\\{print\_ln}\?\1(\?\.{"!\ No\ change\ file."}\?\2)\?;\2\6
\&{end}\8\2\&{if}\1;\6
\\{open\_a\_file}\?\1(\?\\{ada\_file}\?,\39\?\p{TEXT\_IO}\?.\?\p{OUT\_FILE}\?,%
\39\?\\{ada\_output\_file\_name}\?,\39\?\\{opened}\?\2)\?;\6
\&{if}\1\1\?\8\1(\R\?\\{opened}\?\2)\ANT\1(\R\?\\{create\_ada\_file}\?\1(\?%
\\{ada\_output\_file\_name}\?\2)\2)\?\2\8\&{then}\6
\\{fatal\_error}\?\1(\?\.{"ADA\ file\ not\ open"}\?\2)\?;\2\6
\&{end}\8\2\&{if}\1;\6
\\{fn\_length}\?\K\?\p{TEXT\_IO}\?.\?\p{NAME}\?\1(\?\\{ada\_file}\?\2)'\?%
\p{LAST};\5
\\{file\_name}\?\1(\?1\?\to\?\\{fn\_length}\?\2)\K\?\p{TEXT\_IO}\?.\?\p{NAME}\?%
\1(\?\\{ada\_file}\?\2)\?;\5
\\{f\_line}\?\K\?1;\5
\\{phase\_i};\6
\&{stat}\1\6
\4\&{for}\8\\{zo}\?\in\?\\{z\_range}\8\&{loop}\6
\\{max\_tok\_ptr}\?\1(\?\\{zo}\?\2)\K\?\\{tok\_ptr}\?\1(\?\\{zo}\?\2)\?;\2\6
\&{end}\8\&{loop};\2\6
\&{tats}\6
\\{phase\_ii};\5
\&{raise}\8\\{end\_of\_ATANGLE}; \4\&{exception} \6
\4\&{when}\8\\{end\_of\_ATANGLE}\KR\6
\&{stat}\1\6
\X168:Print statistics about memory usage\X;\2\6
\&{tats}\6
\hbox{\4\4}\C{here files should be closed if the operating system requires it}\6
\X169:Print the job \\{history}\X;\2\6
\&{end}\8\\{ATANGLE};\par
\fi

\M167. \P\?\X167:Phase I: Read all the user's text and compress it into \\{tok%
\_mem}\X\S\?\6
\\{phase\_one}\?\K\?\p{TRUE};\5
\\{module\_count}\?\K\?0;\6
\1\&{repeat}\6
\\{next\_control}\?\K\?\\{skip\_ahead};\2\6
\&{until}\?\8\1(\?\\{next\_control}\?=\?\\{new\_module}\?\2)\?;\6
\1\&{while}\?\8\R\?\\{input\_has\_ended}\8\&{loop}\6
\\{scan\_module};\2\6
\&{end}\8\&{loop};\6
\X130:Check that all changes have been read\X;\6
\\{phase\_one}\?\K\?\p{FALSE}\par
\U section~164.\fi

\M168. \P\?\X168:Print statistics about memory usage\X\S\?\6
\\{print\_nl}\?\1(\?\.{"Memory\ usage\ statistics:"}\?\2)\?;\5
\\{new\_ln};\5
\\{print\_int}\?\1(\?\\{name\_ptr}\?,\39\?1\?\2)\?;\5
\\{print}\?\1(\?\.{"\ names,\ "}\?\2)\?;\5
\\{print\_int}\?\1(\?\\{text\_ptr}\?,\39\?1\?\2)\?;\5
\\{print\_ln}\?\1(\?\.{"\ replacement\ texts;"}\?\2)\?;\5
\\{print\_int}\?\1(\?\\{byte\_ptr}\?\1(\?0\?\2),\39\?1\?\2)\?;\6
\1\&{for}\8\\{wo}\?\in\?1\?\to\?\\{ww}\?-\?1\8\&{loop}\6
\\{print}\?\1(\?\.{"+"}\?\2)\?;\5
\\{print\_int}\?\1(\?\\{byte\_ptr}\?\1(\?\\{wo}\?\2),\39\?1\?\2)\?;\2\6
\&{end}\8\&{loop};\6
\\{print}\?\1(\?\.{"\ bytes,\ "}\?\2)\?;\5
\\{print\_int}\?\1(\?\\{max\_tok\_ptr}\?\1(\?0\?\2),\39\?1\?\2)\?;\6
\1\&{for}\8\\{zo}\?\in\?1\?\to\?\\{zz}\?-\?1\8\&{loop}\6
\\{print}\?\1(\?\.{"+"}\?\2)\?;\5
\\{print\_int}\?\1(\?\\{max\_tok\_ptr}\?\1(\?\\{zo}\?\2),\39\?1\?\2)\?;\2\6
\&{end}\8\&{loop};\6
\\{print}\?\1(\?\.{"\ tokens."}\?\2)\?\par
\U section~166\*.\fi

\M169. Some implementations may wish to pass the \\{history} value to the
operating system so that it can be used to govern whether or not other
programs are started. Here we simply report the history to the user.

\Y\P\?\4\X169:Print the job \\{history}\X\S\?\6
\&{case}\1\8\\{history}\1\8\&{is}\8\6
\4\&{when}\8\\{spotless}\KR\8\\{print\_nl}\?\1(\?\.{"(No\ errors\ were\
found.)"}\?\2)\?;\6
\4\&{when}\8\\{harmless\_message}\KR\8\\{print\_nl}\?\1(\?\.{"(Did\ you\ see\
the\ warning\ message\ above?)"}\?\2)\?;\6
\4\&{when}\8\\{error\_message}\KR\8\\{print\_nl}\?\1(\?\.{"(Pardon\ me,\ but\ I%
\ think\ I\ spotted\ something\ wrong.)"}\?\2)\?;\6
\4\&{when}\8\\{fatal\_message}\KR\8\\{print\_nl}\?\1(\?\.{"(That\ was\ a\ fatal%
\ error,\ my\ friend.)"}\?\2)\?;\2\6
\4\2\&{end}\8\&{case};\C{there are no other cases}\6
\\{new\_ln}\par
\U section~166\*.\fi

\N170\*.  System-dependent changes.
This module should be replaced, if necessary, by changes to the program
that are necessary to make \.{ATANGLE} work at a particular installation.
It is usually best to design your change file so that all changes to
previous modules preserve the module numbering; then everybody's version
will be consistent with the printed program. More extensive changes,
which introduce new modules, can be inserted here; then only the index
itself will get a new module number.

\Y\P\O{int\_number\_io.ads}\?\6
\?\&{with}\8\p{TEXT\_IO};\6
\&{package}\8\1\\{int\_number\_io}\?\8\&{is}\8\&{new}\?\8\p{TEXT\_IO}\?.\?%
\p{INTEGER\_IO}\?\1(\?\p{INTEGER}\?\2)\?\2;\par
\fi

\N171\*.  Index.
Here is a cross-reference table for the \.{ATANGLE} processor.
All modules in which an identifier is
used are listed with that identifier, except that reserved words are
indexed only when they appear in format definitions, and the appearances
of identifiers in module names are not indexed.
Underlined entries of subprograms and packages correspond to sections where
this entity is specified, whereas entries in italic type correspond to
the section where the entity's body is stated.
For any other identifier underlined entries correspond to where the
identifier was declared. Error messages and
a few other things like ``ASCII code'' are indexed here too.




\fi


\ch 1\*, 3\*, 4\*, 5\*, 6\*, 7\*, 166\*, 170\*, 171\*.
\inx
\:\.{\AT!a is ignored in ADA text}, 151.
\:\.{\AT!d is ignored in ADA text}, 151.
\:\.{\AT!f is ignored in ADA text}, 151.
\:\|{a}, \[84], \[149].
\:\\{ac}, 166\*.
\:\\{Ada}, 166\*.
\:{ADA attributes}, 13--14.
\:\.{ADA file not open}, 166\*.
\:\.{ADA text flushed...}, 161.
\:\\{ada\_file}, \[25], 26, 94, 108, 166\*.
\:\\{ada\_file\_name}, 23, \[26].
\:\\{ada\_output\_file\_name}, \[165], 166\*.
\:\.{Ambiguous prefix}, 69.
\:\\{and\_sign}, 16, 102.
\:\\{app}, \[96], 98--99.
\:\\{app\_repl}, \[90], 149--154.
\:\\{Argument}, 166\*.
\:\\{Argument\_Count}, 166\*.
\:{ASCII code}, 12, 72.
\:{ASCII code constant}, 72.
\:\\{ASCII\_code}, \[12], 14, 16, 27--28, 39, 49, 62, 132, 134, 149.
\:\\{ascii\_tok}, \[72], 90, 102, 132, 140, 149.
\:\\{ATANGLE}, \\{166\*}.
\:\.{AWEB file ended...}, 124.
\:\.{AWEB file not open}, 166\*.
\:\|{b}, \[84], \[94], \[149].
\:\\{bal}, \[84], 90, \[134], 135, \[149], 150.
\:\\{banner}, \[1\*], 164.
\:\\{begin\_ada}, \[131], 132, 151, 160.
\:\\{begin\_comment}, \[72], 112, 132.
\:\p{BOOLEAN}, 23--24, 26, 28--30, 92, 109, 116, 119, 136, 165.
\:\\{box\_sign}, 16, 102, 138.
\:\\{brace\_level}, \[79], 80, 95, 112.
\:\\{break\_ptr}, \[91], 92--95, 98--99, 113.
\:\\{buf\_size}, \[10], 27--29, 153.
\:\\{buffer}, \[27], 28--29, 32--33, 49, 51--52, 54, 56, 59, 119, 121,
123--130, 133--135, 137--140, 143--144, 146, 151--154.
\:\\{buffer\_index}, \[27], 29, 32, 49, 51, 116, 118.
\:\\{buffer\_type}, \[27], 91, 97, 118.
\:\\{byte\_field}, \[76].
\:\\{byte\_index}, \[39], 40, 51, 65, 69, 84, 101.
\:\\{byte\_mem}, 38, \[39], 40, 45, 48, 51, 54, 59, 61, 65--67, 69, 84, 101,
104.
\:\\{byte\_memory}, \[39].
\:\\{byte\_ptr}, \[40], 59, 66, 87, 168.
\:\\{byte\_start}, 38, \[39], 40--41, 48--49, 54, 59, 61, 66--67, 76, 78, 87,
104, 136.
\:\|{c}, \[29], \[51], \[65], \[69], \[132], \[133], \[134], \[137].
\:\.{Can't output ASCII code n}, 101.
\:\\{carriage\_return}, 16--17, 28.
\:\.{Change file ended...}, 122, 124, 129.
\:\.{Change file entry did not match}, 130.
\:\\{change\_buffer}, \[118], 119--120, 123--124, 128, 130.
\:\\{change\_changing}, \[117], 124, 126, 129.
\:\\{change\_file}, \[22], 33, 116, 118, 121--122, 124, 129, 166\*.
\:\\{change\_input\_file\_name}, \[165].
\:\\{change\_limit}, \[118], 119--120, 123--124, 128, 130.
\:\\{changing}, 33, \[116], 117--118, 120, 124, 126--127, 130.
\:\p{CHARACTER}, 13, 15.
\:\\{character\_tok}, \[72], 90, 102, 146, 149.
\:\\{check\_break}, \[94], 98.
\:\\{check\_change}, \\{124}, 128.
\:\\{chop\_hash}, \[49], 58, 60.
\:\\{chop\_hash\_index}, \[49], 51.
\:\\{chopped\_id}, \[49], 51, 56, 61.
\:\p{CLOSE}, 108.
\:\\{Command\_Line}, 166\*.
\:\\{comment\_output}, \[92], 93--94, 112.
\:\\{compare}, \[63], 65, 69.
\:\\{compress}, \[138].
\:\\{confusion}, \[36], 86.
\:\\{continue}, 101, 121, 149, 151, 158--159.
\:\\{control\_code}, 131, \\{132}, 133, 136, 140.
\:\\{control\_text}, \[131], 132, 140.
\:\.{Could not open file `x'}, 108.
\:\\{count}, \[69].
\:\p{CREATE}, 26.
\:\\{create\_ada\_file}, \[23], \\{26}, 108, 166\*.
\:\\{cur\_byte}, \[76], 80--82, 84, 87, 90.
\:\\{cur\_char}, \[101], 104--105, 108, 110--112.
\:\\{cur\_end}, \[76], 80--82, 84, 87.
\:\\{cur\_mod}, \[76], 80--81, 84.
\:\\{cur\_module}, \[136], 141, 151, 160.
\:\\{cur\_name}, \[76], 80--82.
\:\\{cur\_repl}, \[76], 77, 80--82.
\:\\{cur\_repl\_text}, \[148], 149, 155, 160, 163.
\:\\{cur\_state}, \[76], 81--82.
\:\\{cur\_val}, \[83], 84, 86, 104, 112.
\:\|{d}, \[137].
\:\\{data\_structures}, \\{3\*}, 4\*--7\*, 166\*.
\:\.{data\_structures.adb}, 3\*.
\:\.{data\_structures.ads}, 3\*.
\:\&{declare}, 8.
\:\\{decr}, \[9], 29, 82, 88, 90, 96, 112, 135, 139, 143, 149--150.
\:\\{default\_name}, 23, \[24].
\:\\{define\_macro}, \\{155}, 158--159.
\:\\{definition}, \[131], 132, 151, 158.
\:\.{Definition flushed...}, 158.
\:\\{double\_dot}, \[72], 102, 138.
\:\\{double\_star}, 16, 102, 138.
\:\\{eight\_bits}, \[38], 39, 50--51, 72, 79, 84, 92, 97--98, 101, 131--134,
137, 149, 155.
\:\&{end}, 8--9.
\:\\{end\_comment}, \[72], 112, 132.
\:\\{end\_field}, \[76].
\:\\{end\_of\_ATANGLE}, \[35], \\{166\*}.
\:\p{END\_OF\_FILE}, 29.
\:\p{END\_OF\_LINE}, 29.
\:\\{equal}, \[63], 65--67.
\:\\{equiv}, 38, \[39], 44--46, 49, 58, 60--61, 66, 81, 85, 87, 155, 163.
\:\\{equivalence\_sign}, 16, 102, 138, 158--159, 161.
\:\\{err\_print}, \[31], 57, 65, 69, 94, 105--106, 108, 111--112, 117,
121--122, 124--125, 129--130, 134--135, 137, 140, 143--144, 149--154, 158--159,
161.
\:\\{error}, 29, \[31], \\{32}, 35, 61, 85, 87, 95, 101, 125, 150.
\:\\{error\_message}, \[11], 169.
\:\\{extension}, \[63], 67.
\:\.{Extra )}, 149.
\:\.{Extra \AT!\}}, 112.
\:\|{f}, \[24], \[29].
\:\\{f\_line}, 34, 94, \[107], 108, 166\*.
\:\p{FALSE}, 24, 26, 28--30, 93, 109--110, 112, 117--119, 124, 126, 136--137,
165--167.
\:\\{fatal\_error}, \[35], 36--37, 108, 166\*.
\:\\{fatal\_message}, \[11], 169.
\:\p{FILE}, 24, 26.
\:\.{File name should have length > 0}, 108.
\:\p{FILE\_MODE}, 23--24.
\:\\{file\_name}, 34, \[107], 108, 166\*.
\:\p{FILE\_TYPE}, 13.
\:\\{final\_limit}, \[29].
\:\p{FIRST}, 13.
\:\\{flush\_buffer}, \\{94}, 95, 113.
\:\\{fn\_length}, 34, \[107], 108, 166\*.
\:\\{force\_line}, \[72], 101, 132.
\:\\{form\_feed}, 16, 28.
\:\\{format}, \[131], 132, 151, 158.
\:\\{found}, 53--54, 65.
\:\\{frac}, \[97], 98--99, 101, 111.
\:\.{Fraction too long}, 111.
\:\p{GET}, 29.
\:\\{get\_fraction}, 101, 110--111.
\:\\{get\_line}, 116, \\{127}, 133--134, 137, 143.
\:\\{get\_next}, 136, \\{137}, 149, 158--159, 161.
\:\\{get\_output}, 83, \\{84}, 91, 100--102, 105--106, 108, 110--112.
\:\\{greater}, \[63], 65, 67, 69.
\:\\{greater\_or\_equal}, 16, 102, 138.
\:\|{h}, \[51].
\:\\{harmless\_message}, \[11], 169.
\:\\{hash}, 40, \[49], 53.
\:\\{hash\_index}, \[49], 51.
\:\\{hash\_size}, \[10], 49, 52, 56.
\:\\{hashing}, \\{5\*}, 7\*.
\:\.{hashing.adb}, 5\*.
\:\.{hashing.ads}, 5\*.
\:\\{history}, \[11], 169.
\:\\{history\_type}, \[11].
\:\.{Hmm... n of the preceding...}, 125.
\:\|{i}, \[18], \[51].
\:\\{id\_first}, \[49], 51--52, 54, 56, 59, 136, 139.
\:\\{id\_loc}, \[49], 51--52, 54, 56, 59, 136, 139.
\:\\{id\_lookup}, 49, \[50], \\{51}, 136, 151, 155.
\:\\{ident}, \[97], 98--99, 102, 104, 110.
\:\\{identifier}, \[83], 86, 104, 136, 139, 151, 158.
\:\.{Identifier conflict...}, 61.
\:\\{ignore}, \[131], 132--133, 140.
\:\\{ilk}, 38, \[39], 44--45, 49, 55, 57--59, 82, 86--87.
\:\.{Improper \AT! within control text}, 140.
\:\.{Improper ASCII-code sequence...}, 140.
\:\.{Improper based number...}, 137.
\:\p{IN\_FILE}, 166\*.
\:\.{Incompatible module names}, 65.
\:\\{incr}, \[9], 29, 52, 54, 56, 59, 61, 66--67, 69, 81, 84, 87, 90, 94, 96,
104--106, 108, 110--112, 121--122, 124, 128--129, 133--135, 137--140, 143--144,
149, 152--154, 157.
\:\.{Input ended in mid-comment}, 134.
\:\.{Input ended in section name}, 143.
\:\.{Input line too long}, 29.
\:\\{input\_has\_ended}, \[116], 124, 126, 128, 133--134, 137, 143, 167.
\:\\{input\_ln}, \[28], \\{29}, 121--122, 124, 128--129.
\:\\{input\_output}, 3\*, \\{4\*}, 5\*--7\*, 166\*.
\:\.{input\_output.adb}, 4\*.
\:\.{input\_output.ads}, 4\*.
\:\\{input\_phase}, \\{7\*}, 166\*.
\:\.{input\_phase.adb}, 7\*.
\:\.{input\_phase.ads}, 7\*.
\:\\{int\_number\_io}, 3\*--4\*, 6\*--7\*, 20, 166\*, \\{170\*}.
\:\.{int\_number\_io.ads}, 170\*.
\:\p{INTEGER}, 12, 15, 27, 38--39, 49, 62, 76, 83, 91, 96, 101, 107, 116, 124,
145, 166\*, 170\*.
\:\p{INTEGER\_IO}, 170\*.
\:\p{IS\_OPEN}, 24, 26, 29.
\:\|{j}, \[34], \[65], \[69], \[101], \[145].
\:\\{join}, \[72], 98, 101, 132.
\:\\{jump\_out}, 31, \[35].
\:\|{k}, \[33], \[48], \[51], \[65], \[69], \[84], \[94], \[96], \[98], \[101],
\[137].
\:\\{kj}, \[66].
\:\\{kl}, \[61].
\:\|{l}, \[32], \[51], \[65], \[69].
\:\p{LAST}, 13, 166\*.
\:\\{last\_unnamed}, \[70], 163.
\:\\{left\_arrow}, 16, 102, 138.
\:\\{left\_label\_bracket}, 16, 102, 138.
\:\\{length}, \[40], 53.
\:\\{less}, \[63], 65--67, 69.
\:\\{less\_or\_equal}, 16, 102, 138.
\:\\{limit}, 28--29, 33, \[116], 119, 121--123, 125--130, 133--134, 137--138,
143, 152--154.
\:\\{line}, 33, 93--94, \[116], 117, 121--122, 124, 126, 128--130.
\:\\{line\_feed}, 16, 28.
\:\\{line\_length}, \[10], 91, 94, 97, 101, 105--106, 111, 113.
\:\\{lines\_dont\_match}, \\{119}, 124.
\:\\{link}, 38, \[39], 40, 45, 49, 53, 70.
\:\\{llink}, \[45], 65--66, 69.
\:\\{loc}, 29, 33, \[116], 121, 125--127, 129--130, 133--135, 137--140,
143--144, 146, 151--154, 158.
\:\.{Long line must be truncated}, 94.
\:\\{long\_name\_index}, \[62], 65, 69, 137.
\:\\{longest\_name}, \[10], 62, 143, 145.
\:\&{loop}, 9.
\:\\{mark\_error}, \[11], 32.
\:\\{mark\_fatal}, \[11], 35.
\:\\{mark\_harmless}, \[11], 100, 145.
\:\\{max\_bytes}, \[10], 39, 59, 66.
\:\\{max\_id\_length}, \[10], 104.
\:\\{max\_names}, \[10], 39, 59, 66, 87.
\:\\{max\_texts}, \[10], 39, 70, 87, 149.
\:\\{max\_tok\_ptr}, \[42], 88, 166\*, 168.
\:\\{max\_toks}, \[10], 39, 74, 90.
\:\\{misc}, \[92], 93, 97--99, 101--102, 110, 112--113.
\:\.{Missing n )}, 150.
\:\&{mitpo}, \[8].
\:\\{mod\_field}, \[76].
\:\\{mod\_lookup}, 62--63, \[64], \\{65}, 141--142.
\:\\{mod\_range}, \[76], 156.
\:\\{mod\_text}, \[62], 65--69, 137, 141--145.
\:\p{MODE}, 24, 26.
\:\\{mode\_of\_file}, 23, \[24].
\:\\{module\_count}, 132, \[156], 157, 162, 167.
\:\\{module\_flag}, 70, 82, 163.
\:\\{module\_name}, \[131], 132, 136, 140, 148, 151, 158, 160.
\:\\{module\_number}, \[83], 84, 112.
\:\|{n}, \[124].
\:\p{NAME}, 24, 26, 166\*.
\:\.{Name does not match}, 69.
\:\p{NAME\_ERROR}, 24, 26.
\:\\{name\_field}, \[76].
\:\\{name\_index}, \[39], 40, 69.
\:\\{name\_info}, \[39], 45.
\:\\{name\_pointer}, \[40], 47--48, 50--51, 64--65, 68--69, 76, 81, 136, 155,
157.
\:\\{name\_ptr}, \[40], 48, 51, 53, 55, 57, 59, 66, 87--90, 168.
\:\p{NEW\_LINE}, 20.
\:\\{new\_ln}, \[20], 33, 94--95, 101, 125, 150, 166\*, 168--169.
\:\\{new\_module}, \[131], 132--133, 137, 151, 161, 167.
\:\\{next\_control}, \[148], 149, 158--161, 167.
\:\.{No change file.}, 166\*.
\:\.{No output was specified}, 100.
\:\.{No parameter given for macro}, 87.
\:\\{normal}, \[44], 49--51, 55, 57--59, 86, 151.
\:\.{Not present: <section name>}, 85.
\:\\{not\_equal}, 16, 102, 138.
\:\\{not\_found}, 61.
\:\\{not\_sign}, 16, 102.
\:\\{null\_array}, \[40], 41--43, 49.
\:\\{num\_or\_id}, \[92], 98--99, 102.
\:\p{OPEN}, 24.
\:\\{open\_a\_file}, \[23], \\{24}, 165--166\*.
\:\\{opened}, \[165], 166\*.
\:\&{optim}, \[8].
\:\\{or\_sign}, 16, 102.
\:\\{other\_line}, \[116], 117, 126, 130.
\:\\{out\_buf}, 34, \[91], 93--94, 96--97.
\:\\{out\_buf\_size}, \[10], 91, 96.
\:\\{out\_buffer\_index}, \[91], 94, 96.
\:\\{out\_contrib}, \[97], 98, 101--102, 104--106, 110--112.
\:\p{OUT\_FILE}, 26, 166\*.
\:\\{out\_ptr}, 34, \[91], 92--96, 98--99, 113.
\:\\{out\_state}, \[92], 93, 98--99, 101--102, 105, 112--113.
\:\.{Output file name didn't end}, 154.
\:\\{output\_phase}, \\{6\*}, 166\*.
\:\.{output\_phase.adb}, 6\*.
\:\.{output\_phase.ads}, 6\*.
\:\\{output\_state}, \\{76}.
\:\\{output\_tok}, \[72], 101, 108, 132, 151, 154.
\:\\{overflow}, \[37], 59, 66, 74, 81, 87, 90, 149.
\:\|{p}, \[48], \[51], \[65], \[69], \[81], \[155], \[157].
\:\p{PACK}, 39.
\:\\{param}, \[72], 84, 90, 149.
\:\\{parametric}, \[44], 50, 82, 86, 148--149, 159.
\:\\{phase\_i}, \[7\*], \\{164}, 166\*.
\:\\{phase\_ii}, \[6\*], \\{114}, 166\*.
\:\\{phase\_one}, \[30], 32, 167.
\:\\{pop\_level}, \\{82}, 84, 87--88.
\:\p{POS}, 14.
\:\\{prefix}, \[63], 67.
\:\\{prefix\_lookup}, \[68], \\{69}, 141.
\:\\{prime\_the\_change\_buffer}, \\{120}, 126, 129.
\:\\{print}, \[20], 33--34, 48, 61, 85, 90, 94--95, 101, 125, 132, 145, 150,
168.
\:\\{print\_id}, \[47], \\{48}, 85, 87.
\:\\{print\_int}, \[20], 33--34, 94--95, 101, 125, 132, 150, 168.
\:\\{print\_ln}, \[20], 33--34, 164, 166\*, 168.
\:\\{print\_nl}, \[20], 29, 31, 35, 61, 85, 87, 100, 145, 168--169.
\:\.{Program ended at brace level n}, 95.
\:\\{push\_level}, \\{81}, 85--86, 89.
\:\p{PUT}, 20, 94.
\:\|{q}, \[51], \[65], \[69].
\:\|{r}, \[69].
\:\&{repeat}, \[9].
\:\\{repl\_field}, \[76].
\:\\{restart}, 84--87, 89, 99, 127, 137, 140.
\:\\{reswitch}, 101, 105, 110--112.
\:\\{right\_arrow}, 16, 102, 138.
\:\\{right\_label\_bracket}, 16, 102, 138.
\:\\{rlink}, \[45], 46, 65--66, 69.
\:\|{s}, \[51].
\:\\{scan\_module}, 156, \\{157}, 167.
\:\\{scan\_repl}, 148, \\{149}, 155, 160.
\:\\{scanning\_based\_number}, \[109], 110, 136--137.
\:\\{scanning\_character}, \[136], 137, 140, 146.
\:\.{Section ended in mid-comment}, 135.
\:\.{Section name didn't end}, 144.
\:\.{Section name too long}, 145.
\:\\{semi\_ptr}, \[91], 93--95, 98, 112.
\:\\{send\_out}, 97, \\{98}, 100--102, 104--106, 110--113.
\:\\{send\_the\_output}, 100, \\{101}.
\:\\{send\_val}, \\{96}, 100, 102, 112.
\:\\{set\_element\_sign}, 16, 102.
\:\\{simple}, \[44], 50, 86--87, 148, 158.
\:\\{sixteen\_bits}, \[38], 39, 44, 49, 64--65, 68--70, 73--74, 76, 83--84, 98,
148--149.
\:\\{skip\_ahead}, \\{133}, 140, 158, 161, 167.
\:\\{skip\_comment}, \\{134}, 137.
\:\p{SKIP\_LINE}, 29.
\:\.{Sorry, x capacity exceeded}, 37.
\:\\{spotless}, \[11], 169.
\:\\{stack}, \[76], 81--82.
\:\\{stack\_ptr}, \[76], 80--82, 84, 87, 101, 105--106, 108.
\:\\{stack\_size}, \[10], 76, 81.
\:\p{STANDARD\_OUTPUT}, 20.
\:\&{stat}, \[8].
\:\p{STATUS\_ERROR}, 24, 26.
\:\\{store\_two\_bytes}, \[73], \\{74}, 90, 162.
\:\\{str}, \[97], 98, 102, 105--106, 110, 112--113.
\:\p{STRING}, 23--24, 26, 107.
\:\.{String didn't end}, 152.
\:\.{String too long}, 105.
\:\\{success}, 23, \[24].
\:{system dependencies}, \[1\*], 13, 17, 20--21, 24, 26, 28, 33, 165, 169--170%
\*.
\:\|{t}, \[51], \[98], \[149], \[155].
\:\\{tab\_mark}, 16, 33, 132, 135, 137, 143--144.
\:\&{tats}, \[8].
\:\\{temp\_line}, \[116], 117.
\:\\{text\_char}, \[13], 14--15, 18, 20, 29.
\:\\{text\_file}, \[13], 22--25, 28--29.
\:\\{text\_index}, \[39], 42.
\:\p{TEXT\_IO}, 3\*--7\*, 13, 20, 23, \\{24}, \\{26}, 29, 94, 108, 166\*, 170\*.
\:\\{text\_link}, 38, \[39], 42, 70--71, 80, 82, 87, 100, 155, 163.
\:\\{text\_pointer}, \[42], 70, 76, 148.
\:\\{text\_ptr}, \[42], 78, 87--88, 149, 168.
\:\.{This can't happen}, 36.
\:\.{This identifier has already...}, 57.
\:\.{This identifier was defined...}, 57.
\:\\{tok\_mem}, 38, \[39], 42, 70, 74, 76--78, 84, 87, 90, 149.
\:\\{tok\_ptr}, \[42], 74, 78, 87--88, 90, 149, 166\*.
\:\\{tok\_start}, 38, \[39], 42--43, 70, 76, 80--82, 87--88, 149.
\:\\{token\_index}, \[39], 42.
\:\\{token\_info}, \[39].
\:\\{token\_memory}, \[39].
\:\p{TRUE}, 24, 26, 28--30, 112, 116--117, 119, 124, 126, 128, 130, 136--137,
140, 146, 167.
\:\|{u}, \[96].
\:\\{unambig\_length}, \[10], 44, 49, 56, 61.
\:\\{unbreakable}, \[92], 99, 101, 105, 112.
\:\&{until}, \[9].
\:\\{update\_terminal}, \[21], 32, 94, 100, 132.
\:{uppercase}, 56, 61, 102, 110--111.
\:\.{Use == for macros}, 159.
\:\p{USE\_ERROR}, 24, 26.
\:\|{v}, \[96], \[98].
\:\p{VAL}, 14.
\:\\{verbatim}, \[72], 101, 106, 132, 151, 153.
\:\.{Verbatim string didn't end}, 153.
\:\.{Verbatim string too long}, 106.
\:\|{w}, \[48], \[51], \[65], \[69], \[84], \[101].
\:\\{w\_range}, \[39], 40, 48, 51, 65, 69, 84, 101.
\:\\{web\_file}, \[22], 33, 116, 118, 124, 128, 130, 166\*.
\:\\{web\_input\_file\_name}, \[165].
\:\.{Where is the match...}, 121, 125, 129.
\:\\{wo}, \[168].
\:\\{ww}, 10, \[39], 40--41, 48, 54, 59, 61, 66--67, 87, 104, 168.
\:\|{x}, \[74].
\:\\{xchr}, \[14], 15, 17--18, 33--34, 48, 61, 94, 108, 145, 151.
\:\\{xord}, \[14], 18, 29.
\:\.{You should double \AT! signs}, 137, 152--154.
\:\|{z}, \[42].
\:\\{z\_range}, \[39], 42, 77, 166\*.
\:\\{zo}, \[77], 80--82, 84, 87, 90, \[166\*], \[168].
\:\\{zz}, 10, \[39], 42--43, 77, 80--82, 87--88, 149, 168.
\fin
\:\X112:Cases involving \.{@\{} and \.{@\}}\X
\U section~101.
\:\X102:Cases like \.{/=} and \.{:=}\X
\U section~101.
\:\X110:Cases related to constants, possibly leading to \\{get\_fraction} or %
\\{reswitch}\X
\U section~101.
\:\X104:Cases related to identifiers\X
\U section~101.
\:\X60:Check for ambiguity and update secondary hash\X
\U section~59.
\:\X145:Check for overlong name\X
\U section~143.
\:\X61:Check if \|q conflicts with \|p\X
\U section~60.
\:\X130:Check that all changes have been read\X
\U section~167.
\:\X161:Check that \?=\? or \?\S\? follows this module name, otherwise %
\&{return}\X
\U section~160.
\:\X54:Compare name \|p with current identifier, \&{goto}\8\\{found} if equal\X
\U section~53.
\:\X138:Compress two-symbol combinations like `\.{:=}'\X
\U section~137.
\:\X52:Compute the hash code \|h\X
\U section~51.
\:\X53:Compute the name location \|p\X
\U section~51.
\:\X56:Compute the secondary hash code \|h and put the first characters into
the auxiliary array \\{chopped\_id}\X
\U section~55.
\:\X152:Copy a string from the buffer to \\{tok\_mem}\X
\U section~149.
\:\X154:Copy output file-name from the buffer to \\{tok\_mem}\X
\U section~151.
\:\X90:Copy the parameter into \\{tok\_mem}\X
\U section~87.
\:\X153:Copy verbatim string from the buffer to \\{tok\_mem}\X
\U section~151.
\:\X135:Do special things when \|c\?=\?\.{"@"}\?,\?\.{"\\"}\?,\?\.{"\{"}\?,\?%
\.{"\}"}; \&{return} at end\X
\U section~134.
\:\X95:Empty the last line from the buffer\X
\U section~100.
\:\X59:Enter a new identifier into the table at position \|p\X
\U section~55.
\:\X66:Enter a new module name into the tree\X
\U section~65.
\:\X86:Expand macro \|a and \&{return}, or \&{goto}\8\\{restart} if no output
found\X
\U section~84.
\:\X85:Expand module \|a\?-\bn{8}{24000}\?, \&{goto}\8\\{restart}\X
\U section~84.
\:\X113:Force a line break\X
\U sections~101, 108, and~112(3).
\:\X139:Get an identifier\X
\U section~137.
\:\X140:Get control code and possible module name\X
\U section~137.
\:\X99:Get the buffer ready for appending the new information\X
\U section~98.
\:\X108:Get the \\{file\_name} and redirect further output\X
\U section~101.
\:\X57:Give double-definition error and change \|p to type \|t\X
\U section~55.
\:\X144:If end of name, \&{exit}\X
\U section~143.
\:\X125:If the current line starts with \.{@y}, report any discrepancies and %
\&{return}\X
\U section~124.
\:\X159:If the next text is `\?(\?\#\?)\S\?', call \\{define\_macro} and %
\&{goto}\8\\{continue}\X
\U section~158.
\:\X146:If this is the beginning of an character token then set \\{scanning%
\_character} to \p{TRUE} and \?\&{return}\?\8\\{character\_tok}; otherwise do
nothing\X
\U section~137.
\:\X151:In cases that \|a is a non-ASCII token (\\{identifier}, \\{module%
\_name}, etc.), either process it and change \|a to a byte that should be
stored, or \&{goto}\8\\{continue} if \|a should be ignored, or \?\&{exit}\8%
\&{when}\?\8\|a signals the end of this replacement text\X
\U section~149.
\:\X126:Initialize the input system\X
\U section~164.
\:\X93:Initialize the output buffer\X
\U section~100.
\:\X80:Initialize the output stacks\X
\U section~100.
\:\X162:Insert the module number into \\{tok\_mem}\X
\U section~160.
\:\X150:Make sure the parentheses balance\X
\U section~149.
\:\X123:Move \\{buffer} and \\{limit} to \\{change\_buffer} and \\{change%
\_limit}\X
\U sections~120 and~124.
\:\X103:Other printable characters\X
\U section~101.
\:\X167:Phase I: Read all the user's text and compress it into \\{tok\_mem}\X
\U section~164.
\:\X100:Phase II: Output the contents of the compressed tables\X
\U section~114.
\:\X33:Print error location based on input buffer\X
\U section~32.
\:\X34:Print error location based on output buffer\X
\U section~32.
\:\X168:Print statistics about memory usage\X
\U section~166\*.
\:\X169:Print the job \\{history}\X
\U section~166\*.
\:\X48, 74:Procedures in \\{data\_structures}\X
\U section~3\*.
\:\X51, 65, 69:Procedures in \\{hashing}\X
\U section~5\*.
\:\X24, 26, 29, 32:Procedures in \\{input\_output}\X
\U section~4\*.
\:\X119, 120, 124, 127, 132, 133, 134, 137, 149, 155, 157, 164:Procedures in %
\\{input\_phase}\X
\U section~7\*.
\:\X81, 82, 84, 94, 96, 98, 101, 114:Procedures in \\{output\_phase}\X
\U section~6\*.
\:\X87:Put a parameter on the parameter stack, or \&{goto}\8\\{restart} if
error occurs\X
\U section~86.
\:\X143:Put module name into \\{mod\_text}\?(\?1\?\to\?\|k\?)\?\X
\U section~141.
\:\X129:Read from \\{change\_file} and maybe turn off \\{changing}\X
\U section~127.
\:\X128:Read from \\{web\_file} and maybe turn on \\{changing}\X
\U section~127.
\:\X88:Remove a parameter from the parameter stack\X
\U section~82.
\:\X58:Remove \|p from secondary hash table\X
\U section~57.
\:\X160:Scan the \ADA\ part of the current module\X
\U section~157.
\:\X158:Scan the definition part of the current module\X
\U section~157.
\:\X141:Scan the module name and make \\{cur\_module} point to it\X
\U section~140.
\:\X105:Send a string, \&{goto}\8\\{reswitch}\X
\U section~101.
\:\X106:Send verbatim string\X
\U section~101.
\:\X67:Set \(\|c to the result of comparing the given name to name \|p\X
\U sections~65 and~69.
\:\X15, 17, 18, 41, 43, 46, 71, 142:Set initial values\X
\U section~3\*.
\:\X121:Skip over comment lines in the change file; \&{return} if end of file\X
\U section~120.
\:\X122:Skip to the next nonblank line; \&{return} if end of file\X
\U section~120.
\:\X111:Special code to finish real constants\X
\U section~101.
\:\X47, 73:Specification of procedures visible from \\{data\_structures}\X
\U section~3\*.
\:\X50, 64, 68:Specification of procedures visible from \\{hashing}\X
\U section~5\*.
\:\X23, 28, 31:Specification of procedures visible from \\{input\_output}\X
\U section~4\*.
\:\X89:Start scanning current macro parameter, \&{goto}\8\\{restart}\X
\U section~84.
\:\X10, 11, 12, 13, 14, 16, 22, 25, 27, 30, 35, 38, 39, 40, 42, 44, 45, 49, 62,
63, 70, 72, 76, 77, 79, 83, 91, 92, 97, 107, 109, 116, 118, 131, 136, 148, 156,
165:Types and objects visible from \\{data\_structures}\X
\U section~3\*.
\:\X163:Update the data structure so that the replacement text is accessible\X
\U section~160.
\:\X55:Update the tables and check for possible errors\X
\U section~51.
\con
